% book example for classicthesis.sty
\documentclass[11pt,a4paper,footinclude=true,headinclude=true]{scrbook} % KOMA-Script book
\usepackage[T1]{fontenc}                
\usepackage{lipsum}
\usepackage[linedheaders,parts,pdfspacing]{classicthesis} % ,manychapters
%\usepackage[osf]{libertine}

%行距
\renewcommand{\baselinestretch}{1.15}
\usepackage[top=2.25cm,bottom=1.5cm,hmargin=2cm]{geometry}

%packages
\usepackage{dsfont}
\usepackage{polynom}
\usepackage{empheq}
%\usepackage[latin1]{inputenc}
\usepackage[utf8]{inputenc}
%\usepackage[T1]{fontenc}
%\usepackage[ngerman]{babel}
\usepackage{subfiles}
\usepackage{graphicx}
\usepackage{systeme}
\usepackage{mathtools, stmaryrd}
\usepackage{tikz-cd}
\usepackage{extpfeil}
%\usetikzlibrary{decorations.markings}
\usepackage{CJK}
\usepackage{amsmath, amsthm, amssymb, amsfonts, mathrsfs, bm, fancyhdr, color, comment, graphicx, environ}
\usepackage{xcolor}
\usepackage{mdframed}
\usepackage[shortlabels]{enumitem}
\usepackage{indentfirst}
\usepackage{hyperref}
\usepackage{ifthen}
\usepackage{enumerate}

%custom commands: symbols
\newcommand{\trianglerightneq}{\mathrel{\ooalign{\raisebox{-0.5ex}{\reflectbox{\rotatebox{90}{$\nshortmid$}}}\cr$\triangleright$\cr}\mkern-3mu}}
\newcommand{\triangleleftneq}{\mathrel{\reflectbox{$\trianglerightneq$}}}
\renewcommand\qedsymbol{$\square$}
\DeclareMathOperator{\E}{E} % expectation operator
\newcommand\myeq[1]{\stackrel{\textnormal{#1}}{=}}

\newcommand{\verteq}{\rotatebox{90}{$\,=$}}
\newcommand{\acton}{\rotatebox{180}{$\,\circlearrowright$}}
\newcommand{\equalto}[2]{\underset{\scriptstyle\overset{\mkern4mu\verteq}{#2}}{#1}}
\newcommand*{\longhookrightarrow}{\ensuremath{\lhook\joinrel\relbar\joinrel\rightarrow}}

\DeclareMathOperator{\Ann}{Ann}
\DeclareMathOperator{\Ass}{Ass}
\DeclareMathOperator{\Supp}{Supp}
\DeclareMathOperator{\WeakAss}{\widetilde{Ass}}
\let\Im\relax
\DeclareMathOperator{\Im}{Im}
\DeclareMathOperator{\Spec}{Spec}
\DeclareMathOperator{\SPEC}{\mathbf{Spec}}
\DeclareMathOperator{\Sp}{sp}
\DeclareMathOperator{\maxSpec}{mSpec}
\DeclareMathOperator{\Hom}{Hom}
\DeclareMathOperator{\Soc}{Soc}
\DeclareMathOperator{\Ht}{ht}
\let\AA\relax
\DeclareMathOperator{\AA}{\mathbb{A}}
\DeclareMathOperator{\V}{\mathbf{V}}
\DeclareMathOperator{\Aut}{Aut}
\DeclareMathOperator{\Char}{char}
\DeclareMathOperator{\Frac}{Frac}
\DeclareMathOperator{\Proj}{Proj}
\DeclareMathOperator{\PROJ}{\mathbf{Proj}}
\DeclareMathOperator{\stimes}{\text{\footnotesize\textcircled{s}}}
\DeclareMathOperator{\End}{End}
\DeclareMathOperator{\Ker}{Ker}
\DeclareMathOperator{\Coker}{coker}
\DeclareMathOperator{\LCM}{LCM}
\DeclareMathOperator{\CDiv}{CDiv}
\DeclareMathOperator{\Div}{Div}
\DeclareMathOperator{\Pic}{Pic}
\DeclareMathOperator{\id}{id}
\DeclareMathOperator{\Cl}{Cl}
\DeclareMathOperator{\CaCl}{CaCl}
\DeclareMathOperator{\dv}{div}
\DeclareMathOperator{\Der}{Der}
%\DeclareMathOperator{\Gr}{Gr}
\DeclareMathOperator{\Gal}{Gal}
\DeclareMathOperator{\nil}{nil}
\DeclareMathOperator{\pr}{pr}
\DeclareMathOperator{\trdeg}{trdeg}
\DeclareMathOperator{\rank}{rank}
\DeclareMathOperator{\cd}{cd}
\DeclareMathOperator{\codim}{codim}
\DeclareMathOperator{\sgn}{sgn}
\DeclareMathOperator{\Ob}{Ob}
\DeclareMathOperator{\GL}{GL}
\newcommand{\EE}{\mathscr{E}}
\newcommand{\FF}{\mathscr{F}}
\newcommand{\GG}{\mathscr{G}}
\newcommand{\HH}{\mathscr{H}}
\newcommand{\II}{\mathscr{I}}
\newcommand{\JJ}{\mathscr{J}}
\newcommand{\KK}{\mathscr{K}}
\newcommand{\LL}{\mathscr{L}}
\newcommand{\MM}{\mathscr{M}}
\newcommand{\NN}{\mathcal{N}}
\newcommand{\OO}{\mathcal{O}}
\newcommand{\Ss}{\mathscr{S}}
\newcommand{\TT}{\mathcal{T}}
\newcommand{\Af}{\mathfrak{A}}
\newcommand{\Bf}{\mathfrak{B}}
\newcommand{\Aa}{\mathscr{A}}
\newcommand{\PP}{\mathbb{P}}
\let\aa\relax
\newcommand{\aa}{\mathfrak{a}}
\newcommand{\bb}{\mathfrak{b}}
\newcommand{\dd}{\mathfrak{d}}
\newcommand{\pp}{\mathfrak{p}}
\newcommand{\qq}{\mathfrak{q}}
\newcommand{\nn}{\mathfrak{n}}
\newcommand{\mm}{\mathfrak{m}}
\newcommand{\red}{\mathrm{red}}
\newcommand{\reg}{\mathrm{reg}}
\newcommand{\Shf}{\mathfrak{Shf}}
\newcommand{\Psh}{\mathfrak{Psh}}
\newcommand{\LRS}{\mathfrak{LRS}}
\newcommand{\Sch}{\mathfrak{Sch}}
\newcommand{\Var}{\mathfrak{Var}}
\newcommand{\Ab}{\mathfrak{Ab}}
\newcommand{\Ring}{\mathfrak{Ring}}
\newcommand{\Top}{\mathfrak{Top}}
\newcommand{\<}{\langle}
\let\>\relax
\newcommand{\>}{\rangle}
\DeclareMathOperator{\In}{in}
\DeclareMathOperator{\Ext}{Ext}
\DeclareMathOperator{\Spe}{Sp\acute{e}}
\DeclareMathOperator{\HHom}{\mathscr{H}om}
\newcommand{\isoto}{\overset{\sim}{\to}}
\newcommand{\isolongto}{\overset{\sim}{\longrightarrow}}
\newcommand{\spn}{\mathchar`-\mathrm{span}}
\newcommand{\GR}{\mathfrak{Gr}}
\newcommand{\fg}{\mathfrak{f.g.}}
\newcommand{\MOD}{\mathfrak{Mod}}
%\newcommand{\qgr}{\mathsf{qgr}\mathchar`-}
%\newcommand{\uqgr}{\underline{\mathsf{qgr}}\mathchar`-}
\newcommand{\qcoh}{\mathfrak{Qch}}
\newcommand{\ALG}{\mathfrak{Alg}}
\newcommand{\coh}{\mathfrak{Coh}}
\newcommand{\vect}{\mathsf{vect}\mathchar`-}
\newcommand{\imm}[1][imm]{\hspace{0.75ex}\raisebox{0.58ex}{%

\begin{tikzpicture}[commutative diagrams/every diagram]
\draw[commutative diagrams/.cd, every arrow, every label,hook,{#1}] (0,0ex) -- (2.25ex,0ex);
\end{tikzpicture}}\hspace{0.75ex}}

\newcommand{\dashto}[2]{\smash{\hspace{-0.7em}\begin{tikzcd}[column sep=small,ampersand replacement=\&] {#1} \rar[dashed] \& {#2} \end{tikzcd}\hspace{-0.7em}}}

\tikzset{commutative diagrams/.cd,
mysymbol/.style={start anchor=center,end anchor=center,draw=none}
}
\newcommand\FibSq[2][\square]{%
  \arrow[mysymbol]{#2}[description]{#1}}
\newcommand\Comm[2][\acton]{%
  \arrow[mysymbol]{#2}[description]{#1}}



%custom commands: theorem
\newtheorem{innercustomthm}{Theorem}
\newenvironment{theorem}[2][]
{\renewcommand\theinnercustomthm{#2}\innercustomthm[#1]
\begin{mdframed}[backgroundcolor=yellow!20]\ }
{\end{mdframed}\endinnercustomthm}
\newtheorem*{innertheorem*}{Theorem}
\newenvironment{theorem*}
{\begin{innertheorem*}\begin{mdframed}[backgroundcolor=gray!20]\ }
{\end{mdframed}\end{innertheorem*}}
\newtheorem{innercustomprop}{Proposition}
\newenvironment{proposition}[2][]
{\renewcommand\theinnercustomprop{#2}\innercustomprop[#1] 
\begin{mdframed}[backgroundcolor=gray!20]\ }
{\end{mdframed}\endinnercustomprop}
\newtheorem{innercustomlemma}{Lemma}
\newenvironment{lemma}[2][]
{\renewcommand\theinnercustomlemma{#2}\innercustomlemma[#1]
\begin{mdframed}[backgroundcolor=gray!20]\ }
{\end{mdframed}\endinnercustomprop}
\newtheorem{innercustomcol}{Corollary}
\newenvironment{corollary}[2][]
{\renewcommand\theinnercustomcol{#2}\innercustomcol[#1]
\begin{mdframed}[backgroundcolor=gray!20]\ }
{\end{mdframed}\endinnercustomcol}
\newtheorem{innercustomalg}{Algorithm}
\newenvironment{algorithm}[2][]
{\renewcommand\theinnercustomalg{#2}\innercustomalg[#1]
\begin{mdframed}[backgroundcolor=gray!20]\ }
{\end{mdframed}\endinnercustomalg}

%custom command
\newenvironment{definition}[2][Definition]
    { \begin{mdframed}[backgroundcolor=gray!20] \textbf{#1 #2.}\ }
    {  \end{mdframed}}
\newenvironment{definition*}[2][Definition]
    { \begin{mdframed}[backgroundcolor=blue!20] \textbf{#1.}\ }
    {  \end{mdframed}}
\newenvironment{remark}[2][Remark]
    { \begin{mdframed}[backgroundcolor=gray!20] \textbf{#1 #2.} }
    {  \end{mdframed}}
\newenvironment{remark*}[2][Remark]
    { \begin{mdframed}[backgroundcolor=gray!20] \textbf{#1.}\ }
    {  \end{mdframed}}    
\newenvironment{nt}[2][Note]
    { \begin{mdframed}[backgroundcolor=gray!20] \textbf{#1.}\ }
    {  \end{mdframed}}
\newenvironment{fact}[2][Fact]
    { \begin{mdframed}[backgroundcolor=gray!20] \textbf{#1 #2.} }
    {  \end{mdframed}}    
\newenvironment{example}[2][Example]
    { \begin{mdframed}[backgroundcolor=gray!20] \textbf{#1 #2.} }
    {  \end{mdframed}}
\newenvironment{exe}[2][Exercise]
    { \begin{mdframed}[backgroundcolor=gray!20] \textbf{#1 #2.}\\ }
    {  \end{mdframed}}
\newenvironment{problem}[2][Problem]
    { \begin{mdframed}[backgroundcolor=gray!20] \textbf{#1 #2.} }
    {  \end{mdframed}}
\newenvironment{notation}[2][Notation]
    { \begin{mdframed}[backgroundcolor=gray!20] \textbf{#1.} }
    {  \end{mdframed}}   
\newenvironment{axiom}[2][Axiom]
    { \begin{mdframed}[backgroundcolor=red!20] \textbf{#1 #2.} }
    {  \end{mdframed}}     
\newenvironment{motivation}[2][Motivation]
    { \begin{mdframed}[backgroundcolor=red!20] \textbf{#1 #2.} }
    {  \end{mdframed}}    
\newenvironment{proofthm}[2][Proof of Theorem]
    { \begin{mdframed}[backgroundcolor=gray!20] \textbf{#1 #2.} }
    {  \end{mdframed}}   
\newenvironment{summary}[2][]
    { \begin{mdframed}[backgroundcolor=green!20] \textbf{#1 #2.} }
    {  \end{mdframed}}    

%Setup
\hypersetup{
    colorlinks=true,
    linkcolor=blue,
    filecolor=magenta,      
    urlcolor=blue,}
    
\makeatletter
\@addtoreset{chapter}{part}
\makeatletter    


\begin{document}
%	\pagestyle{scrheadings}
%	\manualmark
%	\markboth{\spacedlowsmallcaps{\contentsname}}{\spacedlowsmallcaps{\contentsname}}
	
	\tableofcontents 

	%\automark[section]{chapter}
	%\renewcommand{\chaptermark}[1]{\markboth{\spacedlowsmallcaps{#1}}{\spacedlowsmallcaps{#1}}}
	%\renewcommand{\sectionmark}[1]{\markright{\thesection\enspace\spacedlowsmallcaps{#1}}}

% use \cleardoublepage here to avoid problems with pdfbookmark
\cleardoublepage\part{Varieties}

\chapter{Affine Varieties}

\begin{summary}{Reference}
    \begin{itemize}
        \item \textit{Introduction to Commutative Algebra} - Atiyah \& McDonald
        \item \textit{Commutative Algebra} - Matsumura
        \item (Ch II) \textit{Algebraic Geometry} - Lei Fu 
        \item (Ch II) \textit{My Way to Algebraic Geometry} - Marco lo Giudice
        \item (Ch II) \textit{Foundations of Algebraic Geometry} - Ravi Vakil
    \end{itemize}
\end{summary}

\begin{theorem}{1.3. (Hilbert's Nullstellensatz)}
    Let $k$ be an algebraically closed field, and let $\mathfrak{a}$ be an ideal in $A=k[x_1,\cdots,x_n]$, and let $f\in A$ be a polynomial which vanishes at all points of $Z(\mathfrak{a})$, i.e. $f\in I(Z(\mathfrak{a}))$. Then $f^r\in\mathfrak{a}$ for some integer $r>0$.
\end{theorem}
\begin{proof}
    Suppose on the contrary that $f\notin\sqrt{\mathfrak{a}}$, we must show that $f\notin I(Z(\mathfrak{a}))$. Recall that 
    \begin{align*}
        \sqrt{a}=\bigcap\{\pp\triangleleft A:\text{prime ideal with}\ \pp\supset\mathfrak{a}\}.
    \end{align*}
    We pick a prime ideal $\pp\triangleleft A$ with $\pp\supset\mathfrak{a}$ for which $f\notin\pp$. Consider the natural ring homomorphism $A\twoheadrightarrow A/\pp=:B$, and denote the image of $f$ by $\overline{f}\coloneqq f+\pp$. Note that $\overline{f}\neq0$ in $B$, so we may consider the localization $C\coloneqq B_{\overline{f}}$. Let $\mm\triangleleft C$ be a maximal ideal, then we have the homomorphism
    \begin{align*}
        k\xhookrightarrow{}k[x_1,\cdots,x_n]=A\twoheadrightarrow A/\pp=B\xhookrightarrow{(B:\ \text{integral})}B_{\overline{f}}=C\twoheadrightarrow C/\mm. \tag{$\star$}
    \end{align*}
    Now, $C=B_{\overline{f}}=B[\frac{1}{\overline{f}}]$ is a finitely generated  $B$-algebra, and thus, a finitely generated $k$-algebra. So $C/\mm$ is a finitely generated $k$-algebra. By \textit{Weak Nullstellensatz (AM, 7.10)}, $[C/\mm:k]<\infty$, i.e. a finite field extension over $k$. But $k$ is algebraically closed, whence $C/\mm=k$. Back to the homomorphisms in $(\star)$, we suppose
    \begin{align*}
        A\longrightarrow C/\mm=k,\quad x_i\longmapsto a_i,
    \end{align*}
    and pick $P\coloneqq(a_1,\cdots,a_n)\in\AA^n$. Now, to show that $f\notin I(Z(\mathfrak{a}))$, it suffices to show that $P\in Z(\mathfrak{a})$ and $f(P)\neq0$. The latter is clear, since $\frac{1}{\overline{f}}\neq0$. For any $g\in\mathfrak{a}(\subset\pp)$, the homomorphisms in $(\star)$ yield:
    \begin{align*}
        A\longrightarrow A/\pp\longrightarrow k,\quad g \longmapsto\overline{g}=0\longmapsto0=g(a_1,\cdots,a_n)=g(P).
    \end{align*}
    Thus, $P\in Z(\mathfrak{a})$, and therefore, $f\notin I(Z(\mathfrak{a}))$.
\end{proof}

\begin{theorem}{1.11. (Krull's Hauptidealsatz)}
    Let $A$ be a Noetherian ring, and $f\in A$ be an element which is neither a zero divisor nor a unit. Then every minimal prime ideal $\pp\triangleleft A$ containing $f$ has height $1$.
\end{theorem}

\begin{theorem}{1.12}
    A Noetherian integral domain $A$ is a UFD if and only if every prime ideal of height $1$ is principal.
\end{theorem}
\begin{proof}
    This is \textit{(Matsumura, Theorem 47, p.141)}. We use the fact that
    \begin{align*}
        \textit{A Noetherian integral domain is a UFD if and only if every irreducible element is prime}.
    \end{align*}
    
    ($\Rightarrow$) Suppose $A$ is a UFD, let $\pp\triangleleft A$ be any prime ideal with $\Ht\pp=1$. (May assume $\pp\neq0$, otherwise it is trivial). Pick a nonzero element $a\in\pp$, and factorize $a$ into irreducible elements $a=ua_1^{r_1}\cdots a_n^{r_n}$ (where $a_i\in A$ irreducible, $u\in A^\times$, $r_i>0$). Since $\pp$ is a prime ideal, we have $a_i\in\pp$ for some $i$. From the fact, $a_i\in A$ is a prime element, so $\<a_i\>\triangleleft A$ is a prime ideal contained in $\pp$. Note that $\Ht\<a_i\>=1=\Ht\pp$, we must have $\pp=\<a_i\>$.\\
    
    ($\Leftarrow$) Suppose the converse holds, it suffices to show that every irreducible element is prime. Let $0\neq a\in A$ be an irreducible element, and choose a minimal prime ideal $\pp\triangleleft A$ containing $a$. As $A$ is Noetherian, it follows from \textit{(1.11)} that $\Ht\pp=1$. By our assumption, such $\pp$ must be principal, say $\pp=\<b\>$. Then, $b\in A$ is a prime element, thus, irreducible. Since $a\in\pp=\<b\>$, and $a\in A$ is irreducible, they must be associated, whence $\<a\>=\<b\>=\pp$. Therefore, $a\in A$ is a prime element.
\end{proof}

\chapter{Projective Varieties}
\begin{notation}{(Localizations of Graded Rings)}
    Let $S$ be any graded ring, $f\in S$ be any homogeneous element, and $\pp\triangleleft S$ be any homogeneous prime ideal.
    \begin{itemize}
        \item We denote $S_f$ by the localization of $S$ on the multiplicatively closed subset $T_f\coloneqq\{f^n\}_{n\geq0}$, i.e. $S_f\coloneqq T_f^{-1}S$.
        \item In fact, $S_f$ is a graded $S$-module, with the natural grading $\deg\frac{a}{f^n}\coloneqq\deg a-n$ (where $a\in S^h$). We denote $S_{(f)}\coloneqq(S_f)_0$, the degree $0$ part of $S_f$.
        \item Similarly, $S_\pp$ also has a natural grading structure by $\deg\frac{f}{g}\coloneqq\deg f-\deg g$ (where $f,g\in S^h$, with $g\notin\pp$). We denote $S_{(\pp)}\coloneqq(S_\pp)_0$.
    \end{itemize}
\end{notation}

\begin{fact}{2.A}
    For any graded ring $S$ and any homogeneous element $f\in S^h$, we have a natural isomorphism $S_f\simeq S_{(f)}[f,f^{-1}]$.
\end{fact}

\begin{exe}{2.1. (Homogeneous Nullstellensatz)}
    If $\mathfrak{a}\triangleleft S$ is a homogeneous ideal, and $f\in S$ is any homogeneous polynomial of positive degree, such that $f(P)=0$ for every $P\in Z_+(\mathfrak{a})$ (in $\PP^n$), i.e. $f\in I_+(Z_+(\mathfrak{a}))$. Then, $f^m\in\mathfrak{a}$ for some $q>0$.
\end{exe}
\begin{proof}
    Change the problem into $\AA^{n+1}$ and use the usual \textit{Nullstellensatz}.
\end{proof}

\begin{exe}{2.2}
    For a homogeneous ideal $\mathfrak{a}\triangleleft S$, T.F.A.E.:
    \begin{itemize}
        \item[(i)] $Z_+(\mathfrak{a})=\emptyset$
        \item[(ii)] $\sqrt{\mathfrak{a}}=$ either $S$ or $S_+$, the \textit{irrelevant} maximal ideal.
        \item[(iii)] $\mathfrak{a}\supset S_d$ for some $d>0$, (and thus,  $\mathfrak{a}\supset S_d$ for all large $d>0$). 
    \end{itemize}
\end{exe}

\chapter{Morphisms}

\begin{fact}{3.A}
    Let $A$ be an integral domain with $K=\Frac A$. Then, for any prime ideal $\pp\triangleleft A$, $A_\pp$ can be viewed as a subring of $K$, and that
    \begin{align*}
        A=\bigcap\{A_\pp:\pp\triangleleft A\ \text{prime ideal}\}=\bigcap\{A_\mm:\mm\triangleleft A\ \text{maximal ideal}\}.
    \end{align*}
\end{fact}
\begin{proof}
    This is \textit{(Matsumura, Lemma 2, p.8)}. Note that the containment $\subset$ is clear (as $A$ is an integral domain).
\end{proof}

\begin{fact}{3.B}
     Fix a graded ring $S$. Recall that for any homogeneous prime ideal $\pp\triangleleft S$, the localization $S_\pp$ is a graded $S$-module, which is also a local ring with maximal ideal $\pp S_\pp$.
     \begin{itemize}
         \item Show that its degree $0$ part $S_{(\pp)}$ is also a local ring with maximal ideal $(\pp S_\pp)\cap S_{(\pp)}$.
         \item In particular, if $S$ is an integral domain, and $\pp=(0)$ the zero ideal, show that $S_{((0))}=(\Frac S)_0$ is a field.
     \end{itemize}
\end{fact}

\cleardoublepage\part{Schemes}

\chapter{Sheaves of Modules}

\begin{definition}{(Sheaves of Modules)}
    \begin{itemize}
        \item Let $(X,\OO_X)$ be a ringed space. A \textbf{sheaf of $\OO_X$-modules} (or simply an \textbf{$\OO_X$-module}) is a sheaf $\FF$ on $X$, such that for each open set $U\subset X$, the group $\FF(U)$ is an $\OO_X(U)$-module, and for each inclusion of open sets $V\subset U$, the restriction homomorphism $\FF(U)\rightarrow\FF(V)$ is compatible with the module structures via the ring homomorphism $\OO_X(U)\rightarrow\OO_X(V)$.
        \item A \textbf{morphism} $\FF\to\GG$ of sheaves of $\OO_X$-module is a morphism of sheaves, such that for each open set $U\subset X$, the map $\FF(U)\to\GG(U)$ is an $\OO_X(U)$-module homomorphism.
    \end{itemize}
\end{definition}
 
\begin{definition}{(Associated Sheaf on $\Spec A$)}
    Fix a ring $A$ and a graded $A$-module $M$. We define the \textit{sheaf associated} to $M$ on $\Spec A$, denoted by $\widetilde{M}$, as follows. For any open subset $U\subset\Spec A$, define
    \begin{align*}
        \widetilde{M}(U)\coloneqq\Big\{s:U\to\coprod_{\pp\in U}M_{\pp}\Big|\ s\ \text{is locally a fraction}\ \frac{m}{f}\ (m\in M,\ f\in A)\Big\}.
    \end{align*}
    This make $\widetilde{M}$ into a sheaf with obvious restriction maps.
\end{definition} 
    
\begin{proposition}{5.1}
    Let $\widetilde{M}$ be the sheaf on $X=\Spec A$ associated to $M\in \MOD_A$. Then,
    \begin{itemize}
        \item[(a)] $\widetilde{M}$ is an $\OO_X$-module.
        \item[(b)] For each $\pp\in X$, we have $(\widetilde{M})_\pp\simeq M_\pp$.
        \item[(c)] For any $f\in A$, we have $\widetilde{M}(D(f))\simeq M_f$ as $A_f$-modules.
        \item[(d)] In particular, $\Gamma(X,\widetilde{M})=M$.
    \end{itemize}
\end{proposition}

\begin{proposition}{5.2}\newline
    Given $A\xrightarrow[\Ring]{\phi}B$ and its corresponding morphism $Y=\Spec B\xrightarrow[\Sch]{f}\Spec A=X$. Then,
    \begin{itemize}
        \item[(a)] The map $M\mapsto\widetilde{M}$ gives an exact, fully faithful functor $\MOD_{A}\to\MOD_{\OO_X}$.
        \item[(b)] If $M,N\in \MOD_A$, then $(M\otimes_AN)\widetilde{}\simeq\widetilde{M}\otimes_{\OO_X}\widetilde{N}$.
        \item[(c)] If $M_i\in \MOD_A\ (i\in I)$, then $(\bigoplus_{i\in I}M_i)\widetilde{}\simeq\bigoplus_{i\in I}\widetilde{M_i}$.
        \item[(d)] For any $N\in \MOD_B$,  $f_\ast(\widetilde{N})\simeq\,_A\widetilde{N}$ as $\OO_X$-modules, where $N=\,_AN\in \MOD_A$.
        \item[(e)] For any $M\in \MOD_A$,  $f^\ast(\widetilde{M})\simeq(M\otimes_AB)\widetilde{}$ as $\OO_Y$-modules.
    \end{itemize}
\end{proposition}
\begin{proof}\mbox{}
    \begin{itemize}
        \item[(a)] \textbf{(Exactness).} Given a short exact sequence $0\to M'\to M\to M''\to0$ of $A$-modules. Since localization is an exact functor, by \textit{(5.1b)}, we have
        \begin{center}
            \begin{tikzcd}
                0\rar[] &M'_\pp\rar[]\dar["\simeq"] &M_\pp\rar[]\dar["\simeq"] &M''_\pp\rar[]\dar["\simeq"] &0 \\
                0\rar[] &(\widetilde{M}')_\pp\rar[] &(\widetilde{M})_\pp\rar[] &(\widetilde{M}'')_\pp\rar[] &0
            \end{tikzcd}
        \end{center}
        for every $\pp\in\Spec A$. Therefore, \textit{[1.2c]} implies that the sequence $0\to\widetilde{M}'\to\widetilde{M}\to\widetilde{M}''\to0$ of $\OO_X$-modules is exact.
        \item[(a)] \textbf{(Full Faithfulness).} We claim that $\Hom_{\MOD_A}(M,N)\simeq\Hom_{\MOD_{\OO_X}}(\widetilde{M},\widetilde{N})$.\\
        Clearly, a homomorphism $M\xrightarrow[\MOD_A]{\tilde{\phi}}N$ induces the morphism $\widetilde{M}\xrightarrow[\MOD_{\OO_X}]{\tilde{f}}\widetilde{N}$. Conversely, if we have a morphism $\widetilde{M}\xrightarrow[\MOD_{\OO_X}]{\tilde{f}}\widetilde{N}$, by definition with $U=X$, we have $\tilde{f}_X:M=\widetilde{M}(X)\xrightarrow[\MOD_A]{}\widetilde{N}(X)=N$.
        \item[(b)] For each $\pp\in\Spec A$, we have the isomorphisms
        \begin{align*}
            (M\otimes_AN)\widetilde{}_\pp
            \underset{(5.1b)}{\simeq} (M\otimes_AN)_\pp
            \underset{(\text{AM})}{\simeq}M_\pp\otimes_{A_\pp}N_\pp \underset{(5.1b)}{\simeq}(\widetilde{M})_\pp\otimes_{\OO_{X,\pp}}(\widetilde{N})_\pp
            \underset{(\text{Def.})}{\simeq}(M\otimes_AN)\widetilde{}_\pp
        \end{align*}
        as stalks. Thus, \textit{(1.1)} implies that $(M\otimes_AN)\widetilde{}\simeq\widetilde{M}\otimes_{\OO_X}\widetilde{N}$ as sheaves.
        \item[(c)] For each $\pp\in\Spec A$, we have the isomorphisms
        \begin{align*}
            (\bigoplus_{i\in I}M_i)\widetilde{}_\pp\underset{(5.1b)}{\simeq}(\bigoplus_{i\in I}M_i)_\pp\simeq\bigoplus_{i\in I}(M_i)_\pp\underset{(5.1b)}{\simeq}\bigoplus_{i\in I}(\widetilde{M}_i)_\pp\underset{(\text{Def.})}{\simeq}(\bigoplus_{i\in I}\widetilde{M}_i)_\pp
        \end{align*}
        as stalks. Thus, $(\bigoplus_{i\in I}M_i)\widetilde{}\simeq\bigoplus_{i\in I}\widetilde{M}_i$ as sheaves.
        \item[(d)] Note that $\widetilde{A}=\OO_{\Spec A}=\OO_X$ for any ring $A$. Note that $\widetilde{\cdot}$ is a fully faithful functor by \textit{(a)}, which gives the correspondence diagram:
        \begin{center}
            (i)\qquad
            \begin{tikzcd}
                A\rar["\alpha"]\dar["\phi"] 
                & \,_AN\dar["\tilde{\phi}"] \\
                B\rar["\beta"] & N
            \end{tikzcd}
            $\qquad\longleftrightarrow\qquad$
            \begin{tikzcd}
                \OO_X\rar["\tilde{\alpha}"]\dar["f^{\#}"] & \,_A\widetilde{N}\dar["\tilde{f}^{\#}", dashed]\\
                f_\ast\OO_Y\rar["\tilde{\beta}"] & f_\ast\widetilde{N}
            \end{tikzcd}
            \qquad(ii)
        \end{center}
        where the action $\alpha$ is defined by $(a,n)\mapsto\beta(\phi(a),n)=\phi(a)\cdot n$. Further, for each $\qq\in\Spec B=Y$, the diagram (i) induces the commutative diagram:
        \begin{center}
            \begin{tikzcd}
                A_{\phi^{-1}(\qq)}\dar["\phi_\qq"]\rar[] & (_AN)_{\phi^{-1}(\qq)}\dar["\tilde{\phi}_\qq"]\\ B_\qq\rar[] &N_\qq
            \end{tikzcd}
        \end{center}
        Now, for every open set $U\subset X$, we define the sheaf map by
        \begin{align*}
            \tilde{f}^{\#}_U:\,_A\widetilde{N}(U)\longrightarrow f_\ast\widetilde{N}(U)=\widetilde{N}(f^{-1}U),\quad s\longmapsto\tilde{\phi}_\bullet\circ s\circ f
        \end{align*}
        via $\tilde{\phi}_\qq$ for every $\qq\in f^{-1}(U)$. Then, this makes the diagram (ii) commutative. Now, for each $\qq\in\Spec B$, the map $\tilde{f}^{\#}_\qq:(_A\widetilde{N})_{f(\qq)}\simeq(_AN)_{\phi^{-1}(\qq)}\xrightarrow{\sim}N_\qq\simeq(\widetilde{N})_\qq$ is a natural isomorphism of stalks, and hence, $\tilde{f}^{\#}$ gives an isomorphism $f_\ast(\widetilde{N})\simeq\,_A\widetilde{N}$ of $\OO_X$-modules.
        \item[(e)] Note that $f^{-1}\OO_X=\bigsqcup_{y\in Y}\OO_{X,f(y)}\in\Shf(Y)$; and recall that $f^\ast\FF=f^{-1}\FF\otimes_{f^{-1}\OO_X}\OO_Y$ if $\FF\in\MOD_{\OO_X}$. Now, for each $\qq\in\Spec B=Y$, we see that
        \begin{align*}
            (f^\ast\widetilde{M})_\qq &\underset{(\text{Def.})}{=}M_{\phi^{-1}(\qq)}\otimes_{A_{\phi^{-1}(\qq)}}B_\qq\simeq(M_{\phi_{-1}(\qq)}\otimes_AA_{\phi^{-1}(\qq)})\otimes_{A_{\phi^{-1}(\qq)}}B_\qq\simeq M\otimes_AB_\qq\\ &\simeq(M\otimes_AB)\otimes_BB_\qq\simeq M\otimes_AB_\qq\simeq(M\otimes_AB)_\qq\underset{(5.1b)}{\simeq}(M\otimes_AB)\widetilde{}_\qq.
        \end{align*}
        Therefore, $f^\ast\widetilde{M}\simeq(M\otimes_AB)\widetilde{}$ as $\OO_Y$-modules.
    \end{itemize}
\end{proof}

\begin{definition*}{}
    Fix a scheme $(X,\OO_X)$.
    \begin{itemize}
        \item \textbf{(Quasi-Coherent Sheaves).} A \textit{quasi-coherent sheaf} on $X$ is a sheaf $\FF\in\MOD_{\OO_X}$ such that $X$ can be covered by open affine subsets $U_i=\Spec A_i$ where for each $i,\ \FF|_{U_i}\simeq\widetilde{M}_i$ for some $M_i\in \MOD_{A_i}$.
        \item \textbf{(Coherent Sheaves).} A \textit{coherent sheaf} on $X$ is a quasi-coherent sheaf on $X$ with an additional assumption that each $M_i\in \MOD_{A_i}$ is finitely generated. 
    \end{itemize}
\end{definition*}

\begin{lemma}{5.3}
    Let $\FF$ be a quasi-coherent sheaf on $X=\Spec A$ and let $f\in A$.
    \begin{itemize}
        \item[(a)] If the global section $s\in\Gamma(X,\FF)$ satisfies $s|_{D(f)}=0$ in $\FF(D(f))$, then $f^ns=0$ for some $n>0$.
        \item[(b)] Given $t\in\FF(D(f))$, then $f^nt$ extends to a global section in $\Gamma(X,\FF)$ for some $n>0$.
    \end{itemize}
\end{lemma}
\begin{proof}\mbox{}
    \begin{itemize}
        \item[(I)] Some settings.
        \begin{itemize}
            \item Since $\FF$ is quasi-coherent on $X$, cover $X$ by open affine subsets $V=\Spec B$ such that $\FF|_V\simeq\widetilde{M}$ for some $M\in \MOD_B$.
            \item Note that $\{D(g)\}_{g\in A}$ forms a base for $\TT_X$. Cover $V$ by some $D(g)$'s, then the inclusion $\Spec A_g=D(g)\xhookrightarrow[\Sch]{i}V=\Spec B$ induces $B\xrightarrow[\Ring]{}A_g$. By \textit{(5.2e)}, $\FF|_{D(g)}=i^\ast(\FF|_V)\simeq i^\ast\widetilde{M}\simeq(M\otimes_BA_g)\widetilde{}$.
            \item Consequently, if $\FF$ is quasi-coherent on $X$, then $X$ can be covered by $D(g_i)$'s such that $\FF|_{D(g_i)}\simeq\widetilde{M}_i$ for some $M_i\in \MOD_{A_{g_i}}$. But, $X=\Spec A$ is quasi-compact, so we may choose finitely many $i$'s, say $1\leq i\leq r$.
        \end{itemize}
        \item[(II)] Proof of (a). For each $i$, set $s_i\coloneqq s|_{D(g_i)}\in\FF(D(g_i))=\Gamma(D(g_i),\widetilde{M}_i)\simeq M_i$. Note that on the intersection $D(f)\cap D(g_i)=D(fg_i)$, we have $s_i|_{D(f)}=s|_{D(fg_i)}=0$ in $(\FF|_{D(g_i)})(D(f))\simeq\widetilde{M}_i(D(f))\simeq(M_i)_f$, so $f^{n_i}s=0$ in $M_i=\FF(D(g_i))$ for some $n_i>0$. Here, $n_i$ depends on $i$, but there are only finitely many $i$'s, we may choose $n\geq\max\{n_i:1\leq i\leq r\}$ such that $f^ns_i=0$ for every $i$. Further, these $D(g_i)$'s cover $X$, and thus, $f^ns=0$.
        \item[(III)] Proof of (b). 
        \begin{itemize}
            \item Given $t\in\FF(D(f))$. For each $i$, set $t_i\coloneqq t|_{D(g_i)}\in\FF(D(fg_i))\simeq(M_i)_f$. By definition of localization, there is a section $\tilde{t}_i\in M_i\simeq\FF(D(g_i))$ such that $\tilde{t}_i|_{D(f)}=f^{n_i}t_i$ on $D(fg_i)$ for some $n_i>0$. Again, pick $n\geq\max\{n_i:1\leq i\leq r\}$ such that $\tilde{t}_i|_{D(f)}=f^nt_i$ for every $i$.
            \item Note that on the intersection $D(f)\cap D(g_ig_j)=D(fg_ig_j)$, we have
            \begin{align*}
                \tilde{t}_i|_{D(fg_ig_j)}=f^nt_i|_{D(g_j)}=f^nt|_{D(g_ig_j)}=f^nt_j|_{D(g_i)}=\tilde{t}_j|_{D(fg_ig_j)}.
            \end{align*}
            So, we have $s_{ij}\coloneqq(\tilde{t}_i-\tilde{t}_j)|_{D(g_ig_j)}\in\FF(D(g_ig_j))$ with $s|_{D(f)}=0$. By (a), there exists an $m_{ij}>0$ such that $f^{m_{ij}}s_{ij}=0$ on $D(g_ig_j)$. Again, take $m\geq\max\{m_{ij}:1\leq i,j\leq r\}$ such that $f^ms_{ij}=0$, i.e. $f^m\tilde{t}_i=f^m\tilde{t}_j$ on $D(g_ig_j)$. Now, gluing these sections $\{f^m\tilde{t}_i\}_{i=1}^r$ on $\{D(g_i)\}_{i=1}^r$, we obtain an $s\in\FF(X)$ with $s|_{D(f)}=f^{n+m}t$. (Indeed, on each $D(fg_i)$, we have $s=f^m\tilde{t}_i=f^{n+m}t_i=f^{n+m}t$; on each $D(fg_ig_j)$, we have $f^{n+m}t_i=f^m\tilde{t}_i=f^m\tilde{t}_j=f^{n+m}t_j$).
        \end{itemize}
    \end{itemize}
\end{proof}

\begin{proposition}{5.4}
    Let $(X,\OO_X)$ be a scheme. Then
    \begin{itemize}
        \item[(i)] $\FF\in\MOD_{\OO_X}$ is quasi-coherent $\Leftrightarrow$ For every open affine subset $U=\Spec A\subset X$,  we have $\FF|_U\simeq\widetilde{M}$ for some $M\in \MOD_A$.
        \item[(ii)] If $X$ is Noetherian. Then $\FF$ is coherent $\Leftrightarrow$ The same is true, with the extra condition that $M\in \MOD_A$ is finitely generated.
    \end{itemize}
\end{proposition}
\begin{proof}\mbox{}
    \begin{itemize}
        \item[(i)] $(\Leftarrow)$. Clear. $(\Rightarrow)$. Given an open affine subset $U=\Spec A\subset X$.
        \begin{itemize}
            \item By \textit{(5.3I)}, we may cover $X$ by $D(g_i)$ where $g_i\in B_i$ such that $\FF|_{D(g_i)}\simeq\widetilde{M}_i$ for some $M_i\in\MOD_{(B_i)_{g_i}}$. Since we may cover $U$ by some $D(g_i)\subset U$, it follows that $\FF_U$ is quasi-coherent. Thus, we may reduce to consider the affine case $X=\Spec A$.
            \item Consider $M\coloneqq\Gamma(X,\FF)$. By \textit{[5.3]}, there is a natural map $\alpha:\widetilde{M}\to\FF$. Since $\FF$ is quasi-coherent on $X=\Spec A$, we may cover $X$ by $D(g_i)\simeq\Spec A_{g_i}$ such that $\FF|_{D(g_i)}\simeq\widetilde{M}_i$ for some $M_i\in \MOD_{A_{g_i}}$. Now, 
            \begin{align*}
                M_i\underset{(5.1d)}{=}\Gamma(D(g_i),\widetilde{M}_i)\simeq\Gamma(D(g_i),\FF|_{D(g_i)})=\FF(D(g_i))\underset{(5.3)}{\simeq}M_{g_i}.
            \end{align*}
            Thus, 
            \begin{align*}
                \alpha_{D(g_i)}:\widetilde{M}(D(g_i))\underset{(5.1c)}{\simeq}M_i\xrightarrow[]{\sim}M_{g_i}\simeq\FF(D(g_i))
            \end{align*}
            is an isomorphism for each $i$, and hence, $\alpha:\widetilde{M}\xrightarrow{\sim}\FF$ is an isomorphism of $\OO_X$-modules.
        \end{itemize}
        \item[(ii)] $(\Leftarrow)$. Clear. $(\Rightarrow)$. In addition to (i), $M_i=M_{g_i}\in \MOD_{A_{g_i}}$ is finitely generated for each $i$. Now, since $\Spec A$ is Noetherian, $A$ is a Noetherian ring by \textit{(3.2)}, whence so is $A_{g_i}$ for each $i$ by \textit{(AM, 7.3)}. Further, $M_{g_i}\in \MOD_{A_{g_i}}$ is Noetherian for each $i$ by \textit{(AM, 6.5)}, and thus, $M\in \MOD_A$ is Noetherian as $A=(g_1,\cdots,g_n)$. Therefore, by \textit{(AM, 6.2)}, $M\in \MOD_A$ is finitely generated.
    \end{itemize}
\end{proof}

\begin{remark*}{}
    Note that being \textit{finitely generated} is actually a \textit{local property}. Thus, there is no need to assume $X$ is Noetherian in (ii).
\end{remark*}    

\begin{corollary}{5.5}
    Let $X=\Spec A$.
    \begin{itemize}
        \item[(i)] The functor $M\mapsto\widetilde{M}$ gives an equivalence of categories between the category $\MOD_A$ of $A$-modules and the category $\qcoh(X)$ of quasi-coherent $\OO_X$-modules. Its inversre is the functor $\Gamma(X,\cdot)$.
        \item[(ii)] If $A$ is Noetherian, the same functor also gives an equivalence of categories between the category $\fg\MOD_A$ of finitely generated $A$-modules and the category $\coh(X)$ of coherent $\OO_X$-modules.
    \end{itemize}
\end{corollary}
\begin{proof}
    (i) We note that
    \begin{itemize}
        \item For $M\in\MOD_A$, we have $\widetilde{M}\in\qcoh(X)$ by \textit{(5.1)}.
        \item For $\FF\in\qcoh(X)$, we have $\Gamma(X,\FF)\in\MOD_{\Gamma(X,\OO_X)}=\MOD_A$ by definition.
        \item By \textit{(5.1)}, we have $M\longmapsto\widetilde{M}\longmapsto\Gamma(X,\widetilde{M})=M$.
        \item By \textit{(5.4i)}, we have $\FF\longmapsto\Gamma(X,\FF)\longmapsto\Gamma(X,\FF)\widetilde{}\ \overset{\alpha}{\simeq}\ \FF$.
    \end{itemize}
    (ii) Similarly.
\end{proof}

\begin{proposition}{5.6}
    Let $X$ be an affine scheme, and let $0\to\FF'\xrightarrow{\phi}\FF\xrightarrow{\psi}\FF''\to0$ be an exact sequence of $\OO_X$-modules. If $\FF'\in\qcoh(X)$, then the sequence
    \begin{align*}
        0\longrightarrow\Gamma(X,\FF')\xrightarrow{\phi_X}\Gamma(X,\FF)\xrightarrow{\psi_X}\Gamma(X,\FF'')\longrightarrow0
    \end{align*}
    is exact.
\end{proposition}
\begin{notation}\mbox{}
    \begin{itemize}
        \item For each $k$, denote $U_k\coloneqq D(f_k)$; for each $i$, similarly, $U_i\coloneqq D(g_i)$; for the intersections, $U_{ki}\coloneqq D(f_kg_i)=D(f_k)\cap D(g_i)$, and similarly for $U_{ij}$ and $U_{ijk}$.
        \item For any $\FF\in\MOD_{\OO_X}$ and any section $s$, denote the restriction $s|_k\coloneqq s|_{U_k}\in\FF(U_k)$, and $s|_{ki},\ s|_{ij},\ s|_{ijk}$, similarly. 
        \item For any morphism $\FF\xrightarrow[\MOD_{\OO_X}]{\phi}\GG$, denote $\phi_k$ for the map $\phi_{U_k}:\FF(U_k)\to\GG(U_k)$, and similarly for $\phi_{ki},\ \phi_{ij},\ \phi_{ijk}$.
    \end{itemize}
\end{notation}

\begin{proof}
    From \textit{[1.8]}, $\Gamma(X,\cdot)$ is a left exact functor, it suffices to show that $\psi_X$ is surjective. Given $s''\in\FF''(X)$. Since $\psi$ is surjective, by \textit{[1.3a]}, there exists an open cover $\{D(f_k)\}$ for $X$, and for each $k$, there is a $t_k\in\FF(U_k)$ such that $s''|_k=\psi_k(t_k)$. But, $X$ is affine, thus quasi-compact, we may choose finitely many $k$'s, say $1\leq k\leq r$.\\ \newline
    \textbf{Step I.} For each $k$, there is an $n_k>0$ and a $\hat{t}_k\in\FF(X)$ for which $\psi(\hat{t}_k)=f^{n_k}_ks''$ in $\FF''(X)$.
    \begin{itemize}
        \item[(i)] As $X$ is affine, choose a finite open affine cover $\{D(g_i)\}_{i=1}^p$ for $X$ such that $s''|_i=\psi_i(t_i)$ for some $t_i\in\FF(U_i)$.
        \item[(ii)] Note that $\psi_{ik}=\psi_{ki}$ on its domain $U_{ki}=U_k\cap U_i$. So,
        \begin{align*}
            \psi_{ik}(t_k|_i)=\psi_{ki}(t_k|_i)=\psi_k(t_k)|_i=s''|_{ki}=s''|_{ik}=\psi_i(t_i)|_k=\psi_{ik}(t_i|_k),
        \end{align*}
        whence $t_k|_i-t_i|_k\in\ker\psi_{ik}=\Im\phi_{ik}$. By exactness, we have $\phi_{ik}(s'_{ik})=t_k|_i-t_i|_k$ for some $s'_{ik}\in\FF'(U_{ik})$.
        \item[(iii)] Since $\FF'\in\qcoh(X)$, by \textit{(5.3b)}, we have $f^{n_i}_ks'_{ik}\in\FF'(U_i)$ for some $n_i>0$. Take $n\geq\max\{n_i\}_{i=1}^p$ such that $f^n_ks'_{ik}\in\FF'(U_i)$ for all $i$. Now, set $u_i\coloneqq\phi_i(f^n_ks'_{ik})\in\FF(U_i)$ and $\tilde{t}_i\coloneqq f^n_kt_i+u_i\in\FF(U_i)$. Then, 
        \begin{align*}
            \psi_i(\tilde{t}_i)=f^n_k\psi_i(t_i)+\underbrace{\psi_i\phi_i(f^n_ks'_{ik})}_{=0\ \text{(Exactness)}}=f^n_ks''|_i.
        \end{align*}
        \item[(iv)] For each $i,j$, we have
        \begin{align*}
            \psi_{ij}(\tilde{t}_i|_j)=\psi_i(\tilde{t}_i)|_j=f^n_ks''|_{ij}=f^n_ks''|_{ji}=\psi_j(\tilde{t}_j)|_i=\psi_{ji}(\tilde{t}_j|_i)=\psi_{ij}(\tilde{t}_j|_i)
        \end{align*}
        on $U_{ij}=U_i\cap U_j$, whence $\tilde{t}_i|_j-\tilde{t}_j|_i\in\ker\psi_{ij}=\Im\phi_{ij}$. By exactness, we have $\phi_{ij}(\tilde{s}'_{ij})=\tilde{t}_i|_j-\tilde{t}_j|_i$ for some $\tilde{s}'_{ij}\in\FF'(U_{ij})$.
        \item[(v)] Note that 
        \begin{align*}
            \tilde{t}_i|_k=f^n_kt_i|_k+u_i|_k=f^n_kt_i|_k+\phi_{ik}(f^n_ks'_{ik})=f^n_k[t_i|_k+\phi_{ik}(s'_{ik})]\underset{(ii)}{=}f^n_kt_k|_i\tag{$\star$}
        \end{align*}
        on $U_{ik}=U_i\cap U_k$. Restricting to $U_{ijk}=U_{ij}\cap U_k$, we see that 
        \begin{align*}
            \phi_{ijk}(\tilde{s}'_{ij}|_k)=\phi_{ij}(\tilde{s}'_{ij})|_k=\tilde{t}_i|_{jk}-\tilde{t}_j|_{ik}=\tilde{t}_i|_{kj}-\tilde{t}_j|_{ki}\underset{(\star)}{=}f^n_kt_k|_{ij}-f^n_kt_k|_{ji}=0.
        \end{align*}
        By exactness, $\phi_{ijk}$ is injective, whence $\tilde{s}'_{ij}|_k=0$ on $U_{ij}\cap U_k$. Further, since $\FF'\in\qcoh(X)$, by \textit{(5.3a)}, we have $f^{m_{ij}}_k\tilde{s}'_{ij}=0$ in $\FF'(U_{ij})$ for some $m_{ij}>0$. As before, choose $m\geq\max\{m_{ij}:1\leq i,j\leq p\}$ for which $f^m_k\tilde{s}'_{ij}=0$ for all $i,j$. Then, 
        \begin{align*}
            f^m_k(\tilde{t}_i|_j-\tilde{t}_j|_i)=f^m_k\phi_{ij}(\tilde{s}'_{ij})=\phi_{ij}(f^m_k\tilde{s}'_{ij})=0,\ \text{i.e.}\ f^m_k\tilde{t}_i|_j=f^m_k\tilde{t}_j|_i\ \text{in}\ \FF(U_{ij}).
        \end{align*}
        Now, we may glue the sections $\{f^m_k\tilde{t}_i\}_{i=1}^p$ along $\{U_i\}_{i=1}^p$ to obtain $\hat{t}_k\in\FF(X)$.
        \item[(vi)] Check $\phi_X(\hat{t}_k)=f^{m+n}_ks''$ in $\FF''(X)$. For each $i$, we see that
        \begin{align*}
            \psi_X(\hat{t}_k)|_i=\psi_i(\hat{t}_k|_i)=\psi_i(f^m_k\tilde{t}_i)=f^m_k\psi_i(\tilde{t}_i)\underset{(iii)}{=}f^{m+n}_ks''|_i
        \end{align*}
        on $U_i$, whence $\phi_X(\hat{t}_k)=f^{m+n}_ks''$ on $X$. Here, such $m+n$ is the $n_k$ we required.
    \end{itemize}
    \textbf{Step II.} Finish the proof. Since there are finitely many $k$'s, we may pick $N\geq\max\{n_k\}_{k=1}^r$ such that $\psi_X(\hat{t}_k)=f^N_ks''$ in $\FF''(X)$ for all $k$. Further, $X=\bigcup_{k=1}^rD(f_k)=\bigcup_{k=1}^rD(f^N_k)$, so $\<f^N_1,\cdots,f^N_r\>=A$, whence $1=\sum_{k=1}^ra_kf^N_k$ for some $a_k\in A$. Pick $t\coloneqq\sum_{k=1}^ra_k\hat{t}_k\in\FF(X)$, then we have
    \begin{align*}
        \psi_X(t)=\sum_{k=1}^ra_k\psi_X(\hat{t}_k)\underset{(I)}{=}\sum_{k=1}^ra_kf^N_ks''=1\cdot s''=s''.
    \end{align*}
\end{proof}

\begin{proposition}{5.7}
    Let $X$ be a scheme.
    \begin{itemize}
        \item[(i)] The kernel, cokernel, and image of any morphism $\FF\xrightarrow[\qcoh(X)]{\phi}\GG$ are quasi-coherent.
        \item[(ii)] Any extension of quasi-coherent sheaves is quasi-coherent.
        \item[(iii)] If $X$ is Noetherian, then (i) and (ii) are true for coherent sheaves.
    \end{itemize}
\end{proposition}
\begin{proof}
    The question is local, we may assume $X=\Spec A$ is affine.
    \begin{itemize}
        \item[(i)] By \textit{(5.5)}, we have $\FF=\widetilde{M}$ and $\GG=\widetilde{N}$ for some $M,N\in \MOD_A$. By \textit{(5.2)}, the functor $M\mapsto\widetilde{M}$ is exact and fully faithful, so $\phi$ corresponds to a homomorphism $M\xrightarrow[\MOD_A]{\alpha}N$. Since $0\to\ker\alpha\to M\xrightarrow{\alpha}\Im\alpha\to0$ is exact, the sequence
        \begin{align*}
            0\xrightarrow[]{}\underbrace{(\ker\alpha)\widetilde{}}_{\simeq\ker\phi}\xrightarrow[]{}\widetilde{M}\xrightarrow{\phi}\underbrace{(\Im\alpha)\widetilde{}}_{\simeq\Im\phi}\xrightarrow[]{}0
        \end{align*}
        is exact. Similarly, $\Coker\phi\simeq(\Coker\alpha)\widetilde{}$. Therefore, they are quasi-coherent. 
        \item[(ii)] \textbf{Claim.} If $0\to\FF'\to\FF\to\FF''\to0$ is exact and $\FF',\FF''\in\qcoh(X)$, then $\FF\in\qcoh(X)$.
        
        By \textit{(5.6)}, 
        \begin{align*}
            0\longrightarrow\underbrace{\Gamma(X,\FF')}_{\coloneqq M'}\longrightarrow\underbrace{\Gamma(X,\FF)}_{\coloneqq M}\longrightarrow\underbrace{\Gamma(X,\FF'')}_{\coloneqq M''}\longrightarrow0
        \end{align*}
        is exact. Applying the exact functor $M\mapsto\widetilde{M}$, we have the commutative diagram
        \begin{center}
            \begin{tikzcd}
                0\rar[] &\widetilde{M}'\rar[]\dar["\theta'"] &\widetilde{M}\rar[]\dar["\theta"] &\widetilde{M}''\rar[]\dar["\theta''"] &0\\
                0\rar[] &\FF'\rar[] &\FF\rar[] &\FF''\rar[] &0
            \end{tikzcd}
        \end{center}
        with exact rows. Since $\FF',\ \FF''\in\qcoh(X)$, the maps $\theta'$ and $\theta''$ are isomorphisms. By \textit{Five Lemma}, $\theta$ is an isomorphism. Thus, $\FF\simeq\widetilde{M}\in\qcoh(X)$.
        \item[(iii)] If $X$ is Noetherian and $\FF',\FF''\in\coh(X)$. By \textit{(5.5)}, $M',M''\in\fg\MOD_A$. Further, $0\to M'\to M\to M''\to0$ is exact, so $M$ is finitely generated, whence $\FF=\widetilde{M}\in\coh(X)$.
    \end{itemize}
\end{proof}

\begin{proposition}{5.8}
    Let $f:X\to Y$ be a scheme morphism.
    \begin{itemize}
        \item[(a)] If $\GG\in\qcoh(Y)$, then $f^\ast\GG\in\qcoh(X)$.
        \item[(b)] If $X$ and $Y$ are Noetherian, and $\GG\in\coh(Y)$, then $f^\ast\GG\in\coh(X)$.
        \item[(c)] Assume that either (i) $X$ is Noetherian, or (ii) $f$ is quasi-compact and separated. If $\FF\in\qcoh(X)$, then $f_\ast\FF\in\qcoh(Y)$.
    \end{itemize}
\end{proposition}
\begin{proof}
    For (a) and (b), the questions are local, we may assume the morphism $X=\Spec A\xrightarrow[\Sch]{f}\Spec B=Y$ is induced by the homomorphism $B\xrightarrow[\Ring]{\phi}A$.
    \newline
    (a) If $\GG\in\qcoh(Y)$, by \textit{(5.5)}, $\GG=\widetilde{N}$ for some $N\in \MOD_B$. So, $f^\ast\GG=f^\ast(\widetilde{N})\underset{(5.2e)}{\simeq}(N\otimes_BA)\widetilde{}\in\qcoh(X)$ since $N\otimes_BA\in \MOD_A$.\\
    \newline
    (b) If both $X,Y$ are Noetherian, and $\GG\in\coh(Y)$, in addition to \textit{(a)}, we have $N\in\fg\MOD_B$. There exists an $n>0$ and a surjective map $B^{\otimes n}\twoheadrightarrow N$. Note that the functor $-\otimes_BA$ is right exact, we have $A^{\otimes n}\simeq B^{\otimes n}\otimes_BA\twoheadrightarrow N\otimes_BA$, i.e. $N\otimes_BA\in\fg\MOD_A$. Thus, $f^\ast\GG\simeq(N\otimes_BA)\widetilde{}\in\coh(X)$ by \textit{(5.5)}.\\
    \newline
    (c) This question is local on $Y$ only. 
    \begin{itemize}
        \item Note that if $V\subset Y$ is open affine, then $(f_\ast\FF)|_V$ depends on the restriction of $\FF$ to $f^{-1}(V)\subset X$. Replace $Y$ and $X$ by $V$ and $f^{-1}(V)$, respectively, we may assume $Y\coloneqq\Spec B$ is affine.
        \item In both case (i) and (ii), $X$ is quasi-compact, so we may cover $X$ by finitely many open affines $U_i$ ($1\leq i\leq r$).
        \item \textbf{Case (ii).} Since $f$ is separated, by \textit{[4.3]}, each $U_i\cap U_j$ ($1\leq i,j\leq r$) is affine. Denote them by $U_{ijk}\coloneqq U_i\cap U_j$, (here, $k=1$).
        \item \textbf{Case (i).} As $X$ is Noetherian, each $U_i\cap U_j$ ($1\leq i,j\leq r$) is quasi-compact. We can cover each $U_i\cap U_j$ with finitely many open affines $U_{ijk}$, (here, $k$ depends on both $i,j$).
        \item \textbf{Claim.} If $\FF\in\qcoh(X)$, then $f_\ast\FF\in\qcoh(Y)$.
        
        Recall from the \textit{Sheaf Property (4)} that if $V\subset Y$ is open, to give an $s\in\FF(f^{-1}V)$ is the same as to give a collection $\{s_i\}$ with $s_i\in\FF((f^{-1}V)\cap U_i)$ such that the restrictions on $f^{-1}\cap U_{ijk}$ are equal. So, we have the exact sequence of $\OO_Y$-modules:
        \begin{align*}
            0\longrightarrow f_\ast\FF\longrightarrow\bigoplus_{i=1}^rf_\ast(\FF|_{U_i})\xrightarrow{~\psi~~}\bigoplus_{i,j,k}f_\ast(\FF|_{U_{ijk}}).
        \end{align*}
        Since $\FF\in\qcoh(X)$, we see that 
        \begin{align*}
            \bigoplus_{i=1}^rf_\ast(\FF|_{U_i}) &\underset{\text{(Def.)}}{\simeq}\bigoplus_{i=1}^rf_\ast(\widetilde{M}_i) \underset{(5.2d)}{\simeq}\bigoplus_{i=1}^r\,_B\widetilde{M}_i \underset{(5.2c)}{\simeq}(\bigoplus_{i=1}^r\,_BM_i)\widetilde{}\in\qcoh(Y)
        \end{align*}
        for some $M_i\in \MOD_A$. Similarly, $\bigoplus_{i,j,k}f_\ast(\FF|_{U_{ijk}})\in\qcoh(Y)$. It follows that $\psi$ is a morphism of quasi-coherent sheaves on $Y$, whence $f_\ast\FF=\ker\psi\in\qcoh(Y)$ by \textit{(5.7a)}. 
    \end{itemize}
\end{proof}

\begin{definition}{(Ideal Sheaf)}
    Let $Y$ be a closed subscheme of $X$ and $i:Y\to X$ the inclusion morphism. The \textit{ideal sheaf} of $Y$ is defined to be $\II_Y\coloneqq\ker i^\#=\ker(\OO_X\to i_\ast\OO_Y)$.
\end{definition}

\begin{proposition}{5.9}
    Let $X$ be a scheme.
    \begin{itemize}
        \item[(i)] For any closed subscheme $Y$ of $X$, the ideal sheaf $\II_Y$ is a quasi-coherent sheaf of ideals on $X$.
        \item[(ii)] In addition to (i), if $X$ is Noetherian, then $\II_Y$ is coherent.
        \item[(iii)] Conversely, if $\JJ$ is a quasi-coherent sheaf of ideals on $X$, then $\JJ=\II_Y$ is the ideal sheaf of a uniquely determined closed subscheme $Y$ of $X$.
    \end{itemize}
\end{proposition}
\begin{proof}
    (i) Note that the inclusion morphism $i:Y\xhookrightarrow{}X$ is clearly quasi-compact, and separated (by \textit{(4.6a)} as $i$ is a closed immersion). Further, as $Y$ is a scheme, $\OO_Y\in\qcoh(Y)$, so it follows from \textit{(5.8c)} that $i_\ast\OO_Y\in\qcoh(X)$, whence $i^\#:\OO_X\xrightarrow[\qcoh(X)]{}i_\ast\OO_Y$ is a morphism of $\OO_X$-modules. By \textit{(5.7i)}, we have $\II_Y=\ker i^\#\in\qcoh(X)$.\\ \newline
    (ii) If $X$ is Noetherian, for any open affine subset $U=\Spec A\subset X,\ A$ is Noetherian. Thus, the ideal $I\coloneqq\Gamma(U,\II_Y|_U)\triangleleft\Gamma(U,\OO_X|_U)=A$ is finitely generated. By \textit{(5.5)}, $\widetilde{I}=\Gamma(U,\II_Y|_U)\widetilde{}=\II_Y|_U$. Thus, $\II_Y\in\coh(X)$.\\ \newline
    (iii) Given $\JJ\in\qcoh(X)$ with $\JJ\triangleleft\OO_X$. Consider $Y\coloneqq\Supp(\OO_X/\JJ)$, we claim that $(Y,\OO_X/\JJ)$ is the unique closed subscheme of $X$ with ideal sheaf $\JJ$. The uniqueness is clear. (Why?) Note that the question is local, so we may assume $X=\Spec A$ is affine. Since $\JJ\in\qcoh(X)$, we have $\JJ=\widetilde{\mathfrak{a}}$ for some ideal $\mathfrak{a}\triangleleft A$. Further, we see that
    \begin{align*}
        Y &=\Supp(\OO_X/\JJ)\underset{\text{(Def.)}}{=}\{P\in X:(\OO_X/\JJ)_P\neq0\}\underset{(5.1b)}{=}\{\pp\in X:(A/\mathfrak{a})_\pp\neq0\} \\ &=\Spec(A/\mathfrak{a})\overset{\text{homeom.}}{\simeq}V(\mathfrak{a}).
    \end{align*}
    Thus, $(Y,\OO_X/\JJ)$ is the closed subscheme of $X$ determined by $\mathfrak{a}$. (cf. \textit{3.2.3}).
\end{proof}

\begin{corollary}{5.10}
    If $X=\Spec A$ is an affine scheme.
    \begin{itemize}
        \item[(i)] There is a 1-1 correspondence between ideals $\mathfrak{a}\triangleleft A$ and closed subschemes $Y\subset X$, given by  $\mathfrak{a}\mapsto$ image of $\Spec(A/\mathfrak{a})$ in $X$.
        \item[(ii)] In particular, every closed subscheme of an affine scheme is affine.
    \end{itemize}
\end{corollary}
\begin{proof}
    By \textit{(5.5)}, the quasi-coherent sheaves of ideals on $X=\Spec A$ are 1-1 correspondence with the ideals of $A$.
\end{proof}

\begin{definition}{(Associated Sheaf on $\Proj S$)}
    Fix a graded ring $S$ and a graded $S$-module $M$. We define the \textit{sheaf associated} to $M$ on $\Proj S$, denoted by $\widetilde{M}$, as follows. For any open subset $U\subset\Proj S$, define
    \begin{align*}
        \widetilde{M}(U)\coloneqq\Big\{s:U\to\coprod_{\pp\in U}M_{(\pp)}\Big|\ s\ \text{is locally a fraction}\ \frac{m}{f}\ (m\in M,\ f\in A)\Big\}.
    \end{align*}
    This make $\widetilde{M}$ into a sheaf with obvious restriction maps.
\end{definition}

\begin{proposition}{5.11}
    Fix a graded ring $S$ and a graded $S$-module $M$. Let $X=\Proj S$, then
    \begin{itemize}
        \item[(a)] For any $\pp\in X$, the stalk $(\widetilde{M})_\pp=M_{(\pp)}$.
        \item[(b)] For any homogeneous $f\in S_+$, we have $\widetilde{M}|_{D_+(f)}\simeq\widetilde{M_{(f)}}$ (via $D_+(f)\simeq\Spec S_{(f)}$).
        \item[(c)] $\widetilde{M}\in\qcoh(X)$. If $S$ is Noetherian and $M$ is finiely generated, then $\widetilde{M}\in\coh(X)$. 
    \end{itemize}
\end{proposition}

\begin{proposition}{5.A}
    Fix a graded ring $S$ which is generated by $S_1$ as an $S_0$-algebra, and let $X=\Proj S$.
    \begin{itemize}
        \item[(a)] For any $M,N\in \GR\MOD_S$, and $f\in S_1$, we have $(M\otimes_SN)_{(f)}\simeq M_{(f)}\otimes_{S_{(f)}}N_{(f)}$.
        \item[(b)] Continue \textit{(a)}, we have $(M\otimes_SN)\widetilde{}\simeq\widetilde{M}\otimes_{\OO_X}\widetilde{N}$.
    \end{itemize}
    Let $T$ be another graded ring, generated by $T_1$ as a $T_0$-algebra. Let $\phi:S\to T$ be a graded homomorphism, let $U\subset Y=\Proj T$, and $f:U\xrightarrow[\Sch]{} X$ be determined by $\phi$ (as in \textit{[2.14]}).
    \begin{itemize}
        \item[(c)] For any $M\in \GR\MOD_S,\ f^\ast(\widetilde{M})\simeq(M\otimes_ST)\widetilde{}|_U$ (as $\OO_Y$-modules).
        \item[(d)] For any $N\in \GR\MOD_T,\ f_\ast(\widetilde{N}|_U)\simeq(\,_SN)\widetilde{}$ (as $\OO_X$-modules).
    \end{itemize}
\end{proposition}
\begin{proof}
    (a) Note that $M\otimes_SN$ is a graded $S$-module by $(M\otimes_SN)_d=\bigoplus_{d_1+d_2=d}(M_{d_1}\otimes_SN_{d_2})$. Fix $f\in S_1$, then we have the canonical inclusions $S_{(f)}\xhookrightarrow{}S_f,\ M_{(f)}\xhookrightarrow{}M_f$ and $N_{(f)}\xhookrightarrow{}N_f$, which induce a homomorphism $M_{(f)}\otimes_{S_{(f)}}N_{(f)}\to M_f\otimes_{S_f}N_f$. But, the canonical isomorphism $M_f\otimes_{S_f}N_f\xrightarrow{\sim}(M\otimes_SN)_f$ preserves the degrees, whence this yields a map
    \begin{align*}
        \phi_f:M_{(f)}\otimes_{S_{(f)}}N_{(f)}\longrightarrow(M\otimes_SN)_{(f)},\quad \frac{m}{f^{d_1}}\otimes\frac{n}{f^{d_2}}\longmapsto\frac{m\otimes n}{f^{d_1+d_2}}
    \end{align*}
    for every $m\in M_{d_1},\ n\in N_{d_2}$. It remains to show that $\phi_f$ is an isomorphism.
    \begin{itemize}
        \item \textbf{(Surjectivity).} Suppose $\frac{m\otimes n}{f^d}\in(M\otimes_SN)_{(f)}$, then $\deg(m\otimes n)=d$, say $m\otimes n\in M_{d_1}\otimes_SN_{d_2}$ (with $d_1+d_2=d$). Clearly, we have $\frac{m}{f^{d_1}}\otimes\frac{n}{f^{d_2}}\mapsto\frac{m\otimes n}{f^{d}}$. 
        \item \textbf{(Injectivity).} Suppose $\frac{m}{f^{d_1}}\otimes\frac{n}{f^{d_2}}\mapsto0$ where $m\in M_{d_1}$ and $n\in N_{d_2}$. Then, there exists an $r>0$ for which $f^r(m\otimes n)=0$ in $M\otimes_SN$. Note that if we consider $S_{(f)}$ as an $S$-algebra with the canonical homomorphism $S\to S_{(f)}$ sending $s\in S_d\mapsto\frac{s}{f^d}$, we may obtain another canonical homomorphism $\psi_f:M\otimes_SN\to M_{(f)}\otimes_{S_{(f)}}N_{(f)}$.
        Now, we see that
        \begin{align*}
            \frac{m}{f^{d_1}}\otimes\frac{n}{f^{d_2}}=\frac{f^rm}{f^{r+d_1}}\otimes\frac{n}{f^{d_2}}=\psi_f(f^rm\otimes n)=\psi_f(0)=0,
        \end{align*}
        whence $\phi_f$ is injective.
    \end{itemize}
    (b) Recall from \textit{(2.5b)} that $D_+(f)\simeq\Spec S_{(f)}$ is affine, so $\OO_X|_{D+(f)},\ \widetilde{M}|_{D+(f)}$ and $\widetilde{N}|_{D+(f)}$ are indeed (quasi-coherent) $\OO_{\Spec S_{(f)}}$-modules. It suffices to show that $(M\otimes_SN)\widetilde{}|_{D_+(f)}\simeq(\widetilde{M}\otimes_{\OO_X}\widetilde{N})|_{D_+(f)}$ for every $f\in S_+$. But, $S_1$ generates $S$ (thus, $S_+$) over $S_0$, whence it remains to show for every $f\in S_1$. Now, we see that
    \begin{align*}
        (M\otimes_SN)\widetilde{}|_{D_+(f)} &\underset{(5.11b)}{\simeq}((M\otimes_SN)_{(f)})\widetilde{} \underset{(a)}{\simeq}(M_{(f)}\otimes_{S_{(f)}}N_{(f)})\widetilde{} \underset{(5.2b)}{\simeq}\widetilde{M_{(f)}}\otimes_{\widetilde{S_{(f)}}}\widetilde{N_{(f)}}\\
        &\underset{(5.11b)}{\simeq}\widetilde{M}|_{D_+(f)}\otimes_{\OO_X|_{D_+(f)}}\widetilde{N}|_{D_+(f)}=(\widetilde{M}\otimes_{\OO_X}\widetilde{N})|_{D_+(f)}.
    \end{align*}
    (c) Note that $\{D_+(g):g\in\phi(S_+)\}$ forms an open affine cover of $U$ and that $D_+(g)\simeq\Spec T_{(g)}$ is affine. It suffices to show that $f^\ast(\widetilde{M})|_{D_+(g)}\simeq(M\otimes_ST)\widetilde{}|_{D_+(g)}$ for every $g\coloneqq\phi(h)\in\phi(S_+)\subset T_+$. Indeed, we see that
    \begin{align*}
        (M\otimes_ST)\widetilde{}|_{D_+(g)} &\underset{(5.11b)}{\simeq}\Big((M\otimes_ST)_{(g)}\Big)\widetilde{} \simeq(M_{(h)}\otimes_{S_{(h)}}T_{(g)})\widetilde{} \underset{(5.2e)}{\simeq}f^\ast(\widetilde{M_{(h)}})\\ &\underset{(5.11b)}{\simeq}f^\ast(\widetilde{M}|_{D_+(h)}) =f^\ast(\widetilde{M})|_{D_+(g)}.
    \end{align*}
    (d) Similar as \textit{(c)}, we claim that $f_\ast(\widetilde{N}|_U)|_{D_+(h)}\simeq(\,_SN)\widetilde{}|_{D_+(h)}$ for every $h\in S_+$. Indeed,
    \begin{align*}
        f_\ast(\widetilde{N}|_U)|_{D_+(h)} &=f_\ast(\widetilde{N}|_{D_+(\phi(h))}) \underset{(5.11b)}{\simeq}f_\ast(\widetilde{N_{(\phi(h))}}) \underset{(5.2d)}{\simeq}(\,_{S_{(h)}}N_{(\phi(h))})\widetilde{} =\Big((\,_SN)_{(h)}\Big)\widetilde{} \underset{(5.11b)}{\simeq}(\,_SN)\widetilde{}|_{D_+(h)}.
    \end{align*}
\end{proof}

\begin{definition}{(Twisted Sheaf)}
    Fix a graded ring $S$ and let $X=\Proj S$. Suppose $n\in\mathbb{Z}$.
    \begin{itemize}
        \item The sheaf $\OO_X(n)\coloneqq\widetilde{S(n)}$, where $S(n)\coloneqq\bigoplus_{d\geq n}S_d$. Call $\OO_X(1)$ the \textit{twisting sheaf of Serre}.
        \item For any $\FF\in\MOD_{\OO_X}$, denote $\FF(n)\coloneqq\FF\otimes_{\OO_X}\OO_X(n)$.
    \end{itemize}
\end{definition}

\begin{proposition}{5.12}
    Fix a graded ring $S$ which is generated by $S_1$ as an $S_0$-algebra, and let $X=\Proj S$.
    \begin{itemize}
        \item[(a)] The sheaf $\OO_X(n)$ is an invertible sheaf on $X$ for every $n\in\mathbb{Z}$.
        \item[(b)] For any $M\in \GR\MOD_S,\ \widetilde{M}(n)\simeq\widetilde{M(n)}$. Further, $\OO_X(n)\otimes\OO_X(m)\simeq\OO_X(n+m)$.
        \item[(c)] Let $T,\ \phi$ and $f$ be as in \textit{(5.A)}. Then,  $f^\ast(\OO_X(n))\simeq\OO_Y(n)|_U$ and $f_\ast(\OO_Y(n)|_U)\simeq(f_\ast(\OO_Y|_U))(n)$. 
    \end{itemize}
\end{proposition}
\begin{proof}
    (a) Note that $X$ is covered by $\{D_+(f):f\in S_1\}$ (as $S_1$ generates $S$, and thus, $S_+$ over $S_0$). Given $f\in S_1$, consider $\OO_X(n)|_{D_+(f)}=\widetilde{S(n)}|_{D_+(f)}=\widetilde{S(n)_{(f)}}$. Clearly, $S(n)_{(f)}$ is a free $S_{(f)}$-module of rank $1$, and that we have the canonical isomorphism
    \begin{align*}
        S_{(f)}\isolongto S(n)_{(f)},\quad\frac{a}{f^d}\longmapsto f^{n-d}a\in S_{n-d}\cdot S_d\subset S_n=S(n)_0
    \end{align*}
    for any $d\geq0$ and $a\in S_d$. So, $\OO_X(n)|_{D_+(f)}\simeq\widetilde{S_{(f)}}\simeq\OO_X|_{D_+(f)}$ is a free $\OO_X|_{D_+(f)}$-module of rank $1$, and therefore, $\OO_X(n)$ is a locally free sheaf of rank $1$, i.e. an invertible sheaf on $X$.\\ \newline
    (b) For the first isomorphism,
    \begin{align*}
        \widetilde{M}(n) &\underset{\text{(Def.)}}{=}\widetilde{M}\otimes_{\OO_X}\OO_X(n) =\widetilde{M}\otimes_{\widetilde{S}}\widetilde{S(n)} \underset{(5.Ab)}{\simeq}(M\otimes_SS(n))\widetilde{} \simeq M(n)\widetilde{}. 
    \end{align*}
    For the second, we have
    \begin{align*}
        \OO_X(n)\otimes\OO_X(m)\simeq\widetilde{S(n)}\otimes\widetilde{S(m)} \underset{(5.Ab)}{\simeq}(S(n)\otimes S(m))\widetilde{} =S(n+m)\widetilde{}=\OO_X(n+m).
    \end{align*}
    (c) First, we have
    \begin{align*}
        f^\ast(\OO_X(n))=f^\ast(\widetilde{S(n)}) \underset{(5.Ac)}{\simeq}(S(n)\otimes_ST)\widetilde{}|_U=T(n)\widetilde{}|_U \underset{(b)}{=}\OO_Y(n)|_U.
    \end{align*}
    For the latter, 
    \begin{align*}
        f_\ast(\OO_Y(n)|_U) =f_\ast(\widetilde{T(n)}|_U) \underset{(5.Ad)}{\simeq}(\,_ST(n))\widetilde{} \underset{(b)}{=}\widetilde{(\,_ST)}(n) \underset{(5.Ad)}{\simeq}(f_\ast(\widetilde{T}|_U))(n) =(f_\ast(\OO_Y|_U))(n).
    \end{align*}
\end{proof}

\begin{proposition}{5.B}
    The twisting functor $\FF\mapsto\FF(n)$ is an exact functor between the category $\MOD_{\OO_X}$ for every $n\in\mathbb{Z}$ where $X=\Proj S$ and $S$ is a graded ring generated by $S_1$ as an $S_0$-algebra.
\end{proposition}
\begin{proof}
    Fix $n\in\mathbb{Z}$. Suppose $0\to\FF'\to\FF\to\FF''\to0$ is a short exact sequence of $\OO_X$-modules, we show that the sequence $0\to\FF'(n)\to\FF(n)\to\FF''(n)\to0$ is exact.
    \begin{itemize}
        \item For each $P\in X$, the sequence $0\to\FF'_P\to\FF_P\to\FF''_P\to0$ of stalks is exact. Since the tensor functor $-\otimes_{\OO_{X,P}}\OO_X(n)_P$ (in the category $\MOD_{S_{(P)}}$) is right exact, we have
        \begin{center}
            \begin{tikzcd}
                \FF'_P\otimes_{\OO_{X,P}}\OO_X(n)_P\rar[]\dar[dash,"\verteq"] &\FF_P\otimes_{\OO_{X,P}}\OO_X(n)_P\rar[]\dar[dash,"\verteq"] &\FF''_P\otimes_{\OO_{X,P}}\OO_X(n)_P\rar[]\dar[dash,"\verteq"] &0\\
                \FF'(n)_P\rar[] &\FF(n)_P\rar[] &\FF''(n)_P\rar[] &0.
            \end{tikzcd}
        \end{center}
        Therefore, $\FF'(n)\to\FF(n)\to\FF''(n)\to0$ is exact.
        \item To show the left exactness $0\to\FF'(n)\to\FF(n)$, we have to show that the sheaf morphism $\FF'(n)\to\FF(n)$ is injective. 
        %But this is equivalent to show that $\FF'(n)|_{D_+(f)}\to\FF(n)|_{D_+(f)}$ is injective for every $f\in S$ (in fact, for every $f\in S_1$ since $S_1$ generates $S$ over $S_0$). 
        Indeed, for each $f\in S_1$, $\OO_X(n)|_{D_+(f)}\simeq\OO_X|_{D_+(f)}$ is a free $\OO_X|_{D+(f)}$-module by \textit{(5.12a)}, so $\OO_X(n)(D_+(f))\simeq\OO_X(D_+(f))$ is a free (thus, flat) $\OO_X(D_+(f))$-module. It follows that
        \begin{align*}
            \FF'(n)(D_+(f))\simeq\FF'\otimes\OO_X(D_+(f))\longrightarrow\FF\otimes\OO_X(D_+(f))\simeq\FF(n)(D_+(f))
        \end{align*}
        is injective when viewing $\FF'(n)$ and $\FF(n)$ as presheaves, and thus, is injective when viewing them as sheaves by \textit{[1.4a]}. As this is true for every $f\in S_1$, the fact that these $f\in S_1$ generates $S$ over $S_0$ yields the injection $\FF'(n)\to\FF(n)$ 
    \end{itemize}
\end{proof}

\begin{remark*}{}
    The twisting operation above allows us to define a graded $\OO_X$-module on $X=\Proj S$.
\end{remark*}

\begin{definition*}{}\mbox{}
    Fix a graded ring $S$, let $X=\Proj S$, and let $\FF\in\MOD_{\OO_X}$. We define the \textit{graded $S$-module associated} to $\FF$ as a group to be $\Gamma_\ast(\FF)\coloneqq\bigoplus_{n\in\mathbb{Z}}\Gamma(X,\FF(n))$. We give it a structure of graded $S$-module as follows:\\ 
    If $s\in S_d$, then $s$ determines (in a natural way) a global section $s\in\Gamma(X,\OO_X(d))$. Then, for any $t\in\Gamma(X,\FF(n))$, we define $s\cdot t\in\Gamma(X,\FF(n+d))$ by taking the tensor products $s\otimes t$ via the natural isomorphism $\FF(n)\otimes_{\OO_X}\OO_X(d)\simeq\FF(n+d)$.
\end{definition*}

\begin{proposition}{5.13}
    Suppose $S=A[x_0,\cdots,x_r]$, ($r\geq1$) and let $X=\Proj S=\PP^r_A$. Then $\Gamma_\ast(\OO_X)\simeq S$.
\end{proposition}
\begin{proof}\mbox{}
    \begin{itemize}
        \item Cover $X$ by $\{D_+(x_i)\}_{i=0}^r$. Recall from the \textit{sheaf property} that to give a $t\in\Gamma(X,\OO_X(n))$ is the same as to give a $t_i\in\Gamma(D_+(x_i),\OO_X(n))$ for each $i$ such that $t_i|_{D_+(x_ix_j)}=t_j|_{D_+(x_ix_j)}$ for every $i,j$.
        \item Note that if $f\in S_d$ ($d>0$), then $S(nd)_{(f)}=(S(nd)_f)_0=(S_f(n))_0=(S_f)_{nd}$; in particular, $S(n)_{(x_i)}=(S_{x_i})_n$ for each $i$. It follows that $\OO_X(n)|_{D_+(x_i)}=S(n)\widetilde{}|_{D_+(x_i)}=(S(n)_{(x_i)})\widetilde{}=(S_{x_i,n})\widetilde{}$. Now, we see that
        \begin{align*}
            t_i\in\Gamma(D_+(x_i),\OO_X(n))=\Gamma(D_+(x_i),\OO_X(n)|_{D_+(x_i)})=\Gamma(\Spec S_{(x_i)},(S_{x_i,n})\widetilde{}) \underset{(5.1d)}{=}(S_{x_i})_n\subset S_{x_i}.
        \end{align*}
        Similarly, its restriction $t_i|_j\coloneqq t_i|_{D_+(x_j)}\in\Gamma(D_+(x_ix_j),\OO_X(n))=(S_{x_ix_j})_n\subset S_{x_ix_j}$.
        \item Summing over all $n\in\mathbb{Z}$, we have
        \begin{align*}
            \Gamma_\ast(\OO_X)=\bigoplus_{n\in\mathbb{Z}}\Gamma(X,\OO_X(n)) =\{(t_0,\cdots,t_r):t_i\in S_{x_i}\ \text{with $t_i|_j=t_j|_i$ for every}\ i,j\}.
        \end{align*}
        \item For each $i$, since $x_i$ is not a zero divisor of $S=A[x_0,\cdots,x_r]$, the natural localization maps
        \begin{align*}
            S\to S_{x_i},\ s\mapsto\frac{s}{1}\quad S_{x_i}\to S_{x_ix_j},\ \frac{s}{x_i^n}\mapsto\frac{s}{(x_ix_j)^n}
        \end{align*}
        are injective. Thus, these rings, $S,\ S_{x_i},S_{x_ix_j},\cdots$, are all subrings of $S'\coloneqq S_{x_0\cdots x_r}$, whence $\Gamma_\ast(\OO_X)=\bigcap_{i=0}^rS_{x_i}\subset S'$.
        \item Note. Every homogeneous element $g\in S'$ can be written as $g\coloneqq x_0^{d_0}\cdots x_r^{d_r}f(x_0,\cdots,x_r)$ for some $d_k\in\mathbb{Z}$ and $f\in S^h$ with $x_k\nmid f$ for every $k$. It follows that $g\in S_{x_i}\ \Leftrightarrow\ d_k\geq0$ for every $k\neq i$. Therefore, we have
        \begin{align*}
            g\in\Gamma_\ast(\OO_X)\ \Leftrightarrow\ g\in S_{x_i}\ \text{for all}\ 0\leq i\leq r\ \Leftrightarrow\ d_i\geq0\ \text{for all}\ 0\leq i\leq r\ \Leftrightarrow\ g\in S
        \end{align*}
        for any homogeneous element $g$, whence $\Gamma_\ast(\OO_X)\simeq S$.
    \end{itemize}
\end{proof}

\begin{lemma}{5.14}
    Fix a scheme $X$, let $\LL$ be an invertible sheaf on $X$, let $f\in\Gamma(X,\LL)$, let $X_f=\{x\in X:f_x\notin\mm_x\LL_x\}$ be the open set, and let $\FF\in\qcoh(X)$.
    \begin{itemize}
        \item[(a)] Suppose $X$ is quasi-compact, and $s\in\Gamma(X,\FF)$ with restriction $s|_{X_f}=0$. Then $f^ns=0$ in $\Gamma(X,\FF\otimes\LL^{\otimes n})$ for some $n>0$.
        \item[(b)] Suppose further that $X$ has a finite open affine cover $U_i$ such that $\LL|_{U_i}$ is free for each $i$, and such that $U_i\cap U_j$ is quasi-compact for every $i,j$. Given $t\in\Gamma(X_f,\FF)$, then $f^nt\in\Gamma(X_f,\FF\otimes\LL^{\otimes n})$ extends to a global section $s\in\Gamma(X,\FF\otimes\LL^{\otimes n})$ for some $n>0$.
    \end{itemize}
\end{lemma}
\begin{proof}\mbox{}
    \begin{itemize}
        \item General settings.
        \begin{itemize}
            \item Choose a finite open affine cover $U_i=\Spec A_i$ ($1\leq i\leq r$) of $X$ such that each $\LL|_{U_i}$ is a free $\OO_X$-module (as $X$ is quasi-compact and $\LL$ is locally free).
            \item Since $\LL$ is invertible, let $\psi_i:\LL|_{U_i}\xrightarrow[\MOD_{\OO_X}]{\sim}\OO_X|_{U_i}=\OO_{U_i}$ be the isomorphism.
            \item Since $\FF\in\qcoh(X)$, by \textit{(5.4)}, we have $\FF|_{U_i}\simeq\widetilde{M}_i$ for some $M_i\in \MOD_{A_i}$.
            \item For each $i$, let $g_i\coloneqq\psi_{i,U_i}(f|_i)\in\Gamma(U_i,\OO_{U_i})=A_i$ where $f|_i$ denotes $f|_i\coloneqq f|_{U_i}$ as usual. Then $U_i\cap X_f=D(g_i)$ by \textit{[2.16a]}.
        \end{itemize}
        \item Proof of (a). For each $1\leq i\leq r$, consider $s|_i\in\Gamma(U_i,\FF)=\Gamma(\Spec A_i,\widetilde{M}_i)=M_i$, then we have $(s|_i)|_{X_f}=(s|_{X_f})|_i=0$ on $D(g_i)$. By \textit{(5.3a)}, we have $g_i^{n_i}s|_i=0$ in $\FF(U_i)=M_i$ for some $n_i>0$. Using the isomorphism $\mathds{1}_\FF\times\psi_i^{\otimes n}:\FF\otimes\LL^{\otimes n}|_{U_i}\isoto\FF\otimes\OO_X^{\otimes n}|_{U_i}\simeq\FF|_{U_i}$, we see that
            \begin{align*}
                (\mathds{1}_\FF\times\psi_i^{\otimes n})_{U_i}:M\otimes\LL^{\otimes n}(U_i) &\isolongto M_i\otimes A_i^{\otimes n}\simeq M_i,\\
                s|_i\cdot f^{n_i}|_i &\longmapsto s|_i\cdot g_i^{n_i}=0.
            \end{align*}
            By injectivity, we have $(f^{n_i})s|_i=0$ in $(\FF\otimes\LL^{\otimes n})(U_i)$ for every $1\leq i \leq r$. Pick $n\geq\max\{n_i:1\leq i\leq r\}$ for which $(f^ns)|_i=0$ on $U_i$, then we obtain $f^ns=0$ on $X$.
        \item Proof of (b).
        \begin{itemize}
            \item For each $i$, consider $t|_i=t|_{U_i}\in\FF(X_f\cap U_i)=\FF(D(g_i))\simeq(M_i)_{g_i}$. By definition of localization, there exists a $\tilde{t}_i\in M_i=\FF(U_i)$ such that $\tilde{t}_i|_{X_f}=g_i^{n_i}t|_i$ in $\Gamma(D(g_i),\FF\otimes\LL_{U_i}^{\otimes n_i})$ for some $n_i>0$. Pick $n\geq\max\{n_i:1\leq i\leq r\}$, then we have the natural isomorphism
            \begin{align*}
                \FF\otimes\LL^{\otimes n}|_{U_i}\xrightarrow[\sim]{\mathds{1}_\FF\times\psi_i}\FF\otimes\OO_{U_i}^{\otimes n}\simeq\FF\otimes\OO_{U_i}^{\otimes n_i}\simeq\FF|_{U_i},
            \end{align*}
            whence $\tilde{t}_i|_{X_f}=g_i^nt|_i=\psi_i(f^n|_i)\cdot t|_i$ in $\Gamma(D(g_i),\FF\otimes\OO_{U_i}^{\otimes n})$
            \item Now, on the intersection $D(g_ig_j)=(U_i\cap U_j)\cap X_f$, we have
            \begin{align*}
                \tilde{t}_i|_{D(g_ig_j)}=(\psi_i(f^n)\cdot t|_i)|_j=\psi_{ij}(f^n|_{ij})\cdot t|_{ij}=\psi_{ji}(f^n|_{ji})\cdot t|_{ji}=\cdots=\tilde{t}_j|_{D(g_ig_j)}.
            \end{align*}
            It follows that $s_{ij}\coloneqq\tilde{t}_i|_j-\tilde{t}_j|_i\in\FF(U_i\cap U_j)$ satisfies $s_{ij}|_{X_f}=0$. But $U_i\cap U_j$ is quasi-compact, so we have $f^{m_{ij}}s_{ij}=0$ on $U_i\cap U_j$ for some $m_{ij}>0$ by \textit{(a)}. Again, we can take $m\geq\max\{m_{ij}:1\leq i,j\leq r\}$ for which $f^ms_{ij}=0$, i.e. $f^m\tilde{t}_i|_j=f^m\tilde{t}_j|_i$ in $\Gamma(X,\FF\otimes\LL^{\otimes m})$.
            \item Finally, we may glue the sections $\{f^m\tilde{t}_i\}_{i=1}^r$ on $\{U_i\}_{i=1}^r$ to obtain a global section $s\in\Gamma(X,\FF\otimes\LL^{\otimes (n+m)})$ with $s|_{X_f}=f^{n+m}t$. (Check!)
            \item Indeed, we have $(s|_{X_f})|_i=(s|_{U_i})|_{X_f}=f^m\tilde{t}_i|_{X_f}=f^{m+n}t|_i$ on $U_i$ and 
            \begin{align*}
                (s|_{X_f})|_{ij}=(f^m\tilde{t}_i|_{X_f})|_j=(f^m\tilde{t}_i|_j)|_{X_f}=(f^m\tilde{t}_j|_i)|_{X_f}=(f^m\tilde{t}_j|_{X_f})|_i=(s|_{X_f})|_{ji}
            \end{align*}
            on the intersection $U_i\cap U_j$.
        \end{itemize}
    \end{itemize}
\end{proof}

\begin{proposition}{5.15}
    Fix a graded ring $S$, which is finitely generated by $S_1$ as an $S_0$-algebra, and let $X=\Proj S$. If $\FF\in\qcoh(X)$, then there is a natural isomorphism $\beta:\Gamma_\ast(\FF)\widetilde{}\isoto\FF$.
\end{proposition}
\begin{proof}\mbox{}
    \begin{itemize}
        \item[(I)] Define $\beta:\Gamma_\ast(\FF)\widetilde{}\to\FF$. Note that $\Gamma_\ast(\FF)\widetilde{}\in\qcoh(X)$, and that $X$ has an open affine cover $D_+(f)\simeq\Spec S_{(f)}$ (with $f\in S_1$), whence we have
        \begin{align*}
            \Gamma_\ast(\FF)\widetilde{}|_{D_+(f)} &\simeq\Gamma_\ast(\FF)_{(f)}\simeq\bigoplus_{n\in\mathbb{Z}}\Gamma(X,\FF(n))_{(f)}\simeq\bigoplus_{n\in\mathbb{Z}}\FF|_{D_+(f)}(n)\\ &\simeq\bigoplus_{n\in\mathbb{Z}}\Gamma(D_+(f),\FF|_{D_+(f)}(n))\widetilde{} \simeq\Gamma_\ast(\FF|_{D_+(f)})\widetilde{}.
        \end{align*}
        Further, by \textit{[5.3]} we have 
        \begin{align*}
            \Hom_{\OO_X|_{D_+(f)}}(\Gamma_\ast(\FF|_{D_+(f)})\widetilde{},\FF|_{D_+(f)})\simeq\Hom_{S_{(f)}}(\Gamma_\ast(\FF|_{D_+(f)}),\Gamma(D_+(f),\FF|_{D_+(f)})),
        \end{align*}
        so it suffices to define a map
        \begin{align*}
            \alpha_f:\Gamma_\ast(\FF)_{(f)}=\Gamma_\ast(\FF)\widetilde{}(D_+(f))=\Gamma_\ast(\FF|_{D_+(f)})\to\Gamma(D_+(f),\FF)).
        \end{align*}
         for every $f\in S_1$. Indeed, we define 
         \begin{align*}
             \frac{m}{f^d}\longmapsto m\otimes f^{-d}\in\Gamma(D_+(f),\FF(d)\otimes\OO_X(-d))=\Gamma(D_+(f),\FF)
         \end{align*}
         for any $d\geq0$ and any  $m\in\Gamma(X,\FF(d))$.
         \item[(II)] Claim that such $\alpha_f$ is an $S_{(f)}$-module isomorphism.
         \begin{itemize}
             \item Consider the invertible sheaf $\LL\coloneqq\OO_X(1)$ on $X$, then $f\in S(1)=\Gamma(X,S(1)\widetilde{})=\Gamma(X,\LL)$.
             \item By assumption, say $S=S_0[f_0,\cdots,f_r]$ (where $f_i\in S_1$), thus, $X$ has a finite open affine cover $D_+(f_i)\simeq\Spec S_{(f_i)}$ for which every $\LL|_{D_+(f_i)}$ is a free $\OO_X$-module.
             \item The intersections $D_+(f_i)\cap D_+(f_j)=D_+(f_if_j)\simeq\Spec S_{(f_if_j)}$ are affine, thus, quasi-compact.
             \item Note that $\LL^{\otimes n}=\OO_X(1)^{\otimes n}\underset{(5.12b)}{=}\OO_X(n)$ for every $n\in\mathbb{Z}$.
             \item \textbf{(Injectivity).} If $\frac{m}{f^d}\mapsto0$, it follows from \textit{(5.14a)} that $f^n=0$ in $\Gamma(X,\FF\otimes\LL^{\otimes n})=\Gamma(X,\FF(n))$ for some $n>0$, whence $\frac{m}{f^d}=\frac{f^nm}{f^{n+d}}=0$ in $\Gamma_\ast(\FF)_{(f)}$.
             \item \textbf{(Surjectivity).} Given $t\in\Gamma(D_+(f),\FF)$. By \textit{(5.14b)}, there exist an $n>0$ and a global section $s\in\Gamma(X,\FF\otimes\LL^{\otimes n})$ for which $s|_{D_+(f)}=f^nt$ in $\Gamma(D_+(f),\FF\otimes\LL^{\otimes n})=\Gamma(D_+(f),\FF(n))$. Thus, $\frac{s}{f^n}\mapsto t$.
         \end{itemize}
    \end{itemize}
\end{proof}

\begin{corollary}{5.16}
    Fix a ring $A$.
    \begin{itemize}
        \item[(a)] If $Y$ is a closed subscheme of $\PP_A^r$, then there is a homogeneous ideal $I\triangleleft S=A[x_0,\cdots,x_r]$ such that $Y$ is the closed subscheme determined by $I$.
        \item[(b)] A scheme $Y$ over $\Spec A$ is projective $\Longleftrightarrow\ Y\simeq\Proj S$ for some graded ring $S$, which is finitely generated by $S_1$ as an $S_0$-algebra with $S_0=A$.
    \end{itemize}
\end{corollary}
\begin{proof}
    (a) Consider the inclusion morphism $Y\xhookrightarrow[\Sch]{(i,i^\#)}\PP_A^r$, which is a closed immersion. Let $\II_Y$ be the ideal sheaf of $Y$ on $X=\PP_A^r$, i.e. $\II_Y=\ker i^\#\triangleleft\OO_X$. Consider the short exact sequence of $\OO_X$-modules  $0\to\II_Y\xhookrightarrow{}\OO_X\overset{i^\#}{\twoheadrightarrow}i_\ast\OO_Y\to0$ (where $i^\#$ is surjective since $i$ is a closed immersion). For any $n\in\mathbb{Z}$, note that the twisting functor $\FF\mapsto\FF(n)$ is exact by \textit{(5.B)}, we have $0\to\II_Y(n)\to\OO_X(n)\to(i_\ast\OO_Y)(n)\to0$. Then, the left exact functor $\Gamma(X,-)$ gives $0\to\Gamma(X,\II_Y(n))\to\Gamma(X,\OO_X(n))\to\Gamma(X,i_\ast\OO_Y)$. Taking direct sums, we get $0\to\Gamma_\ast(\II_Y)\to\Gamma_\ast(\OO_X)\to\Gamma_\ast(i_\ast\OO_Y)$. Note that $I\coloneqq\Gamma_\ast(\II_Y)$ is a submodule, and hence, an ideal of $\Gamma_\ast(\OO_X)\underset{(5.13)}{\simeq}S$. By \textit{[3.12]}, $I$ determines a closed subscheme of $X=\Proj S$, with sheaf of ideals $\widetilde{I}$. But, $\II_Y\in\qcoh(X)$ by \textit{(5.9)}, whence $\II_Y\underset{(5.15)}{\simeq}\Gamma_\ast(\II_Y)\widetilde{}=\widetilde{I}$.\\
    \newline
    (b) We see that
    \begin{itemize}
        \item ($\Rightarrow$) If $Y$ is projective over $\Spec A$, the morphism $Y\xrightarrow[\Sch]{}\Spec A$ factors through $Y\xhookrightarrow{i}\PP_A^r\to\Spec A$ for some $r\geq0$ where $i$ is a closed immersion. By \textit{(a)}, $Y\simeq\Proj S/I$ for some homogeneous ideal $I\triangleleft S\coloneqq A[x_0,\cdots,x_r]$. Now, we choose $I'\coloneqq\bigoplus_{d>0}(I\cap S_d)\subset S_+$, then $(S/I')_d=(S/I)_d$ for every $d>0$ and $(S/I')_0=A$; in particular, $S/I'$ is finitely generated by $(S/I')_1$ over $A$ \textbf{(Why?)}. Thus, it follows from \textit{[3.12]} that $Y\simeq\Proj S/I\simeq\Proj S/I'$. 
        
        \item ($\Leftarrow$) For such graded ring $S$, we have $S\simeq A[x_0,\cdots,x_r]/I$ for some $r\geq0$ and some homogeneous ideal $I\triangleleft A[x_0,\cdots,x_r]$. Now, the natural quotient map $A[x_0,\cdots,x_r]\underset{\Ring}{\twoheadrightarrow}A/[x_0,\cdots,x_r]$ induces the morphism $Y=\Proj S\xrightarrow[\Sch]{}\Proj A[x_0,\cdots,x_r]=\PP_A^r$, which is a closed immersion by \textit{[3.12a]}. Therefore, $Y\to\PP_A^r\to\Spec A$ is a projective morphism.
    \end{itemize}
\end{proof}

\begin{definition}{(Twisting Sheaf on $\PP_Y^r$)}
    For any scheme $Y$, we define the \textit{twisting sheaf}$\OO(1)$ on $\PP_Y^r$ to be $g^\ast(\OO(1))$, where $g:\PP_Y^r=\PP_\mathbb{Z}^r\times_\mathbb{Z}Y\to\PP_\mathbb{Z}^r$ is the natural map.
\end{definition}

\begin{definition*}\newline
    \begin{itemize}
        \item \textbf{(Immersion).} A morphism $X\xrightarrow[\Sch]{i}Z$ is an \textit{immersion} if $i$ gives an isomorphism of $X$ with an open subscheme of a closed subscheme of $Z$, i.e. $i:X\isoto U\xhookrightarrow{\text{open}}Y\xhookrightarrow{\text{closed}}Z$.
        \item \textbf{(Very Ample Sheaf).} Let $X$ be a scheme over $Y$. An invertible sheaf $\LL$ on $X$ is \textit{very ample} relative to $Y$ if there is an immersion $i:X\to\PP_Y^r$ for some $r$, such that $i^\ast(\OO_{\PP_Y^r}(1))\simeq\LL$.
    \end{itemize}
\end{definition*}
\begin{remark*}{}
    Note that if $Y=\Spec A$, then the $\OO(1)$ is just the $\OO_X(1)$ defined on $X=\PP_A^r=\Proj A[x_0,\cdots,x_r]$ before, by \textit{(5.12c)}.
\end{remark*}

\begin{remark}{5.16.1}
    Let $Y$ be a Noetherian scheme. Then a $Y$-scheme $X$ is projective $\Longleftrightarrow$ it is proper, and there exists a very ample sheaf on $X$ relative to $Y$.
\end{remark}
\begin{proof}\mbox{}
    \begin{itemize}
        \item ($\Rightarrow$) If $f:X\xrightarrow[\Sch]{} Y$ is projective, then $f$ is proper by \textit{(4.9)}. Further, from the definition of projective morphism, there is a closed immersion (thus, an immersion) $X\xhookrightarrow{i}\PP_Y^r$ for some $r$ that factors through $f$. Then, $i^\ast(\OO(1))$ gives a very ample invertible sheaf on $X$ relative to $Y$.
        \item ($\Leftarrow$) If $f:X\xrightarrow[\Sch]{} Y$ is proper, and $\LL$ is a very ample invertible sheaf on $X$, then $\LL\simeq i^\ast(\OO(1))$ for some immersion $i:X\xhookrightarrow{}\PP_Y^r$. But $i(X)$ is closed by \textit{[4.4]}, so in fact $i$ is a closed immersion, whence $f:X\xhookrightarrow{i}\PP_Y^r\to Y$ is projective.
    \end{itemize} 
\end{proof}

\begin{definition*}{}
    Fix $X\in\Sch$ and $\FF\in\MOD_{\OO_X}$. We say that $\FF$ is \textit{generated by global sections} if there is a family of global sections $\{s_i\}_{i\in I}$, where $s_i\in\Gamma(X,\FF)$, such that for each $x\in X$, there germs $(s_i)_x\in\FF_x$ generate the stalk $\FF_x$ as an $\OO_{X,x}$-module, i.e. $\FF_x=\OO_{X,x}\spn\{(s_i)_x:i\in I\}$.
\end{definition*}

\begin{nt}{}
    $\FF\in\MOD_{\OO_X}$ is generated by global sections $\Longleftrightarrow\ \FF$ can by written as a quotient of a free $\OO_X$-module.
\end{nt}
\begin{proof}
    Suppose $s_i\in\Gamma(X,\FF)$ ($i\in I$), we see that
    \begin{align*}
        \FF\ \text{is generated by}\ \{s_i\}_{i\in I}
        &\Longleftrightarrow\ \FF_x=\OO_{X,x}\spn\{(s_i)_x:i\in I\}\ \text{for every}\ x\in X\\ &\Longleftrightarrow\ \OO_{X,x}^{\oplus I}\xrightarrow[]{\phi_x}\FF_x\to0\ \text{is exact for every}\ x\in X\\ &\Longleftrightarrow\ \OO_X^{\oplus I}\xrightarrow[]{\phi}\FF\to0\ \text{is exact}\\ &\underset{[1.7]}{\Longleftrightarrow}\ \FF\simeq\OO_X^{\oplus I}/\ker\phi.
    \end{align*}
\end{proof}

\begin{theorem}{5.17. (Serre)}
    Let $X$ be a projective scheme over a Noetherian ring $A$, let $\OO_X(1)$ be a very ample invertible sheaf on $X$, and let $\FF\in\coh(X)$. Then there is an $n_0\in\mathbb{Z}$ such that for every $n\geq n_0$, the sheaf $\FF(n)$ can be generated by finitely many global sections.
\end{theorem}
\begin{proof}\mbox{}
    \begin{itemize}
        \item By definition of the very ample sheaf $\OO_X(1)$, there is a closed immersion $i:X\xhookrightarrow{}\PP_A^r$ with $i^\ast\OO_{\PP_A^r}(1)=\OO_X(1)$. Then, we have $i_\ast\FF\in\coh(\PP_A^r)$ by \textit{[5.5c]}; and $i_\ast(\FF(n))=(i_\ast\FF)(n)$ by \textit{(5.12c)}; and that $\FF(n)$ is generated by its global sections $\Longleftrightarrow i_\ast(\FF(n))$ is also generated by its global sections (in fact, their global sections are the same as $(i_\ast\FF(n))(X)=\FF(n)(i^{-1}\PP_A^r)=\FF(n)(X)$. Thus, we may reduce to the case $X=\PP_A^r=\Proj S$ where $S\coloneqq A[x_0,\cdots,x_r]$ (now, the $\OO_X(1)$ is just the one defined before).
        \item Cover $X$ with open affine subsets $\{D_+(x_i)\}_{i=0}^r$. For each $1\leq i\leq r$, write $B_i\coloneqq A[\frac{x_0}{x_i},\cdots,\frac{x_r}{x_i}]$. Since $\FF\in\coh(X)$, we have $\FF|_{D_+(x_i)}\simeq\widetilde{M}_i$ for some $M_i\in\fg\MOD_{B_i}$, say $M_i\coloneqq B_i\spn\{s_{ij}\}_{j=1}^{n_i}$.
        \item By \textit{(5.14b)}, there is a large $n>0$ for which $x_i^ns_{ij}=t_{ij}|_{D_+(x_i)}$ for some $t_{ij}\in\Gamma(X,\FF(n))$. Indeed, to check the conditions, we note that
        \begin{itemize}
            \item $s_{ij}\in M_i=\Gamma(D_+(x_i),\FF)$ where $D_+(x_i)\simeq\Spec S_{(x_i)}=\Spec B_i$ covers $X$.
            \item $\FF\in\qcoh(X)$ since $\FF\in\coh(X)$. 
            \item $X_{x_i}=D_+(x_i)$ where $x_i\in S(1)=\Gamma(X,\OO_X(1))$.
            \item $\OO_X(1)=\OO_{\PP_A^r}(1)$ is invertible by \textit{(5.12a)}.
            \item $D_+(x_i)\cap D_+(x_j)\simeq\Spec S_{(x_ix_j)}$ is affine, and thus, quasi-compact.
        \end{itemize}
        \item Further, set $M_i'\coloneqq M_i(n)$ for each $i$, then
        \begin{align*}
            \FF(n)|_{D_+(x_i)}=(\FF\otimes\OO_X(n))|_{D_+(x_i)}\simeq\widetilde{M}_i\otimes_{\widetilde{B}_i}\widetilde{B}_i(n)\simeq(M_i(n))\widetilde{}=\widetilde{M}_i',
        \end{align*}
        i.e. $\FF(n)$ corresponds to $M_i'\in \MOD_{B_i}$ on $D_+(x_i)$. Then, the map
        \begin{align*}
            x_i^n:\FF\longrightarrow\FF(n)=\FF\otimes\OO_X(1)^{\oplus n},\quad t\longmapsto t\otimes x_i^n
        \end{align*}
        gives an isomorphism $\phi_i:M_i\xrightarrow[\MOD_{B_i}]{\sim}M_i(n)=M_i'$ (with inverse isomorphism induced by $x_i^{-n}$). Thus, these $x_i^ns_{ij}=\phi_i(s_{ij})$ generate $M_i'=\Gamma(D_+(x_i),\FF(n))$. Therefore, these $t_{ij}\in\Gamma(X,\FF(n))$ generate $\FF(n)$.
    \end{itemize}
\end{proof}

\begin{corollary}{5.18}
    Let $X$ be a projective scheme over a Noetherian ring $A$. Then any coherent sheaf $\FF\in\coh(X)$ can be written as a quotient of a sheaf $\EE$, where $\EE$ is a finite direct sum of twisted structure sheaves $\OO_X(n_i)$ for various integers $n_i$.
\end{corollary}
\begin{proof}
    By \textit{(5.17)}, let $\FF(n)$ be generated by finitely many global sections. Then, $\bigoplus_{i=1}^N\OO_X\to\FF(n)\to0$ is exact from the \textit{Note} above. As the twisting functor is an exact functor, $\bigoplus_{i=1}^N\OO_X(-n)\to\FF\to0$ is exact, i.e. $\psi:\bigoplus_{i=1}^N\OO_X(-n)\to\FF$ is surjective, and hence, $\FF\simeq\OO_X(n_i)^{\oplus N}/\ker\psi$.
\end{proof}

\begin{theorem}{5.19}
    Let $A\in\fg\ALG_k$, let $X$ be a projective scheme over $A$, and $\FF\in\coh(X)$. Then $\Gamma(X,\FF)$ is a finitely generated $A$-module. In particular, if $A=k$, then $\Gamma(X,\FF)$ is a finite-dimensional $k$-vector space.
\end{theorem}
\begin{proof}
    (I) Reduction of the proof.
    \begin{itemize}
        \item From \textit{(5.16b)}, write $X=\Proj S$ where $S$ is a graded ring finitely generated by $S_1$ over $S_0=A$.
        \item Set $M\coloneqq\Gamma_\ast(\FF)\in \GR\MOD_S$. By \textit{(5.15)}, $\widetilde{M}\simeq\FF$ (as $\FF\in\qcoh(X)$).
        \item Further, since $\FF\in\coh(X)$, it follows from \textit{(5.17)} that $\FF(n)$ is generated by finitely may global sections for $n\gg0$, say $s_i\in\Gamma(X,\FF(n))$ ($1\leq i\leq m$).
        \item Let $M'$ be the $S$-submodule of $M$ generated by $\{s_i\}_{i=1}^m$, i.e. $M'\in\fg\MOD_S$. Then, the exact sequence $0\to M'\xhookrightarrow{}M$ gives $0\to\widetilde{M}'\xhookrightarrow{}\widetilde{M}=\FF$, and thus, $0\to\widetilde{M}'(n)\xhookrightarrow{}\FF(n)$. But, $\FF(n)$ is generated by $\{s_i\}_{i=1}^n$, so $\widetilde{M}'(n)\simeq\FF(n)$. Apply the twisting functor again, we have $\widetilde{M}'\simeq\FF$, i.e. $\FF$ corresponds to a finitely generated graded $S$-module. 
    \end{itemize}
    (II) Now, it reduces to show that if $M\in\fg\MOD_S$, then $\Gamma(X,\widetilde{M})\in\fg\MOD_A$.
    \begin{itemize}
        \item Since $S_0=A\in\fg\ALG_k$, it is Noetherian. By \textit{(3.F)}, $S$ is Noetherian. Further, $M\in\GR\MOD_S$ is finitely generated, so it follows from \textit{(I,7.4a)} that $\exists$ filtration $0=M^0\subset M^1\subset\cdots\subset M^r=M$ in $\GR\MOD_S$ such that each $M^i/M^{i-1}\simeq(S/\pp_i)(n_i)$ for some $n_i\in\mathbb{Z}$ and homogeneous prime ideal $\pp_i\triangleleft S$.
        \item Such filtration induces a filtration $0=\widetilde{M}^0\subset\widetilde{M}^1\subset\cdots\subset\widetilde{M}^r=\widetilde{M}$ in $\MOD_{\OO_X}$. Note that taking the left exact functor $\Gamma(X,-)$ preserves injectivity, so we obtain a filtration 
        \begin{align*}
            0=\Gamma(X,\widetilde{M}^0)\subset\Gamma(X,\widetilde{M}^1)\subset\cdots\subset\Gamma(X,\widetilde{M}^r)=\Gamma(X,\widetilde{M})
        \end{align*}
        in $\MOD_{S_0}=\MOD_A$.
        \item Moreover, for each $i$, we have the exact sequence
        \begin{align*}
            0\to\Gamma(X,\widetilde{M}^{i-1})\to\Gamma(X,\widetilde{M}^i)\to\Gamma(X,(S/\pp_i)\widetilde{}(n_i)).
        \end{align*}
        So $\Gamma(X,\widetilde{M}^i)/\Gamma(X,\widetilde{M}^{i-1})$ is isomorphic to an $A$-submodule $N^i$ of $\Gamma(X,(S/\pp_i)\widetilde{}(n_i))$, and yields the short exact sequence of $A$-modules:
        \begin{align*}
            0\to\Gamma(X,\widetilde{M}^{i-1})\to\Gamma(X,\widetilde{M}^i)\to N^i\to0.
        \end{align*}
        \item Now, if we show that each $\Gamma(X,(S/\pp_i)\widetilde{}(n_i))\in \MOD_A$ is finitely generated, then its submodule $N^i$ over the the Noetherian ring $A$ is also finitely generated. Inductively, $\Gamma(X,\widetilde{M})\in \MOD_A$ is also finitely generated.
    \end{itemize}
    (III) It remains to show the special case that if $S$ is a graded integral domain, which is finitely generated by $S_1$ over $S_0=A$, where $A\in\fg\ALG_k$ is also an integral domain. Then $\Gamma(X,\OO_X(n))\in\fg\MOD_A$ for every $n\in\mathbb{Z}$.
    \begin{itemize}
        \item Note that the assumption of "integrality" of $S$ holds for the integral domain $S/\pp_i$ above. Further, the finitely generated $k$-algebra $(S/\pp_i)_0=A/(\pp_i\cap A)$ is also an integral domain.
        \item First, suppose $S_1\coloneqq A\spn\{x_j:0\leq j\leq r\}$ where $x_j\in S_1$. For every $n\in\mathbb{Z}$, since $S$ is an integral domain, the map $S(n)\xrightarrow[\Ring]{x_0}S(n+1),\ s\mapsto x_0s$ is clearly injective. (If $x_0s=0$, since $x_0\in S_1,\ x_0\neq0$, so $s=0$). Applying the left exact functor $\Gamma(X,-)$, we have $\Gamma(X,\OO_X(n))\xhookrightarrow{}\Gamma(X,\OO_X(n))$. Thus, it suffices to show that $\Gamma(X,\OO_X(n))\in\fg\MOD_A$ for all large $n\gg0$.
        \item Put $S'\coloneqq\bigoplus_{n\geq0}\Gamma(X,\OO_X(n))$. Then, $S'$ is a graded ring satisfying $S\subset S'\subset\bigcap_{i=0}^rS_{x_i}$. Indeed, we see that
        \begin{itemize}
            \item Clearly, from definition of $\OO_X$, every element of $S_0$ is in fact a global section of $\OO_X$, i.e. $S_0\subset\Gamma(X,\OO_X)$. In particular, $S_n\subset\Gamma(X,\OO_X(n))$ for every $n\geq0$, whence $S\subset S'$.
            \item Note that each $D_+(x_i)$ is affine, so 
            \begin{align*}
                \Gamma(X,\OO_X(n))\xhookrightarrow{}\Gamma(D_+(x_i),\OO_X(n))=\Gamma(D_+(x_i),\OO_X(n)|_{D_+(x_i)})=(S_{x_i})_n\subset S_{x_i},
            \end{align*}
            whence $S'\subset\bigoplus_{n\geq0}\bigcap_{i=0}^r(S_{x_i})_n=\bigcap_{i=0}^rS_{x_i}$.
        \end{itemize}
        \item \textbf{Claim.} $S'$ is integral over $S$.
        \begin{itemize}
            \item Let $\overline{S}$ denote the integral closure of $S$ in $K=\Frac S$, i.e. we show that $S'\subset\overline{S}$. But, $S'=\bigoplus_{d\geq0}S'_d$, so it suffices to show that for any $d\geq0$ and $s'\in S'_d,\ s$ is integral over $S$.
            \item Since $s'\in S'=\bigcap_{i=0}^rS_{x_i}$, for each $i,\ \exists\ n_i>0$ with $x_i^{n_i}s'\in S$. As usual, take $n\geq\{n_i:0\leq i\leq r\}$ such that $x_i^ns'\in S$ for all $i$.
            \item Since $S_1$ is generated by $\{x_i\}_{i=0}^r$ over $A$, for any $n\geq0$, we have $S_m=A\spn\{x_0^{d_0}\cdots x_r^{d_r}:\sum_{i=0}^rd_i=m\}$. Take $n\gg0$, then we may assume that $ys'\in S$ for every $y\in S_n$. (In fact, note that $\deg s'=d\geq0$, we have $ys'\in S_{\geq n}$ for any $y\in S_{\geq n}$).
            \item Inductively, for any $q\geq1$, we have $y\cdot(s')^q\in S_{\geq n}$ for any $y\in S_{\geq n}$. So, $x_i^n(s')^q\in S_{\geq n}$.
            \item Now, take $y=x_0^n$ for example. Then, we have $(s')^q\in\frac{1}{x_0^n}S$ for every $q\geq1$, whence $S[s']\subset\frac{1}{x_0^n}S$. Now, $\frac{1}{x_0^n}S$ is a finitely generated $S$-submodule of $\Frac S'$ such that $S[s']\subset\frac{1}{x_0^n}S\subset\Frac S'$. Therefore, by \textit{(AM, 5.1)}, $s'\in S'\subset\Frac S'$ is integral over $S$.  
        \end{itemize}
        \item Now, note that $S\in\fg\ALG_k$ is an integral domain. Taking $L=K=\Frac S$ in \textit{(I, 3.9)}, we have $\overline{S}\in\fg\MOD_S$. But $S'\subset\overline{S}$ from the claim, so $S'\in\fg\MOD_S$, whence $\Gamma(X,\OO_X(n))=S'_n\in\fg\MOD_A$ for every $n\geq0$.
    \end{itemize}
\end{proof}

\begin{corollary}{5.20}
    Let $f:X\xrightarrow[\Sch]{}Y$ be a projective morphism of schemes of finite type over $k$. If $\FF\in\coh(X)$, then so is $f_\ast\FF\in\coh(Y)$.
\end{corollary}
\begin{proof}\mbox{}
    \begin{itemize}
        \item Similar to \textit{(5.8c)}, this question is local on $Y$, so we may assume $Y=\Spec A$ is affine. Further, $Y$ is of finite type over $k$, so we have $A\in\fg\ALG_k$.
        \item From \textit{(5.16b)}, write $X=\Proj S$, where $S$ is a graded ring, which is finitely generated by $S_1$ over $S_0=A$. Now, $S_0=A\in\fg\ALG_k$ is Noetherian, so is $S$ by \textit{(3.F)}, whence $X=\Proj S$ is a Noetherian scheme by \textit{(3.H)}. By \textit{(5.8c)}, we have $f_\ast\FF\in\qcoh(Y)$, and thus,
        \begin{align*}
            f_\ast\FF\underset{\text{(5.5)}}{=}\Gamma(Y,f_\ast\FF)\widetilde{}\ \underset{\text{(Def. of $f_\ast$)}}{=}\Gamma(X,\FF)\widetilde{}.
        \end{align*}
        \item Finally, since $\FF\in\coh(X)$, by \textit{(5.19)}, we have $\Gamma(X,\FF)\in\fg\MOD_A$. Therefore, $f_\ast\FF\in\coh(Y)$.
    \end{itemize}
\end{proof}

\section*{Exercises}
\begin{exe}{5.3}
    Let $X=\Spec A$ be an affine scheme. Show that the functors $\widetilde{}$ and $\Gamma(X,-)$ are adjoint, i.e. for any $M\in \MOD_A$ and any $\FF\in\MOD_{\OO_X}$, there is a natural isomorphism
    \begin{align*}
        \Hom_{\MOD_A}(M,\Gamma(X,\FF))\simeq\Hom_{\MOD_{\OO_X}}(\widetilde{M},\FF).
    \end{align*}
\end{exe}

\begin{exe}{5.3*}
    Let $X=\Proj S$ where $S$ is a graded ring which is finitely generated by $S_1$ as an $S_0$-algebra. Show that the functors $\widetilde{}$ and $\Gamma_\ast(-)$ are adjoint, i.e. for any $M\in \GR\MOD_S$ and any $\FF\in\MOD_{\OO_X}$, there is a natural isomorphism
    \begin{align*}
        \Hom_{\MOD_S}(M,\Gamma_\ast(\FF))\simeq\Hom_{\MOD_{\OO_X}}(\widetilde{M},\FF).
    \end{align*}
\end{exe}

\chapter{Divisors}
\section{Weil Divisors}
\begin{axiom}{($\star$)}
    $X$ is a Noetherian, integral, separated scheme, which is regular in codimension $1$, (i.e. every local ring $\OO_x$ of dimension $1$ is regular).
\end{axiom}

\begin{definition}{(Weil Divisors)}
    Fix a scheme $X\in\Sch$ satisfying $(\star)$.
    \begin{itemize}
        \item \textbf{(Prime Divisor).} A \textit{prime divisor} on $X$ is a closed integral subscheme $Y$ such that $\codim(Y,X)=1$.
        \item \textbf{(Weil Divisor).} Let
        \begin{align*}
            \Div(X)\coloneqq\text{the free abelian group generated by the prime divisors}.
        \end{align*}
        Its elements are called \textit{Weil divisors}.
        \item \textbf{(Effective Divisors).} An \textit{effective divisor} is a (Weil) divisor $D=\sum_in_iY_i\in\Div(X)$ with $n_i\geq0$ for every $i$.
    \end{itemize}
\end{definition}

\begin{remark*}{}
    Let $X\in\Sch$ satisfy $(\star)$, and $Y$ be a prime divisor.
    \begin{itemize}
        \item As $Y$ is integral, it has a unique generic point $\eta\in Y$. Consider the local ring $\OO_{X,\eta}$. Let $U=\Spec A$ be an open affine neighborhood of $\eta\in X$, say $\eta=\pp\in\Spec A$. By \textit{Axiom}, $A$ is a Noetherian integral domain, so is its localization $A_\pp=\OO_{U,\eta}=\OO_{X,\eta}$. Further, note that 
        \begin{align*}
            \dim\OO_{X,\eta}=\codim(\{\eta\}\bar{},X)=\codim(Y,X)=1,
        \end{align*}
        again from the \textit{Axiom}, $\OO_{X,\eta}$ is regular. Therefore, the local ring $\OO_{X,\eta}$ is a DVR.
        \item Note that $X$ is also an integral scheme, it has a unique generic point $\mu\in X$ (caution: $\mu$ may not be the same as $\eta$). By \textit{(4.A)}, we have $\Frac\OO_{X,\eta}=k(\mu)=:K$, the function field of $X$.
        \item As $\OO_{X,\eta}$ is a DVR in $K,\ \exists$ discrete valuation $\nu_Y:K^\times\to\mathbb{Z}$ such that 
        \begin{gather*}
            \OO_{X,\eta}=\{f\in K:\nu_Y(f)\geq0\}\cup\{0\},\\
            \mm_{X,\eta}=\{f\in K:\nu_Y(f)>0\},\\
            \OO_{X,\eta}^\times=\{f\in K:\nu_Y(f)=0\},
        \end{gather*}
        where $\mm_{X,\eta}\coloneqq\<\pi\>$ is principal with $\nu_Y(\pi)=1$.
        \item Since $X$ is separated, by \textit{[4.5]}, such $Y$ is uniquely determined by $\nu_Y$. \textbf{(Why?)}
        \item Fix $f\in K^\times$, we say that $f$ has a \textit{zero [resp. pole]} along $Y$ of order $|\nu_Y(f)|\in\mathbb{Z}_{>0}$ if $\nu_Y(f)>0$ [resp. $\nu_Y(f)<0$].
    \end{itemize}
\end{remark*}

\begin{remark*}{}
    If $X=\Spec A$ is affine, there is a 1-1 correspondence
    \begin{align*}
        \{Y\subset X:\text{prime divisor}\} &\longleftrightarrow\ \{\pp\in\Spec A:\Ht\pp=1\},\\
        \Spec(A/\pp)\quad &\longmapsfrom\qquad\pp.
    \end{align*}
\end{remark*}
\begin{proof}
    Indeed, we see that
    \begin{itemize}
        \item $Y\subset X$ is a closed subscheme $\Leftrightarrow\ Y\simeq\Spec(A/\mathfrak{a})$ for some $\mathfrak{a}\triangleleft A$.
        \item $Y$ is an integral scheme $\Leftrightarrow\ A/\mathfrak{a}$ is an integral domain $\Leftrightarrow\ \mathfrak{a}=\pp\triangleleft A$ is a prime ideal.
        \item Further, $\codim(Y,X)=\dim A-\dim(A/\pp)=\Ht\pp$.
    \end{itemize}
\end{proof}

\begin{definition}{(Regular Functions)}
    Given an integral scheme $X$. We say
    \begin{itemize}
        \item $f\in K(X)$ is \textit{regular} at $x\in X$ if $f\in\OO_{X,x}\underset{\text{(4.A)}}{\xhookrightarrow{}}K(X)$.
        \item $f\in K(X)$ is \textit{regular} on an open affine subset $U=\Spec A$ of $X$ if $f\in\OO_X(U)=A\underset{\text{[3.6]}}{\xhookrightarrow{}}K(X)$.
    \end{itemize}
\end{definition}

\begin{lemma}{6.1}
    Let $X\in\Sch$ satisfy $(\star)$, and let $f\in K^\times$ be a nonzero function on $X$. Then the valuation $\nu_Y(f)=0$ for all but finitely many prime divisors $Y$ on $X$.
\end{lemma}
\begin{proof}\mbox{}
    \begin{itemize}
        \item Let $U=\Spec A$ be an open affine subset $X$ on which $f$ is regular, i.e. $f\in\OO_X(U)$. Then $Z\coloneqq X\setminus U$ is a proper closed subset of $X$.
        \item Note that if $\codim(Z,X)>1$, then $Z$ cannot contain a subscheme $Y$ with $\codim(Y,X)=1$, i.e. $Z$ cannot contain a prime divisor of $X$; now if $\codim(Z,X)=1$ and if $Y$ a prime divisor of $X$ contained in $Z$, then $Y$ is an irreducible component of $Z$ (as it is an irreducible subset of $X$). Since $X$ is a Noetherian scheme, it is a Noetherian topological space, so is its subspace $Z$ with induced topology. By \textit{(I, 1.5)}, $Z$ has finitely many irreducible components, i.e. there are only finitely many such prime divisors $Y\subset Z$ of $X$, all the others must meet $U$. 
        \item Now, to see the prime divisors which meet $U$, it suffices to consider the prime divisors $Y$ of $U$ as $U\subset X$ is open. It remains to show that there are only finitely many prime divisors $Y$ of $U=\Spec A$ for which $\nu_Y(f)\neq0$.
        \item Recall from the proof of \textit{[2.9]}, every open subset $V\subset Y$ must contain its generic point $\eta$. Then for the open subset $U\subset X$, the inclusion $\OO_X(U)\xhookrightarrow{}K$ can be factored as $\OO_X(U)\to\OO_{X,\eta}\xhookrightarrow{}K$, which gives $\OO_X(U)\xhookrightarrow{}\OO_{X,\eta}$. Now, we may view $f\in\OO_{X,\eta}$, so $\nu_Y(f)\geq0$.
        \item Note that if $Y$ is a prime divisor of $U=\Spec A$, it corresponds to some prime ideal $\pp\triangleleft A$, i.e. $Y=V(\pp)\simeq\Spec(A/\pp)$ from the \textit{Remark}. Then,
        \begin{align*}
            \nu_Y(f)>0\ \Leftrightarrow\ f\in\mm_{X,\eta}=\pp A_\pp\ \Leftrightarrow\ f\in\pp\ \Leftrightarrow\ \sqrt{\<f\>}\subset\sqrt{\pp}=\pp\ \Leftrightarrow\ Y=V(\pp)\subset V(\<f\>).
        \end{align*}
        \item Note that $X$ is integral, so $A=\OO_X(U)$ is a reduced integral domain, i.e. $\<0\>=\nil(A)\in\Spec A$. Now, since $f\in K^\times$ is nonzero, $f\notin\nil(A)$, so $W\coloneqq V(\<f\>)\subsetneq U$ is a proper closed subset (by \textit{[AM, 1.17ii]}). Moreover, $X$ is Noetherian, so is its open affine subset $U$, whence $U$ is a Noetherian topological space. Again by \textit{(I, 1.5)}, $W$ can only contain finitely many irreducible closed subset $Y\subset U$ with $\codim(Y,X)=1$, i.e. there can only be finitely many prime divisor $Y$ for which $\nu_Y(f)>0$.
    \end{itemize}
\end{proof}

\begin{proposition}{6.2}
    Let $A$ be a Noetherian domain. Then 
    \begin{align*}
        A\ \text{is a UFD}\ \Longleftrightarrow\ X=\Spec A\ \text{is normal and}\ \Cl X=0.
    \end{align*}
\end{proposition}
\begin{proof}\mbox{}
    \begin{itemize}
        \item Note that if $A$ is a UFD, then it is integrally closed, whence $X=\Spec A$ is normal.
        \item From \textit{(I.1.12)}, it suffices to show that if $A$ is an integrally closed Noetherian domain. Then,  every prime ideal of height $1$ is principal $\Longleftrightarrow\ \Cl(\Spec A)=0$.
        
        ($\Rightarrow$) Suppose every prime ideal of height $1$ is principal.
        \begin{itemize}
            \item Consider a prime divisor $Y$ of $X$, and let $\pp\in\Spec A$ be its corresponding prime ideal of height $1$. By assumption, say $\pp=\<f\>$ for some $f\in A$. Further, $\pp=\eta\in Y$ is its unique generic point, so $\OO_{X,\eta}=A_\pp$ is a DVR with unique maximal ideal $\mm_\eta=\pp A_\pp=\<\frac{f}{1}\>$. It follows that $\nu_Y(f)=\nu_Y(f)-\nu_Y(1)=\nu_Y(\frac{f}{1})=1$ since 
            \begin{align*}
                \OO_{X,\eta}=\{x\in K^\times:\nu_Y(x)>0\}\amalg\{x\in K^\times:\nu_Y(x)=0\}=\mm_\eta\amalg\OO_{X,\eta}^\times.
            \end{align*}
            \item If $Z\subset X$ is another prime divisor with corresponding prime ideal $\qq=\mu$. Since $f\in A$ and $A$ is integral, we have $\frac{f}{1}\in A_\qq=\OO_{Z,\mu}$, so $\nu_Z(f)=\nu_Z(\frac{f}{1})\geq0$. If $\nu_Z(\frac{f}{1})>0$, then $\frac{f}{1}\in\mm_\mu=\qq A_\qq$, i.e. $f\in\qq$. This implies that $\pp=\<f\>\subset\qq$. But, $\Ht\pp=\Ht\qq=1$, so $\pp=\qq$, whence $Z=\{\mu\}\bar{}=\{\eta\}\bar{}=Y$, a contradiction.
            Thus, we have $\nu_Z(f)=0$ whenever $Z\neq Y$, i.e. $\dv(f)=1\cdot Y$.
            \item Finally, for any divisor $D\in\Div X$, we have
            \begin{align*}
                D\coloneqq\sum_{\text{finite}}n_iY_i=\sum_{\text{finite}}n_i\dv(f_i)=\dv(\prod_{\text{finite}}f_i^{n_i}),
            \end{align*}
            which is a principal divisor. Therefore, $\Cl(\Spec A)=0$.
            
        \end{itemize}
        ($\Leftarrow$) Conversely, suppose $\Cl(\Spec A)=0$ and let $\pp\in\Spec A$ with $\Ht\pp=1$.
        \begin{itemize}
            \item Let $Y=\Spec(A/\pp)$ be the corresponding prime divisor of $X$ to $\pp$. Since $A$ is an integrally closed Noetherian domain, so is $A_\pp$. Further, $\dim A_\pp=\Ht\pp=1$, so $A_\pp$ is a DVR.
            \item For the divisor $Y\in\Div X$, since $\Cl(X)=0,\ \exists\ f\in K^\times$ such that $Y=\dv(f)$. So $\nu_Y(f)=1>0$, and hence, $f\in\pp A_\pp$ (the maximal ideal of $A_\pp$). But $A_\pp$ is a DVR, so its maximal ideal is principal, say $\pp A_\pp=\<\pi\>$. Then, $\nu_Y(\pi)=1$, so $\nu_Y(f/\pi)=0$, i.e. $f/\pi\in A_\pp^\times$. Thus, $\pp A_\pp=\<\pi\>=\<f\>$.
            \item If $\qq\in\Spec A$ is another prime ideal of height $1$, we have another corresponding prime divisor $Z=\Spec(A/\qq)$. As $\dv(f)=Y$, we see that $\nu_Z(f)=0$, i.e. $f\in A_\qq$. It follows that $f\in\bigcap_{\Ht\qq=1}A_\qq\underset{(6.3)}{=}A$. Now, $f\in A\cap\pp A_\pp=\pp$, which implies that $\<f\>\subset\pp$.
            \item For $g\in\pp$, we have $\frac{g}{1}\in\pp A_\pp$, so $\nu_Y(g)=\nu_Y(\frac{g}{1})>0$. By the similar argument as above, we see that if $W\subset X$ is any other prime divisor, then $\nu_W(g)=0$. Thus, we have $\nu_W(g/f)\geq0$, i.e. $g/f\in A_\qq$, for \textit{every} prime divisor $W\subset X$ with corresponding prime ideal $\qq$. Again, we have $g/f\in\bigcap_{\Ht\qq=1}A_\qq\underset{(6.3)}{=}A$, so $g\in\<f\>$. Thus, $\pp\subset\<f\>$, and therefore, $\pp=\<f\>$ is principal.
        \end{itemize}
    \end{itemize}
\end{proof}

\begin{proposition}{6.3}
    Let $A$ be an integrally closed Noetherian domain. Then 
    \begin{align*}
        A=\bigcap\{A_\pp:\pp\triangleleft A\ \text{prime with}\ \Ht\pp=1\}.
    \end{align*}
\end{proposition}

\begin{fact}{6.A}
    Let $X$ be a scheme satisfying $(\star)$, and let $x\in X$. Then
    \begin{align*}
        x\ \text{is not in any prime divisor $Y$ of}\ X\ \Longleftrightarrow\ x=\eta,\ \text{the unique generic point of}\ X.
    \end{align*}
\end{fact}
\begin{proof}
    ($\Leftarrow$) Suppose $\eta\in Y$ for some $Y$. Let $U=\Spec A$ be an open affine neighborhood of $\eta\in X$. Then $U\cap Y$ is a prime divisor of $U$, i.e. $U\cap Y\simeq\Spec(A/\pp)$ for some $\pp\in\Spec A$ with $\Ht\pp=1$. But, $\xi=(\overline{0})\in U\cap Y$, which corresponds to $\pp\in U$, is the generic point of $Y$, so $\eta=\xi$, i.e. $\pp=(0)$, which contradicts to $\Ht\pp=1$.\\
    
    ($\Rightarrow$) Let $U=\Spec A$ be an open affine neighborhood of $x\in X$, say $x=\pp$ in $U$. From the assumption, we have $\pp/\qq\notin\Spec(A/\qq)$ for any $\qq\in\Spec A$ with $\Ht\qq=1$. Now, if $\pp\neq(0)$, say $0\neq a\in\pp$, then $\exists\ \qq\triangleleft A$ minimal prime ideal such that $\<a\>\subset\qq\subset\pp$. By \textit{(I.1.11)}, we have $\Ht\qq=1$, whence $\pp/\qq\in\Spec(A/\qq)$, a contradiction. Therefore, $\pp=(0)$ in $U$, i.e. $x=\eta$.
\end{proof}

\begin{lemma}{6.B}
    Suppose $X=\PP_k^n=\Proj S$ for any field $k$ (i.e. $S_0=k$ and $S_+=\<x_0,\cdots,x_n\>$ is a maximal ideal of $S$). Then, $X$ satisfies $(\star)$, so we may talk about the divisors on $X=\PP_k^n$.
    \begin{itemize}
        \item[(i)] For any subset $Y\subset X$, we define $I_+(Y)\coloneqq\bigcap\{\pp:\pp\in Y\}$ (in particular, $I_+(\emptyset)=S$). Then, we have $I_+(Y_2)\subset I_+(Y_1)$ for any subsets $Y_1\subset Y_2\subset X$.
        \item[(ii)] If $\mathfrak{a}\triangleleft S$ is any homogeneous ideal, then $V_+(\mathfrak{a})=\emptyset\ \Longleftrightarrow\ \sqrt{\mathfrak{a}}=S$ or $S_+$.
        \item[(iii)] If $\mathfrak{a}\triangleleft S$ is any homogeneous ideal with $V_+(\mathfrak{a})\neq\emptyset$, then $I_+(V_+(\mathfrak{a}))=\sqrt{\mathfrak{a}}$.
        \item[(iv)] For any $Y\subset X$, we have $V_+(I_+(Y))=\overline{Y}$.
        \item[(v)] The operations $I_+(-)$ and $V_+(-)$ set up an inclusion reversing 1-1 correspondence:
        \begin{gather*}
            \{\mathfrak{a}\triangleleft S:\text{homogeneous radical ideals with}\ \mathfrak{a}\not\supset S_+\}\ \longleftrightarrow\ \{Y\subset X:\text{closed subset}\}\\
            \mathfrak{a}\longmapsto V_+(\mathfrak{a}),\quad I_+(Y)\longmapsfrom Y.
        \end{gather*}
        (Note that $S\leftrightarrow\emptyset$ and $S_+\mapsto\emptyset$). In particular, the restricts to the 1-1 correspondence between the following two subsets:
        \begin{align*}
            \Proj S=\{\pp\triangleleft S:\text{homogeneous prime ideals with}\ \pp\not\supset S_+\}\ \longleftrightarrow\ \{Y\subset X:\text{irreducible closed subset}\}.
        \end{align*}
        \item[(vi)] Further, the correspondence above yields:
        \begin{align*}
            \{\pp\in\Proj S:\Ht\pp=1\} &\longleftrightarrow\ \{Y\subset\PP_k^n:\text{prime divisor (i.e. irreducible hypersurface)}\}\\
            &\longleftrightarrow\ \{f\in S:\text{homogeneous and irreducible (up to associates)}\},\\
            \<f\> &\longmapsfrom\  V_+(f)\longmapsfrom\ f.
        \end{align*}
        if we give each closed subset the reduced induced subscheme structure.
    \end{itemize}
\end{lemma}
\begin{proof}
    (vi) Clearly, a homogeneous prime ideal with $\Ht\pp=1$ corresponds to a irreducible closed subset with $\codim(Y,X)=1$, and note that $\Ht S_+=\Ht\<x_0,\cdots,x_n\>=n+1>1$. This gives the first correspondence. For the second:
    \begin{itemize}
        \item Given such $\pp$ with $\Ht\pp=1$. As $S=k[x_0,\cdots,x_n]$ is a UDF, it follows from \textit{(I.1.12)} that $\pp=\<f\>$ is principal for some homogeneous $f\in S$. Further, $S$ is an integral domain, so $f$ is irreducible.
        \item Given an irreducible homogeneous polynomial $f\in S$. Since $S$ is a UFD, $\<f\>\triangleleft S$ is a homogeneous prime ideal and obviously $\<f\>\not\supset S_+$, whence $\<f\>\in\Proj S$. Moreover, as $S$ is Noetherian, \textit{(I.1.11)} implies that $\Ht\<f\>=1$.
    \end{itemize}
\end{proof}

\begin{proposition}{6.4}
    Let $X=\PP_k^n$ for any field $k$. For any divisor $D=\sum n_iY_i$, we define the degree of $D$ by $\deg D\coloneqq\sum n_i\deg Y_i$ where $\deg Y_i$ is the degree of the hypersurface $Y_i$, (i.e. if $Y_i=V_+(f_i)$, then $\deg Y_i=\deg f_i$). Let $H=V_+(x_0)$ be the hyperplane. Then:
    \begin{itemize}
        \item[(a)] If $D$ is any divisor of degree $d$, then $D\sim dH$.
        \item[(b)] For any $f\in K^\times,\ \deg\dv(f)=0$.
        \item[(c)] The degree function gives an isomorphism $\deg:\Cl(X)\to\mathbb{Z}$.
    \end{itemize}
\end{proposition}
\begin{proof}
    (b) Suppose $S=k[x_0,\cdots,x_n]$ is the homogeneous coordinate ring of $X=\PP_k^n$.
    \begin{itemize}
        \item Note that the function field $K$ of $X$ is $S_{((0))}=(\Frac S)_0$.
        \item If $g\in S_d$ ($d\geq0$) is a homogeneous polynomial. As $S$ is a UFD, we may factor $g$ into irreducible polynomials $g=ag_1^{n_1}\cdots g_r^{n_r}$ where $a\in k$ and $g_i\in S^h$ irreducible. By \textit{(6.B)}, each $Y_i\coloneqq V_+(g_i)$ is a hypersurface of degree $d_i=\deg g_i$. Now, we may extend the function $\dv:K^\times\to\Div(X)$ to $S^h\to\Div(X)$ by letting $\dv(g)\coloneqq\sum n_iY_i\in\Div(X)$. This makes sence since $\nu_{Y_i}(g)=\nu_{g_i}(g)=n_i$. It follows that 
        \begin{align*}
            \deg\dv(g)=\sum_{i=1}^rn_i\deg Y_i=\sum_{i=1}^rn_id_i=\deg g=d.
        \end{align*}
        \item Now, for any $f\in K^\times=(\Frac S)_0$, write $f\coloneqq\frac{g}{h}$ with $g,h\in S_d$ ($d\geq0$). Then, we have $\dv(f)=\dv(g)-\dv(h)$ (since $\nu_{g_i}(h)=0$ for every $i$ and similarly for $h$), whence
        \begin{align*}
            \deg\dv(f)=\deg\dv(g)-\deg\dv(h)=d-d=0.
        \end{align*}
    \end{itemize}
    
    (a) Suppose $D\in\Div(X)$ is any divisor of degree $d$. 
    \begin{itemize}
        \item We may write $D\coloneqq D_1-D_2$ such that both $D_1, D_2$ are effective divisors of degree $d_1,d_2\geq0$, respectively. 
        \item For any prime divisor $Y_i$, it corresponds to an irreducible homogeneous polynomial $h_i\in S$ from \textit{(6.B)}, i.e. $Y_i=V_+(h_i)$. Now, for $j=1,2$, set $g_j\coloneqq\prod_{\text{finite}}h_i^{n_{ij}}\in S^h$, then we have
        \begin{align*}
            D_j=\sum_{\text{finite}}n_{ij}Y_i=\sum_{\text{finite}}n_{ij}\dv(h_i)=\dv(g_j),
        \end{align*}
        where $\deg g_j=\sum_{\text{finite}}n_{ij}\deg h_i=\deg D_j=d_j$.
        \item Finally, take $f\coloneqq g_1/x_0^dg_2$, then
        \begin{align*}
            \deg f=\deg g_1-(d\deg x_0+\deg g_2)=d_1-d-d_2=0,
        \end{align*}
        whence $f\in(\Frac S)_0=K^\times$. Thus, 
        \begin{align*}
            D-dH=\dv(g_1)-\dv(g_2)-d\dv(x_0)=\dv(f)
        \end{align*}
        is indeed a principal divisor, and therefore, $D\sim dH$.
    \end{itemize}
    
    (c) From above, we obtain a group homomorphism $\deg:\Div(X)\to\mathbb{Z}$.
    \begin{itemize}
        \item By \textit{(b)}, the group of principal divisors is contained in $\ker(\deg)$.
        \item By \textit{(a)}, if $\deg D=0$, then $D\sim0$, i.e. $\ker(\deg)$ is contained in the group of principal divisors. Further, for any $d\in\mathbb{Z}$, we have $dH\mapsto d$. Therefore, this lifts to a group isomorphism $\deg:\Div(X)\isoto\mathbb{Z}$ by \textit{Isomorphism Theorem}.
    \end{itemize}
\end{proof}

\begin{proposition}{6.5}
    Let $X\in\Sch$ satisfy $(\star)$. Suppose $Z\subset X$ is a proper closed subset and $U\coloneqq X\setminus Z$. Then:
    \begin{itemize}
        \item[(a)] There is a surjective homomorphism 
        \begin{align*}
            \Cl(X)\longrightarrow\Cl(U),\quad D=\sum n_iY_i\longmapsto\sum_{Y_i\cap U\neq\emptyset}n_i(Y_i\cap U).
        \end{align*}
        \item[(b)] If $\codim(Z,X)\geq2$, then $\Cl(X)\to\Cl(U)$ is an isomorphism.
        \item[(c)] If $Z\subset X$ is an irreducible subset of codimension $1$, then there is an exact sequence 
        \begin{align*}
            \mathbb{Z}\longrightarrow\Cl(X)\longrightarrow\Cl(U)\longrightarrow0,
        \end{align*}
        where the first map is defined by $1\mapsto1\cdot Z$.
    \end{itemize}
\end{proposition}
\begin{proof}
    (a) We approach by the following steps:
    \begin{itemize}
        \item Let $Y\subset X$ be any prime divisor such that $Y\cap U\neq\emptyset$. I claim that $Y\cap U$ is a prime divisor of $U$. Note that $U$ also satisfies $(\star)$. Since $Y\cap U\xhookrightarrow{}U$ is the base change of the closed immersion $Y\xhookrightarrow{}X$ via $U\xhookrightarrow{}X$, it is also a closed immersion. $Y\cap U$ is also an integral scheme as it is an open subscheme of the integral scheme $Y$. Further, the codimension remains the same on its nonempty open subset, so $Y\cap U$ is indeed a prime divisor of $U$. 
        \item Now, we may define
        \begin{align*}
            \phi:\Div(X)\longrightarrow\Div(U),\quad Y\longmapsto
            \left\{\begin{array}{lc} 
            Y\cap U, &\text{if}\ Y\cap U\neq\emptyset \\ 
            0, &\text{otherwise}.\\ 
        \end{array}\right.
        \end{align*}
        Note that $Z\subsetneq X$ is a proper closed subset and $X\in\Top$ is Noetherian, there are only finitely many possibilities of $Y$ satisfying $Y\subset Z$, i.e. $Y\cap U=\emptyset$. 
        \item I claim that $\phi$ sends principal divisors of $X$ to those of $U$. Let $\xi,\eta$ be generic points of $X,Y$, respectively, where $Y\cap U\neq\emptyset$ and $K=\OO_{X,\xi}$. For any $f\in K^\times$, note that $\Frac\OO_{U,\eta}=\Frac\OO_{X,\eta}\underset{\text{(4.A)}}{=}K$, so $\nu_{Y\cap U}$ and $\nu_Y$ define the same valuation, i.e. $\nu_{Y\cap U}(f)=\nu_Y(f)$. Thus,
        \begin{align*}
            \phi:\dv(f)\longmapsto\sum_{Y_i\cap U\neq\emptyset}\nu_{Y_i}(f)(Y_i\cap U)=\sum_{Y_i\cap U\neq\emptyset}\nu_{Y_i\cap U}(f)(Y_i\cap U)=\dv_U(f).
        \end{align*}
        Therefore, $\phi$ induces a group homomorphism $\Cl(X)\to\Cl(U)$.
        \item Clearly, $\Cl(X)\to\Cl(U)$ is surjective since every prime divisor of $U$ is the restriction of its closure in $X$.
    \end{itemize}
    
    (b) Note that $X$ is a Noetherian topological space, we may decompose $Z$ into finitely many irreducible components $Z\coloneqq\bigcup_{i=1}^rZ_i$. If $Y\subset X$ is a prime divisor with $Y\cap U=\emptyset$, i.e. $Y\subset Z$, then $Y\subset Z_i$ for some $i$. It follows that
    \begin{align*}
        1=\codim(Y,X)\geq\codim(Z_i,X)\geq\codim(Z,X)\geq2,
    \end{align*}
    a contradiction. Hence,
    \begin{align*}
        \phi:\{Y\subset X:\text{prime divisor}\}\longrightarrow\{Y\cap U\subset U:\text{prime divisor}\}
    \end{align*}
    is bijective, and therefore, $\Cl(X)\isoto\Cl(U)$.\\
    
    (c) If there are no prime divisors $Y\subset X$ meeting $U$, then $Z$ is the only prime divisor of $X$. This implies that $\ker(\Cl X\to\Cl U)\subset\Im(\mathbb{Z}\to\Cl X)$. Note that the containment $(\supset)$ is clear, so the sequence is exact at $\Cl(X)$.
\end{proof}

\newpage
\section{Cartier Divisors}
\begin{definition}{(Sheaf of Total Quotient Rings)}
    \begin{itemize}
        \item Given $\Aa\in\Ring_X$. For each $U\in\TT_X$, let $\Ss(U)$ be a multiplicatively closed subset of $\Aa(U)$, then $U\mapsto\Ss(U)$ is called the \textit{multiplicatively closed subsheaf} of $\Aa$. We denote by $\Ss^{-1}\Aa\in\Ring_X$ the sheafification of $U\mapsto\Ss(U)^{-1}\Aa(U)$. Note that $(\Ss^{-1}\Aa)_x\simeq(\Ss_x)^{-1}\Aa_x$ for every $x\in X$.
        \item For $A\in\Ring$, we denote by $A_\reg$ the set of elements of $A$ which are not zero divisors, i.e.
        \begin{align*}
            A_\reg\coloneqq\{a\in A:sa\neq0\ \text{for any nonzero}\ s\in A\}.
        \end{align*}
        \item Given $\Aa\in\Ring_X$ and for any $U\in\TT_X$, we define
        \begin{align*}
            \Aa_\reg(U)\coloneqq\{s\in\Aa(U):s_x\in(\Aa_x)_\reg\ \text{for every}\ x\in U\},
        \end{align*}
        which is a multiplicatively closed subsheaf of $\Aa$. (Note: In general, $\Aa_\reg(U)\subsetneq\Aa(U)_\reg$). We define the \textit{sheaf of total quotient rings} of $\Aa$ by $\Frac\Aa\coloneqq(\Aa_\reg)^{-1}\Aa$.
        \item If $(X,\OO_X)\in\Sch$, we let $\KK_X\coloneqq\Frac\OO_X$. Note that there is a canonical monomorphism $\OO_X\xhookrightarrow{}\KK_X$.
        \item We denote by $\KK_X^\ast\in\Ab_X$ the sheaf $U\mapsto\KK_X(U)^\times$. Similarly, we can define $\OO_X^\ast$. Via the canonical monomorphism, we may regard $\OO_X^\ast$ as a subsheaf of $\KK_X^\ast$, and there is an exact sequence in $\Ab_X$:
        \begin{align*}
            1\longrightarrow\OO_X^\ast\longrightarrow\KK_X^\ast\longrightarrow\KK_X^\ast/\OO_X^\ast\longrightarrow1.
        \end{align*}
    \end{itemize}
\end{definition}

\begin{definition}{(Cartier Divisors)}
    \begin{itemize}
        \item A \textit{Cartier divisor} on $X\in\Sch$ is a global section $f\in\Gamma(X,\KK_X^\ast/\OO_X^\ast)$. Equivalently, we may represent $f$ as $\{(U_i,f_i)\}$ where $\{U_i\}$ is an open cover of $X$ and each $f_i\in\KK_X^\ast(U_i)$ such that $f_i/f_j\in\OO_X^\ast(U_i\cap U_j)$ for every $i,j$. We denote the group of all Cartier divisors on $X$ by $\CDiv(X)$.
        \item A Cartier divisor $f\in\CDiv(X)$ is \textit{principal} if $f$ is in the image of $\Gamma(X,\KK_X^\ast)\to\Gamma(X,\KK_X^\ast/\OO_X^\ast)$.
        \item Two Cartier divisors $f,g\in\CDiv(X)$ are \textit{linearly equivalent} if $f/g$ is principal.
    \end{itemize}
\end{definition}

\begin{proposition}{6.C}
    \begin{itemize}
        \item[(i)] If $X\in\Top$ is irreducible, then the constant presheaf $\FF:U\mapsto A$ (where $A\in\Ab,\Ring$ or $\MOD$) is already a sheaf.
        \item[(ii)] Suppose $X\in\Top$ is irreducible and has an open cover $\{U_i\}_{i\in I}$. If $\GG\in\Ab_X$ (or $\Ring_X,\MOD_X$) is such that each $\GG|_{U_i}$ is a constant sheaf, then so is $\GG$.
        \item[(iii)] If $X\in\Sch$ is an integral scheme, then $\KK_X\in\ALG_X$ is the constant sheaf of $K=K(X)$. 
    \end{itemize}
\end{proposition}
\begin{proof}[Proof of (i)]
    Note that if $X$ is irreducible, then $\bigcap_{U\in\TT_X}U=\{\eta\}$ where $\eta\in X$ is the unique generic point by \textit{[2.9]}. Thus, any open sets have nonempty intersection. Further, if $V\subset U$ are nonempty open sets, we have $\FF(U)=A=\FF(V)$, i.e. the restriction map is $\mathds{1}_A$. Now, given any $\emptyset\neq U\in\TT_X$ with an open cover $\{V_i\}$ of $U$:
    \begin{itemize}
        \item[(S3)] Suppose $s\in\FF(U)$ with $s|_i=0$ for every $i$. Clearly, $s=\mathds{1}_A(s)=s|_i=0$ (as $U\cap V_i\neq\emptyset$, i.e. $\FF(U\cap V_i)=A$).
        \item[(S4)] Suppose we have $s_i\in\FF(V_i)$ such that $s_i|_j=s_j|_i$ for every $i,j$. Since $U\cap V_i\cap V_j\neq\emptyset$, we have $s_i=\mathds{1}_A(s_i)=s_i|_j=s_j|_i=\mathds{1}_A(s_j)=s_j$ in $A$. So, $s\coloneqq s_i$ for every $i$ in $A=\FF(V_i)$.
    \end{itemize}
    Therefore, $\FF$ is indeed a sheaf.
\end{proof}
\begin{proof}[Proof of (ii)]
    For each $i$, say $\GG|_{U_i}\simeq\FF_i$ where $\FF_i$ denotes the constant sheaf of some $A_i$. For any $i,j$, since $U_i\cap U_j\neq\emptyset$, we have $A_i\simeq\FF_i(U_i\cap U_j)\simeq\GG(U_i\cap U_j)\simeq\FF_j(U_i\cap U_j)\simeq A_j$. Put $A\coloneqq A_i$ for every $i$, and let $\FF$ be the constant sheaf of $A$. I claim that $\GG\simeq\FF$.\\
    
    We define a morphism $\psi:\GG\to\FF$ as follows: for each $\emptyset\neq U\in\TT_X$, let
    \begin{align*}
        \psi_U:\GG(U)\to\GG(\underbrace{U\cap U_i}_{\neq\emptyset})\isolongto\FF(U\cap U_i)=A=\FF(U).
    \end{align*}
    For any $P\in X$, say $P\in U_i$, then $\psi_P:\GG_P=(\GG|_{U_i})_P\isoto A=\FF_P$ is an isomorphism. Thus, $\psi$ defines an isomorphism.
\end{proof}
\begin{proof}[Proof of (iii)]
    Choose an open affine cover $\{U_i\}$ of $X$, say each $U_i=\Spec A_i$ with $A_i$ an integral domain. Since $X$ is integral, $\OO_{X,x}$ is an integral domain for every $x\in X$, whence $\OO_\reg(U_i)=A_i\setminus\{0\}$. It follows that $\,'\KK(U_i)=\Frac A_i\underset{\text{[3.6]}}{\simeq}K$ (as rings). Now, for any $x\in U_i$, we have
    \begin{align*}
        K=\,'\KK(U_i)\longrightarrow\KK_x=\Frac\OO_{X,x}\underset{\text{(4.A)}}{\xhookrightarrow{~~~~~}}\Frac\OO_{X,\eta}=K,
    \end{align*}
    where the first map is also injective (since $K$ is a field). Hence, $\,'\KK|_{U_i}$ is the constant presheaf $K$. However, $U_i$ is irreducible (as $X$ is), by \textit{(i)}, we have $\KK|_{U_i}=\,'\KK|_{U_i}$. Finally, by \textit{(ii)}, $\KK$ is the constant sheaf of $K$.
\end{proof}

\begin{proposition}{6.11}
    Let $X$ be an integral, separated, Noetherian, \textit{locally factorial} (i.e. each local ring $\OO_{X,x}$ is a UFD) scheme. Then the groups $\Div(X)$ and $\CDiv(X)$ are isomorphic. Further, the principal Weil divisors correspond to the principal Cartier divisors under this isomorphism.
\end{proposition}
\begin{proof}
    Note that for each $x\in X$, the stalk $\OO_x$ is a UFD, and thus, integrally closed, whence $X$ is normal. So, $X$ satisfies $(\star)$, and we may define Weil divisors on $X$. \\
    
    (I) Given $\{(U_i,f_i)\}\in\CDiv(X)$, we construct a $D\in\Div(X)$.
    \begin{itemize}
        \item Note that each $f_i\in\KK^\ast(X)\underset{\text{(6.C)}}{=}K^\times$. For each prime divisor $Y$ with its generic point $\xi$, we take the coefficient to be $c_Y\coloneqq\nu_Y(f_i)$ for any $i$ with $U_i\cap Y\neq\emptyset$. 
        \item \textbf{(Well-defined).} If there is another $j$ with $U_j\cap Y\neq\emptyset$, we have $f_i/f_j\in\OO^\ast(U_i\cap U_j)=\OO(U_i\cap U_j)^\times\xhookrightarrow{}\OO_\xi^\times$, so $\nu_Y(f_i/f_j)=0$, i.e. $\nu_Y(f_j)=\nu_Y(f_i)=c_Y$.
        \item Thus, we define $D\coloneqq\sum c_YY$, which is indeed a finite sum (since $X$ is Noetherian, we may choose finitely many $U_i$, and apply \textit{(6.1)}).
    \end{itemize}
    (II) Given $D\coloneqq\sum n_iY_i\in\Div(X)$, we construct a Cartier divisor.
    \begin{itemize}
        \item Given any $x\in X$. If $x\in Y_i$, we have the following diagram.
        \begin{center}
            \begin{tikzcd}
                \Spec(A_\pp\otimes_AA/\qq_i)\FibSq{dr}\rar[]\dar["\text{(v)}"] &\Spec(A/\qq_i)\rar[equal,"\text{(iii)}"]\dar[] &U\cap Y_i\rar[equal,"\text{(3.3)}"] &U\times_XY_i\FibSq{dr}\rar[]\dar["\text{(ii)}"] &Y_i\dar[hook,"\text{closed}"]\\
                T_x=\Spec\OO_x\rar["\text{(iv)}"'] &\Spec A\arrow[rr,equal,"\text{(i)}"'] & &U\rar[hook,"\text{open}"'] &X
            \end{tikzcd}
        \end{center}
        Note that
        \begin{itemize}
            \item[(i)] Choose an open affine neighborhood $U=\Spec A$ of $x\in X$, then $A$ is an integral domain.
            \item[(ii)] $Y_i$ is a prime divisor with generic point $\xi_i$. By \textit{[3.11a]}, the map (ii) is a closed immersion. 
            \item[(iii)] Now, $U\cap Y_i\xhookrightarrow{}U$ is a closed subscheme of the integral scheme $\Spec A$, so such $\qq_i\triangleleft A$ determines $U\cap Y_i$. Further, $U\cap Y_i$ is an open subscheme of the integral scheme $Y_i$, so $A/\qq_i$ is an integral domain, i.e. $\qq_i\in\Spec A$.
            \item[(iv)] Let $x=\pp\in\Spec A=U$, then the map (iv) is just the map induced by $A\xhookrightarrow{}A_\pp$.
            \item[(v)] Again, as (ii) is a closed immersion, so is the base extension (v) via the map (iv). Further, $A_\pp\otimes_AA/\qq_i\simeq A_\pp/\qq_iA/\pp$ is an integral domain, whence $Y_i\cap T_x=\Spec(A_\pp\otimes_AA/\qq_i)$ is a closed integral subscheme of $T_x$.
        \end{itemize}
        Moreover, we see that
        \begin{align*}
            \codim(Y_i\cap T_x,T_x) &=\codim(\qq_iA_\pp,A_\pp)=\codim(A/\qq_i,A)=\codim(U\cap Y_i,U)\\&=\codim(Y_i,X)=1.
        \end{align*}
        Thus, $U\cap Y_i$ is a prime divisor of $U$ with generic point $\qq_i$. Further, $x\in U\cap Y_i$, so $Y_i\cap T_x$ is also a prime divisor of $T_x$. Therefore, $D$ restricts to a local divisor $D_x\coloneqq\sum_{Y_i\ni x}n_i(Y_i\cap T_x)\in\Div(T_x)$.
        \item Note that $X$ is a Noetherian domain, so is $T_x$, and thus, $\OO_x$. Further, $\OO_x$ is a UFD by assumption. It follows from \textit{(6.2)} that $\Cl(T_x)=0$. So, $\exists\ f_x\in K^\times$ such that $D_x=\dv(f_x)$ in $\Div(T_x)$.
        \item Now, consider $\dv(f_x)=\sum\nu_{Y_i}(f_x)Y_i$ in $\Div(X)$, which also restricts to 
        \begin{align*}
            \sum_{Y_i\ni x}\nu_{Y_i\cap T_x}(f_x)(Y_i\cap T_x)=\sum_{Y_i\ni x}\nu_{Y_i}(f_x)(Y_i\cap T_x)=D_x\quad\text{in}\ \Div(T_x)
        \end{align*}
        as the valuations $((A_\pp)_{\qq_i},\nu_{Y_i\cap T_x})$ and $(A_{\qq_i},\nu_{Y_i})$ are the same. Therefore, $\dv(f_x)$ and $D$ only differ on those prime divisors $Y_i\not\ni x$. Hence, there are only finitely many $Y_i\not\ni x$ for which $\nu_{Y_i}(f_x)\neq0$ or $n_i\neq0$.
        \item In fact, the space $T_x=\bigcap\{U_x:U_x\in\NN_x(X)\}$, so we may replace $T_x$ by some open neighborhood $U_x$ so that $D$ and $\dv(f_x)$ have the same restriction to $U_x$, (e.g. the open affine $U$ chosen above). Now, cover $X$ by these open sets $\{U_x\}_{x\in X}$. Indeed, each $f_x\in K^\times=\KK^\ast(U_x)$, we claim that $\{(U_x,f_x)\}_{x\in X}$ gives a well-defined Cartier divisor.
        \item \textbf{(Well-defined).} Suppose $f,f'\in K^\times$ give the same Weil divisor on some open set $U\subset X$, i.e. $\nu_{Y_i}(f)=\nu_{Y_i}(f')$ for every $Y_i$ with $Y_i\cap U\neq\emptyset$. Then, we have $\nu_{Y_i}(f/f')=0$, i.e. $f/f'\in\OO_{\xi_i}^\times$.
        \begin{itemize}
            \item For any open affine subset $V=\Spec B$ of $X$, say $\xi_i=\qq_i\in\Spec B$. Then $\OO_{\xi_i}^\times=B_{\qq_i}^\times$ and $\Ht\qq_i=1$. Since $B=\OO(V)$ is a normal Noetherian domain, we have
            \begin{align*}
                f/f'\in\bigcap_{Y_i\cap V\neq\emptyset}\OO_{\xi_i}=\bigcap_{\Ht\qq_i=1}B_{\qq_i}\underset{\text{(6.3)}}{=}B
            \end{align*}
            \item Now, let $\{V_i\}$ be an open affine cover of $U$, then we have $(f/f')|_i\in\OO(V_i)$ for each $i$. Further, for each $i,j$, their restriction to $\OO(V_i\cap V_j)$ agree, so we may glue $\{(f/f')|_i\}$ to obtain an $s\in\OO(U)$. In fact, $s=f/f'$ since each $(f/f')|_i\in\Frac\OO(V_i)=\KK(V_i)=K$. Thus, $f/f'\in\OO(U)$. Similarly, $f'/f\in\OO(U)$, so $f/f'\in\OO(U)^\times=\OO^\ast(U)$.
        \end{itemize}
        Now, if we replace $(f,f',U)$ by $(f_x,f_{x'},U_x\cap U_{x'})$, then we have $f_x/f_{x'}\in\OO^\ast(U_x\cap U_{x'})$, whence $\{(U_x,f_x)\}\in\CDiv(X)$. If there is another $\{(V_i,f'_i)\}$ representing this Cartier divisor, replacing $(f,f',U)$ by $(f_x,f'_i,U_x\cap V_i)$ shows that $f_x/f'_i\in\OO^\ast(U_x\cap V_i)$. Thus, such Cartier divisor is independent of the choice of the open cover.
    \end{itemize}
    (III) Finally, we show that the two constructions are inverse to each other.
    \begin{itemize}
        \item For any $D\in\Div(X)$,
        \begin{align*}
            D=\sum n_iY_i &\longmapsto\ \{(U_x,f_x)\}\ \text{with $\nu_{Y_i}(f_x)=n_i$ if $Y_i\ni x$}\\ &\longmapsto\ \sum c_{Y_i}Y_i\ \text{with $c_{Y_i}=\nu_{Y_i}(f_x)$ if $Y_i\cap U_x\neq\emptyset$}\\ &=\sum n_iY_i=D.
        \end{align*}
        \item For any $C\in\CDiv(X)$,
        \begin{align*}
            C=\{(U_i,f_i)\} &\longmapsto\ D=\sum c_Y\cdot Y\ \text{with $c_Y=\nu_Y(f_i)$ if $U_i\cap Y\neq\emptyset$}\\ &\longmapsto\ \{(U_x,f_x)\}\ \text{with $\nu_Y(f_x)=n_i=c_Y$ if $Y\ni x$}
        \end{align*}
        Note that if $Y\cap(U_i\cap U_x)\neq\emptyset$, then $\nu_Y(f_x)=c_Y=\nu_Y(f_i)$, whence $\nu_Y(f_x/f_i)=0$, i.e. $f_x/f_i\in\OO^\ast(U_i\cap U_x)$. Therefore, $\{(U_x,f_x)\}=\{(U_i,f_i)\}=C$.
    \end{itemize}
\end{proof}

\begin{remark}{6.11.1}
    \begin{itemize}
        \item \textbf{(Regular Schemes).} A scheme $X$ is \textit{regular} if each stalk $\OO_{X,x}$ is a regular local ring.
        \item Recall that a regular local ring is a UFD, whence \textit{(6.11)} also applies to any regular integral separated Noetherian scheme.
    \end{itemize}
\end{remark}

\begin{remark}{6.11.2}
    \begin{itemize}
        \item $D\in\Div(X)$ is a \textit{locally principal} Weil divisor if $X$ has an open cover $\{U_i\}$ such that each restriction $D|_{U_i}\in\Div(U_i)$ is principal.
        \item If $X$ is a normal scheme (not necessarily locally factorial), then the proof in \textit{(6.11)} shows that the Cartier divisors are the same as locally principal Weil divisors. Thus, in this case, the map $\CDiv(X)\to\Div(X)$ is injective.
    \end{itemize}
\end{remark}

\newpage
\section{Invertible Sheaves}
\begin{proposition}{6.12}
    If $\LL$ and $\MM$ are invertible sheaves on a ringed space $X$, so is $\LL\otimes_{\OO_X}\MM$. If $\LL$ is any invertible sheaf on $X$, then there exists an invertible sheaf $\LL^{-1}$ on $X$ such that $\LL\otimes_{\OO_X}\LL^{-1}\simeq\OO_X$.
\end{proposition}
\begin{proof}
    For the first statement, note that $\OO_X\otimes_{\OO_X}\OO_X\simeq\OO_X$. For the latter, we take $\LL^{-1}\coloneqq\LL^\vee$, then we have $\LL^\vee\otimes_{\OO_X}\LL\simeq\HHom_{\OO_X}(\LL,\LL)=\OO_X$ by \textit{[5.1]}.
\end{proof}

\begin{definition}{(Picard Group)}\newline
    For any ringed space $X$, we define the \textit{Picard group} on $X$, denoted by $\Pic(X)$, to be the group of isomorphism classes of invertible sheaves on $X$, under the operation $\otimes$. Note that \textit{(6.12)} shows that $(\Pic(X),\otimes)$ is a group with identity $\OO_X$.
\end{definition}

\begin{definition}{(Associated Sheaf to Cartier Divisors)}
    \newline Fix a scheme $X$ and let $D\in\CDiv(X)$ be represented by $\{(U_i,f_i)\}$. We define $\LL(D)$ to be the $\OO_X$-submodule of $\KK_X$ generated locally by $f_i^{-1}$ on $U_i$. (Indeed, this is well-defined, since $f_i/f_j\in\OO_X^\ast(U_i\cap U_j)$, both $f_i^{-1}$ and $f_j^{-1}$ generated the same $\OO_X$-submodule). In particular, we have $\LL(0)=\OO_X$.
\end{definition}

\begin{lemma}{6.D}
    Given $X\in\Sch$ and any $U\in\TT_X$. Then 
    \begin{align*}
        f\in\KK_X(U)\ \text{is a unit}\ \Longleftrightarrow\ f_x\in\KK_{X,x}\ \text{is $\OO_{X,x}$-torsion-free for every}\ x\in U.
    \end{align*}
\end{lemma}
\begin{proof}
    The case $U=\emptyset$ is clear, so we may assume $U\neq\emptyset$.
    \begin{itemize}
        \item Suppose $f\in\KK(U)$ such that $f(x)\in(\:'\KK_x)^\times$ for every $x\in U$. (Note $\KK_x=\:'\KK_x$). Take $g(x)\coloneqq(f(x))^{-1}$ for each $x\in U$, this defines $g\in\:'\KK(U)$. We check that $g\in\KK(U)$. Indeed, if $x\in U$, then $\exists\ V\in\NN_x(X)$ with $V\subset U$ and $\exists\ s\in\:'\KK(V)$ such that $f(y)=s_y$ in $\:'\KK_y$ for every $y\in V$. Now, if $t\in\:'\KK(W)$ with $t_y=(s_y)^{-1}$ in $\:'\KK_y$ for every $y\in V\cap W$. Then, $g(y)=t_y$ for every $y\in V\cap W$, which shows that $g\in\KK(U)$. Further, it is easy to see that $g\in\KK(U)^\times$. Therefore, $f\in\KK(U)^\times\ \Longleftrightarrow\ f(x)\in\KK_x^\times$ for every $x\in U$. Hence, it reduces to show the following claim.
        \item \textit{Claim}. $f(x)\in(\:'\KK_x)^\times$ for every $x\in U\ \Longleftrightarrow\ f(x)$ is $\OO_x$-torsion-free, i.e. $\Ann_{\OO_x}(f(x))=0$ for every $x\in U$.
        
        ($\Rightarrow$) Given $x\in U$, if $s\in\Ann_{\OO_x}(f(x))$, then $s\cdot f(x)=0$ in $\:'\KK_x$, thus $s=0$ (since $f(x)$ is a unit).
        ($\Leftarrow$) Given $x\in U$, then $\exists\ V\in\NN_x(X)$ with $V\subset U$ such that $f(y)=(\frac{a}{s})_y$ for every $y\in V$, (where $a\in\OO_X(V),s\in\OO_{X,\reg}(V)$). Since $f(y)$ is $\OO_y$-torsion-free for each $y\in V$, so is each $a_y\in\OO_y$. Thus, $a_y\in\OO_y$ is regular for every $y\in V$, i.e. $a\in\OO_{X,\reg}(V)$. Thus, $f(x)=(\frac{a}{s})_x\in\:'\KK_x$ is a unit.
    \end{itemize}
\end{proof}

\begin{corollary}{6.E}
    Let $X$ be a scheme and $\LL$ be an invertible $\OO_X$-submodule of $\KK_X$. Then, for each $x\in X,\ \LL_x$ is an $\OO_{X,x}$-submodule of $\KK_{X,x}$ generated by a unit.
\end{corollary}
\newpage
\begin{lemma}{6.F}
    \begin{itemize}
        \item[(i)] \textbf{(Algebraic Fact).} Given $A\in\Ring,\ B\in\ALG_A$. Suppose $M,N$ are $A$-submodules of $B$ such that one of them is generated by a regular element of $B$ over $A$. Then, the natural homomorphism $M\otimes_AN\underset{\ALG_A}{\twoheadrightarrow}M\cdot N,\ m\otimes n\mapsto mn$, is an isomorphism.
        \item[(ii)] For $X\in\Sch$ and $\OO_X$-submodules $\FF,\GG$ of $\KK_X$, we define $\FF\cdot\GG\in\ALG_{\OO_X}$ to be the sheafification of $U\mapsto\FF(U)\cdot\GG(U)$. Now, if both $\FF$ and $\GG$ are invertible sheaves, then the canonical morphism $\FF\otimes_{\OO_X}\GG\underset{\ALG_{\OO_X}}{\twoheadrightarrow}\FF\cdot\GG$ is an isomorphism.
    \end{itemize}
\end{lemma}
\begin{proof}
    (ii) follows from \textit{(6.E)} and \textit{(i)}.
\end{proof}


\begin{proposition}{6.13}
    Let $X$ be a scheme. Then:
    \begin{itemize}
        \item[(a)] For any $D\in\CDiv(X)$, the subsheaf $\LL(D)$ is an invertible sheaf on $X$. The map $D\mapsto\LL(D)$ gives a 1-1 correspondence between
        \begin{align*}
            \CDiv(X)\ \longleftrightarrow\ \{\text{invertible subsheaves of}\ \KK_X\}.
        \end{align*}
        \item[(b)] $\LL(D_1-D_2)\simeq\LL(D_1)\otimes\LL(D_2)^{-1}$.
        \item[(c)] $D_1\sim D_2\ \Longleftrightarrow\ \LL(D_1)\simeq\LL(D_2)$ as invertible $\OO_X$-modules.
    \end{itemize}
\end{proposition}
\begin{proof}
    (a) We approach by the following steps:
    \begin{itemize}
        \item[(I)] Given any $D=\{(U_i,f_i)\}\in\CDiv(X)$, show that $\LL(D)$ is invertible.
        
        For each $i$, consider the morphism $\phi_i:\OO_X|_{U_i}\to\LL(D)|_{U_i}$ given by
        \begin{align*}
            (\phi_i)_V:\OO_X(V)\longrightarrow\LL(D)(V),\quad1\longmapsto f_i^{-1}|_V\quad\text{for every}\ V\subset U_i.
        \end{align*}
        This induces the map $(\phi_i)_x:\OO_x\to\LL(D)_x,\ 1_x\mapsto(f_i^{-1})x$ for each $x\in U_i$. From the definition of $\LL(D)$, we have $\LL(D)_x=\OO_x\spn\{(f_i^{-1})_x\}$ and that $(f_i^{-1})_x,(f_j^{-1})_x$ generate the same $\OO_x$-submodule (as $(f_i^{-1})_x/(f_j^{-1})_x$ is invertible in $\LL(D)_x$). It follows that $(\phi_i)_x$ is an isomorphism for each $x\in U_i$, whence $\phi_i:\OO_X|_{U_i}\isoto\LL(D)|_{U_i}$ for every $i$. Therefore, $\LL(D)$ is an invertible $\OO_X$-module.
        \item[(II)] For any invertible $\OO_X$-submodule $\LL$ of $\KK$, construct a Cartier divisor.
        \begin{itemize}
            \item Cover $X$ with an open cover $\{U_i\}$ such that $\LL|_{U_i}\simeq\OO_X|_{U_i}$ for all $i$. For each $i$, pick an $f_i\in\LL(U_i)\subset\KK(U_i)$ corresponding to $1\in\OO(U_i)$ under this isomorphism.
            \item \textit{Check.} $f_i\in\KK^\ast(U_i)$ for each $i$. From \textit{(6.D)} it suffices to show that $(f_i)_x\in\LL_x\subset\KK_x$ is $\OO_x$-torsion-free for every $x\in U_i$. But this is clear since $\LL_x\simeq\OO_x$.
            \item \textit{Claim.} $\{(U_i,f_i)\}\in\CDiv(X)$. Indeed, for any $i,j$, since $\LL|_{U_i}\simeq\OO_X|_{U_i}$ and $\LL|_{U_j}\simeq\OO_X|_{U_j}$, both $f_i|_j$ and $f_j|_i$ generate $\LL(U_{ij})$ over $\OO_X(U_{ij})$, (here $U_{ij}\coloneqq U_i\cap U_j$). So, $\exists\ u\in\OO_X(U_{ij})$ such that $f_i|_j=u\cdot f_j|_i$, i.e. $(f_i^{-1|_j})/(f_j^{-1}|_i)\in\OO_X(U_{ij})^\times=\OO_X^\ast(U_{ij})$.
            \item \textit{Check.} $\{(U_i,f_i^{-1})\}$ is independent of the choice of the open cover $\{U_i\}$. Thus, $\{(U_i,f_i^{-1})\}\in\CDiv(X)$ is a well-defined Cartier divisor.
        \end{itemize}
        \item[(III)] Check the constructions are inverse to each other.
        \begin{itemize}
            \item For any Cartier divisor $D\in\CDiv(X)$,
            \begin{align*}
                D=\{(U_i,f_i)\} &\longmapsto\ \LL(D):\ \text{generated by $\{f_i^{-1}\}$ with}\ 1\in\OO_X|_{U_i}\xleftrightarrow{\sim}f_i^{-1}\in\LL|_{U_i}\\
                &\longmapsto\ \{(U_i,(f_i^{-1})^{-1})\}=D.
            \end{align*}
            \item For any invertible $\OO_X$-submodule $\LL\subset\KK$ with $f_i\in\LL|_{U_i}\xleftrightarrow{\sim}1\in\OO_X|_{U_i}$, we have
            \begin{align*}
                \LL\ \longmapsto\ \{(U_i,f_i^{-1})\}=D\ \longmapsto\ \LL(D):\ \text{generated by}\ \{(f_i^{-1})^{-1}\}=\{f_i\},
            \end{align*}
            whence $\LL(D)=\LL$.
        \end{itemize}
    \end{itemize}
    %\mbox{}
    
    (b) By choosing a common refinement of open covers, may assume $D_1=\{(U_i,f_i)\}$ and $D_2=\{(U_i,g_i)\}$ (where $f_i,g_i\in\KK^\ast(U_i)$). Then, we have $D_1+D_2=\{(U_i,f_ig_i)\}$. Now, for each $x\in X$, we have
    \begin{align*}
        \LL(D_1+D_2)_x=\OO_x\spn\{(f_ig_i)^{-1}_x\}=\OO_x\spn\{(f_i)^{-1}_x(g_i)^{-1}_x\}=(\LL(D_1)\LL(D_2))_x,
    \end{align*}
    whence $\LL(D_1+D_2)=\LL(D_1)\cdot\LL(D_2)\underset{\text{(6.F)}}{\simeq}\LL(D_1)\otimes_{\OO_X}\LL(D_2)$.\\
    
    (c) From \textit{(b)}, it suffices to show that $D\coloneqq D_1-D_2$ is principal $\Longleftrightarrow\ \LL(D)\simeq\OO_X$.
    \begin{itemize}
        \item ($\Rightarrow$) Clearly, if $D\sim0$, then $\LL(D)\simeq\OO_X$. Indeed, if $D=\{(X,f)\}$, then $\LL(D)$ is globally generated by $f^{-1}$.
        \item ($\Leftarrow$) If $\LL(D)\simeq\OO_X$, say $f\in\LL(D)(X)\longleftrightarrow1\in\OO_X(X)$. Thus, $\{(X,f^{-1})\}$ is principal.
    \end{itemize}
\end{proof}

\begin{corollary}{6.14}
    For any scheme $X\in\Sch$, we have an injective group homomorphism
    \begin{align*}
        \CaCl(X)\longrightarrow\Pic(X),\quad D\longmapsto\LL(D),
    \end{align*}
    where $\CaCl(X)$ denotes the group of Cartier divisors modulo linear equivalence.
\end{corollary}

\begin{proposition}{6.15}
    If $X$ is an integral scheme, then the homomorphism in \textit{(6.14)} is an isomorphism.
\end{proposition}
\begin{proof}\mbox{}
    \begin{itemize}
        \item From \textit{6.14}, it suffices to show that such homomorphism $\CaCl(X)\to\Pic(X)$ is surjective. But, by \textit{(6.13a)}, we have only to show that every invertible sheaf is isomorphic to a subsheaf of $\KK_X$.
        \item Note that $X$ is integral, it follows from \textit{(6.Ciii)} that $\KK_X$ is the constant sheaf of $K=K(X)$.
        \item For any $\LL\in\Pic(X)$, we consider $\MM\coloneqq\LL\otimes_{\OO_X}\KK_X\in\MOD_{\OO_X}$. Since $\LL$ is invertible, $\exists\ \{U_i\}$ open cover of $X$ such that $\LL|_{U_i}\simeq\OO_X|_{U_i}$. Then, $\MM|_{U_i}\simeq\LL|_{U_i}\otimes_{\OO_X|_{U_i}}\KK_X|_{U_i}\simeq\KK|_{U_i}$ is the constant sheaf of $K$ for every $i$. Further, $X$ is irreducible, so $\MM\simeq\KK_X$ is the constant sheaf of $K$ by \textit{(6.Cii)}.
        \item Therefore, we have the natural injective morphism $\LL\xhookrightarrow{}\LL\otimes_{\OO_X}\KK_X\simeq\KK_X$, i.e. we may regard $\LL$ as a subsheaf of $\KK_X$.
    \end{itemize}
\end{proof}

\begin{corollary}{6.16}
    If $X$ is an integral, separated, Noetherian, locally factorial scheme, then there is a natural isomorphism $\Cl(X)\isoto\Pic(X)$.
\end{corollary}
\begin{proof}
    Indeed, $\Cl(X)\underset{\text{(6.11)}}{\simeq}\CaCl(X)\underset{\text{(6.15)}}{\simeq}\Pic(X)$.
\end{proof}

\begin{corollary}{6.17}
    If $X=\PP_k^n$ for any field $k$, then every invertible sheaf on $X$ is isomorphic to $\OO_X(l)$ for some $l\in\mathbb{Z}$.
\end{corollary}
\begin{proof}
    Note that $X=\PP_k^n$ satisfies the condition of \textit{(6.11)}, so we have the following 1-1 correspondence:
    \begin{center}
        \begin{tikzcd}
            \mathbb{Z}\rar[leftrightarrow,"\sim","\text{(6.4)}"'] &\Cl(X)\rar[leftrightarrow,"\sim","\text{(6.11)}"'] &\CaCl(X)\rar[leftrightarrow,"\sim","\text{(6.15)}"'] &\Pic(X),\\
            1\rar[mapsto] &H=V_+(x_0)\rar[mapsto] &D=\{(U_i,\frac{x_i}{x_0})\}_{i=0}^n\rar[mapsto] &\LL(D).
        \end{tikzcd}
    \end{center}
    Indeed, we note that
    \begin{itemize}
        \item[(6.4):] Write $S\coloneqq k[x_0,\cdots,x_n]$, then the function field $K=S_{((0))}=(\Frac S)_0$. Further, $x_0$ corresponds to the prime ideal $\pp=(x_0)$ of height $1$, so we have the DVR given by $S_{(\pp)}=k[\{\frac{x_j}{x_i}:i\neq0\}]$ with its maximal ideal generated by $\frac{x_0}{x_i}$. Thus, $\nu_H(\frac{x_0}{x_i})=1$.
        \item[(6.11):] Set $U_i\coloneqq D_+(x_i)$, the usual open cover of $\Proj S$. Then,
        \begin{align*}
            \frac{x_i}{x_0}\in(\Frac S)_0^\times=K^\times=\KK_X^\ast(U_i)\quad\text{and}\quad\frac{x_i}{x_j}\in k[\{\frac{x_kx_l}{x_ix_j}:0\leq k,l\leq n\}]^\times=S_{(x_ix_j)}^\times\simeq\OO_X^\ast(U_i\cap U_j).
        \end{align*}
        Further, $\Gamma_\ast(\OO_X)\simeq S$ by \textit{(5.13)}, i.e. $\OO_X(d)(X)=S_d$ for every $d\geq0$.
        \item[(6.15):] Now, the subsheaf $\LL(D)\subset\KK_X$ is generated by $\{\frac{x_0}{x_i}\}_{i=1}^n$.
    \end{itemize}
    We show that $\LL(D)\simeq\OO_X(1)$ so that the generator $1\in\mathbb{Z}$ corresponds to the generator $\OO_X(1)\in\Pic(X)$. We consider the morphism $\psi:\LL(D)\to\OO_X(1)$ given by 
    \begin{align*}
        \psi_{U_i}:\LL(D)(U_i)\longrightarrow\OO_X(1)(U_i)\simeq(S_{x_i})_1,\quad\frac{x_0}{x_i}\longmapsto\frac{x_0}{1}
    \end{align*}
    for each $i$, which is clearly an isomorphism by the definition of $\LL(D)$. Thus, $\psi:\LL(D)\isoto\OO_X(1)$, and therefore, we have the group isomorphism $(\mathbb{Z},+)\isoto(\Pic(X),\otimes),\ l\mapsto\OO_X(l)$.
\end{proof}

\begin{definition}{(Effective Cartier Divisors)}
    \begin{itemize}
        \item A Cartier divisor $\{(U_i,f_i)\}\in\CDiv(X)$ is \textit{effective} if $f_i\in\Gamma(U_i,\OO_X|_{U_i})$ for every $i$.
        \item In this case, we define the \textit{associated subscheme} $Y$ to the closed subscheme of codimension $1$ defined by the sheaf of ideals $\II$ which is locally generated by $f_i$.
    \end{itemize}
\end{definition}

\begin{remark}{6.17.1}
    If $X\in\Sch$ satisfies the condition of \textit{(6.11)}, then the definition above gives a 1-1 correspondence under the isomorphism in \textit{(6.11)}:
    \begin{align*}
        \{D\in\Div(X):\text{effective}\}\ \longleftrightarrow\ \{D\in\CDiv(X):\text{effective}\}.
    \end{align*}
\end{remark}

\begin{proposition}{6.18}
    Let $D\in\CDiv(X)$ be an effective Cartier divisor and let $Y$ be the associated \textit{locally principal} closed subscheme (i.e. whose sheaf of ideals is locally generated by a single element). Then $\II_Y\simeq\LL(-D)$.
\end{proposition}
\begin{proof}
    Note that $\LL(-D)$ is a subsheaf of $\KK_X$ locally generated by $f_i$ on $U_i$. As $D$ is effective, we have $f_i\in\OO_X(U_i)$ for each $i$, which implies that $\LL(-D)$ is in fact a subsheaf of $\OO_X$. Therefore, $\LL(-D)\simeq\II_Y$, the ideal sheaf of $Y$.
\end{proof}

\chapter{Projective Morphisms}
\section{Morphisms to $\PP^n$}

\begin{theorem}{7.1}
    Let $A\in\Ring$ and $X\in\Sch_A$.
    \begin{itemize}
        \item[(a)] If $X\xrightarrow[\Sch_A]{\phi}\PP_A^n$ is an $A$-morphism, then $\phi^\ast(\OO(1))$ is an invertible sheaf on $X$ generated by the global sections $s_i=\phi^\ast(x_i)$ ($0\leq i\leq n$).
        \item[(b)] Conversely, if $\LL$ is an invertible sheaf on $X$, and if $\{s_i\}_{i=0}^n$ are the global sections generating $\LL$, then there exists a unique $A$-morphism $X\xrightarrow[\Sch_A]{\phi}\PP_A^n$ such that $\LL\simeq\phi^\ast(\OO(1))$ and $s_i=\phi^\ast(x_i)$ under this isomorphism.
    \end{itemize}
\end{theorem}
\begin{proof}
    (a) Write $S=A[x_0,\cdots,x_n]$.
    \begin{itemize}
        \item From \textit{(5.12)}, $\OO(1)$ on $\PP_A^n$ is invertible, and from \textit{(5.13)}, $S\simeq\Gamma_\ast(\OO)$, so $x_i\in S_1=\Gamma(\PP_A^n,\OO(1))$ for each $i$.
        \item For any $\pp\in\PP_A^n=\Proj S$, we have 
        \begin{align*}
            \OO(1)_\pp=S(1)_{(\pp)}=(S_\pp)_1=(S_\pp)_0\spn\{\frac{x_i}{1}\}_{i=0}^n=\OO_\pp\spn\{\frac{x_i}{1}\}_{i=0}^n,
        \end{align*}
        i.e. $\OO(1)$ is generated by global sections $x_0,\cdots,x_n$.
        \item Recall from the \textit{proof of [5.1d(i)]} that if $\EE\in\MOD_{\OO_Y}$ is locally free of rank $r$ covered by $\{V_i\}$, then $f^\ast\EE\in\MOD_{\OO_X}$ is also locally free of rank $r$ covered by $\{U_i\}$ where $U_i=f^{-1}(V_i)$ under the morphism $(X,\OO_X)\xrightarrow[\LRS]{f}(Y,\OO_Y)$. Therefore, the $\OO_X$-module $\LL=\phi^\ast(\OO(1))$ is indeed an invertible sheaf.
        \item For any $x\in X$, we see that
        \begin{align*}
            \LL_x\simeq\OO(1)_{\phi(x)}\otimes_{\OO_{\phi(x)}}\OO_{X,x}=\OO_{X,x}\spn\{\frac{x_i}{1}\otimes1\}_{i=0}^n,
        \end{align*}
        whence $\LL$ is generated by the global sections $s_i=\phi^\ast(x_i)$.
    \end{itemize}
    
    (b) We construct such morphism in the following steps:
    \begin{itemize}
        \item For each $i$, set $X_i\coloneqq\{P\in X:(s_i)_P\notin\mm_P\LL_P\}$, which is an open subset of $X$ (as we have seen before). Then, $\{X_i\}_{i=0}^n$ is an open cover of $X$. Indeed, for any $P\in X$, since $\{s_i\}_{i=0}^n$ generates $\LL$, we have $\LL_P=\OO_{X,P}\spn\{(s_i)_P\}_{i=0}^n$, say $1=\sum_{i=0}^na_i(s_i)_P$ (where $a_i\in\OO_{X,P}$). There must be at least an $i$ for which $(s_i)_P\notin\mm_P\LL_P$, i.e. $P\in X_i$, otherwise we have $1=\sum_{i=0}^na_i(s_i)_P\in\mm_P\LL_P$, a contradiction.
        \item Let $U_i\coloneqq D_+(x_i)$ ($0\leq i\leq n$) be the usual open affine cover of $Y=\PP_A^n=\Proj S$. We first construct a morphism $X_i\xrightarrow[\Sch]{\phi_i}U_i$.
        \begin{itemize}
            \item We have a trivialization of the invertible sheaf
            \begin{align*}
                \psi_i:\OO_X|_{X_i}\isolongto\LL|_{X_i},\quad1\longmapsto s_i
            \end{align*}
            on the open cover $\{X_i\}$.
            \item Note that $U_i\simeq\Spec A[y_0,\cdots,y_n]$ where $y_j=\frac{x_j}{x_i}$ with $y_i=1$ ommited. We define a ring homomorphism
            \begin{align*}
                f_i:A[y_0,\cdots,y_n]\longrightarrow\Gamma(X_i,\OO_{X_i}),\quad y_j\longmapsto\frac{s_j}{s_i}
            \end{align*}
            and making it $A$-linear. Note that for any $P\in X_i$, we have $(s_i)_P\notin\mm_P\LL_P\simeq\mm_P\OO_{X_i,P}$ (as $\LL$ is invertible), which implies that $(\frac{s_j}{s_i})_P=\frac{(s_j)_P}{(s_i)_P}$ is defined in $\OO_{X_i,P}$. Thus, $\frac{s_j}{s_i}\in\Gamma(X_i,\OO_{X_i})$, i.e. $f_i$ is well-defined.
            \item From \textit{[2.4]}, $f_i$ determines the morphism $\phi_i:X_i\xrightarrow[\Sch]{}\Spec A[y_0,\cdots,y_n]\simeq U_i$ uniquely.
            \item \textit{Check.} We may glue $\{\phi_i\}_{i=0}^n$ along $\{X_i\}_{i=0}^n$ to obtain a morphism $X\xrightarrow[\Sch]{\phi}Y=\PP_A^n$, i.e. it suffices to show that $\phi_i|_j=\phi_j|_i$ for every $i,j$. Write $y_k\coloneqq\frac{x_k}{x_i}$ and $y_l'\coloneqq\frac{x_l}{x_j}$, then
            \begin{gather*}
                A[y_0,\cdots,y_n]_{y_j}\simeq A[x_0,\cdots,x_n]_{(x_ix_j)}\simeq A[y_0',\cdots,y_n']_{y_i'},\\
                U_{ij}=D_+(x_ix_j)\simeq\Spec A[x_0,\cdots,x_n]_{(x_ix_j)}.
            \end{gather*}
           We have the correspondence of morphisms and ring homomorphisms:
            \begin{center}
                \begin{tikzcd}
                    U_i &X_i\lar["\phi_i"'] &A[y_0,\cdots,y_n]\rar["f_i"]\dar[] &\Gamma(X_i,\OO_{X_i})\dar[]\\
                    U_{ij}\uar[hook]\dar[hook] &X_{ij}\lar[shift left,"\phi_j|_i"]\lar[shift right,"\phi_i|_j"']\uar[hook]\dar[hook] &A[x_0,\cdots,x_n]_{(x_ix_j)}\rar[shift left,"f_i|_j"]\rar[shift right,"f_j|_i"'] &\Gamma(X_{ij},\OO_{X_{ij}})\\
                    U_j &X_j\lar["\phi_j"'] &A[y_0',\cdots,y_n']\rar["f_j"]\uar[] &\Gamma(X_j,\OO_{X_j})\uar[]
                \end{tikzcd}
            \end{center}
            where the both restrictions $f_i|_j:\frac{y_k}{y_j}=\frac{x_ix_k}{x_ix_j}\mapsto\frac{s_is_k}{s_is_j}$ and $f_j|_i:\frac{y_l'}{y_i'}=\frac{x_lx_j}{x_ix_j}\mapsto\frac{s_ls_j}{s_is_j}$ are compatible with $f_{ij}:=\frac{x_kx_l}{x_ix_j}\mapsto=\frac{s_ks_l}{s_ix_j}$. So, $f_{ij}$ determines the middle morphism $U_{ij}\leftarrow X_{ij}$ uniquely, and hence, $\phi_i|_j=\phi_j|_i$.
        \end{itemize}
        \item \textit{Claim.} $\LL\simeq\phi^\ast(\OO(1))$ as invertible $\OO_X$-modules. Note that both $\LL$ and $\phi^\ast(\OO(1))$ are locally free, so it suffices to prove on the open cover $\{X_i\}_{i=0}^n$ of $X$. We recall some facts:
        \begin{itemize}
            \item For $f\in S_1$, we have $S(n)_{(f)}=f^n\cdot S_{(f)}$ (as $S_{(f)}$-modules).
            \item If $X\xrightarrow[\Sch]{\phi}Y$ is a morphism, and $s\in\Gamma(Y,\LL)$ is a global section, there is a canonical morphism $\phi^{!}\LL\xrightarrow[\MOD(\OO_Y)]{}\phi_\ast(\phi^\ast\LL)$. The pull back of the section $s$ is defined to be
            \begin{align*}
                \phi^\ast(s)\coloneqq\phi^{!}_Y(s)\in\Gamma(Y,f_\ast(f^\ast\LL))=\Gamma(X,f^\ast\LL).
            \end{align*}
        \end{itemize}
        Now, we have the isomorphisms
        \begin{align*}
            \LL|_{X_i}\xrightarrow[\sim]{\psi_i^{-1}}\OO_X|_{X_i}\xrightarrow[\sim]{s_i\cdot}s_i\cdot\OO_X|_{X_i}\simeq\phi_i^\ast(x_i\cdot\OO|_{U_i})\simeq\phi_i^\ast(\OO(1)|_{U_i})\simeq(\phi^\ast\OO(1))|_{X_i}
        \end{align*}
        on each $X_i$, which glue to the isomorphism $\LL\simeq\phi^\ast\OO(1)$ of $\OO_X$-modules.
        
        
        % \begin{itemize}
        %     \item As $U_i\simeq\Spec S_{(x_i)}$ is affine, we have
        %     \begin{align*}
        %         \Hom_{\OO_Y|_{U_i}}(\OO(1)|_{U_i},\phi_\ast\LL|_{U_i}) &\underset{\text{[5.3]}}{\simeq}\Hom_{S_{(x_i)}}(S(1)_{(x_i)},\Gamma(U_i,\phi_\ast\LL|_{U_i}))\\ &=\Hom_{S_{(x_i)}}((S_{x_i})_1,\Gamma(X_i,\LL)).
        %     \end{align*}
        %     Clearly, we have an $S_{(x_i)}$-module homomorphism
        %     \begin{align*}
        %         (S_{x_i})_1\longrightarrow\Gamma(X_i,\LL),\quad\frac{x_j}{1}\ (j\neq i)\ \longmapsto\frac{s_j}{s_i},
        %     \end{align*}
        %     This uniquely determines a morphism $\OO(1)|_{U_i}\xrightarrow[\OO_Y|_{U_i}]{\psi_i}\phi_\ast\LL|_{U_i}$ for each $i$. Patching these $\psi$ together, we get a morphism $\OO(1)\xrightarrow[\MOD_{\OO_Y}]{\psi}\phi_\ast\LL$.
        %     \item Further, from the \textit{Adjoint Property}, we have
        %     \begin{align*}
        %         \Hom_{\MOD_{\OO_Y}}(\OO(1),\phi_\ast\LL)\simeq\Hom_{\MOD_{\OO_X}}(\phi^\ast(\OO(1)),\LL),
        %     \end{align*}
        %     so $\phi$ corresponds to a morphism $\phi^\ast(\OO(1))\xrightarrow[\MOD_{\OO_X}]{\theta}\LL$ uniquely. Now, for each $x\in X$, say $\phi(x)=\pp\in\Proj S$, then the stalk map
        %     \begin{align*}
        %         \theta_x:\phi^\ast(\OO(1))_x=\OO_Y(1)_\pp\otimes_{\OO_{Y,\pp}}\OO_{X,x}\longrightarrow\LL_x
        %     \end{align*}
        %     is given by $\frac{x_i}{1}\otimes1\mapsto(s_i)_x$, which is clearly an isomorphism as $\{s_i\}_{i=0}^n$ generates $\LL$ globally. Therefore, $\theta:\phi^\ast(\OO(1))\isoto\LL$ is an isomorphism and $s_i$ can be identified with $\phi^\ast(x_i)$ via $\theta$.
        % \end{itemize}
        \item Finally, any morphism with these properties must be the one we constructed, so such $A$-morphism $\phi$ is unique. 
    \end{itemize}
\end{proof}

\begin{proposition}{7.2}
    Let $X\xrightarrow[\Sch]{\phi}\PP_A^n$ be an $A$-morphism corresponding to an invertible sheaf $\LL\in\MOD_{\OO_X}$ with generating sections $s_0,\cdots,s_n\in\Gamma(X,\LL)$ as in \textit{(7.1)}. Then $\phi$ is a closed immersion if and only if
    \begin{itemize}
        \item[(1)] Each open set $X_i=X_{s_i}$ is affine. 
        \item[(2)] For each $i$, the ring homomorphism 
        \begin{align*}
            f_i:A[y_0,\cdots,y_n]\longrightarrow\Gamma(X_i,\OO_{X_i}),\quad y_j\longmapsto\frac{s_j}{s_i}
        \end{align*}
        constructed above is surjective.
    \end{itemize}
\end{proposition}
\begin{proof}
    ($\Rightarrow$) Suppose $\phi$ is a closed immersion. Then, we have the fibre square:
    \begin{center}
        \begin{tikzcd}
            X_i\FibSq{dr}\rar[hook,"\text{open}"]\dar["\phi_i"'] &X\dar["\phi"',"\text{closed}"]\\
            U_i\rar[hook,"\text{open}"'] &\PP_A^n\\
        \end{tikzcd}
    \end{center}
    Then, $X_i=U_i\cap X$ is a closed subscheme of $U_i$. Note that $U_i\simeq\Spec S_{(x_i)}$ is affine, by \textit{(5.10)}, $X_i$ is also affine, determined uniquely by an ideal $\mathfrak{a}\triangleleft S_{(x_i)}$, so that the ring homomorphism $f_i$ is given by $S_{(x_i)}\twoheadrightarrow S_{(x_i)}/\mathfrak{a}$, and thus, is surjective.\\
    
    ($\Leftarrow$) Suppose (1) and (2) are satisfied. Then the morphism $\phi:X_i\to U_i$ is a closed immersion by \textit{[2.18c]} (since $\phi_i$ is induced by $f_i$). Note that $X_i=\phi^{-1}(U_i)$ covers $X$ and that \textit{closed immersions are local on the base}, so $\phi:X\to\PP_A^n$ is a closed immersion.
\end{proof}

\begin{fact}{7.A}
    Let $X\xrightarrow[\Sch]{f}Y$ be a closed morphism of schemes of finite type over any field $k$. Then,
    \begin{align*}
        f\ \text{is injective}\ \Longleftrightarrow\ f\ \text{is injective on closed points}.
    \end{align*}
\end{fact}

We will use this fact in the following proposition.

\newpage
\begin{proposition}{7.3}
    Let $X\in\Sch_k$ be a projective scheme over an algebraically closed field $k$, and let $\phi:X\to\PP_k^n$ be a $k$-morphism corresponding to $\LL$ and $s_0,\cdots,s_n\in\Gamma(X,\LL)$ as above. Let $V\subset\Gamma(X,\LL)$ be the $k$-vector subspace spanned by $\{s_i\}_{i=0}^n$. Then $\phi$ is a closed immersion if and only if 
    \begin{itemize}
        \item[(1)] Elements of $V$ "separate points", i.e. for any two distinct closed points $P,Q\in X$, there is an $s\in V$ such that $s_P\in\mm_P\LL_P$ but $s_Q\notin\mm_Q\LL_Q$, or vice versa.
        \item[(2)] Elements of $V$ "separate tangent vectors", i.e. for each closed point $P\in X$, we have
        \begin{align*}
            \mm_P\LL_P/\mm_P^2\LL_P=k\spn\{\bar{s}_P:s\in V,\ s_P\in\mm_P\LL_P\}.
        \end{align*}
    \end{itemize}
\end{proposition}
\begin{proof}
    ($\Rightarrow$) If $\phi$ is a closed immersion, we may consider $X$ as a closed subscheme of $\PP_k^n$. From \textit{(5.16)}, $X\simeq\Proj S/I$ for some homogeneous ideal $I\triangleleft S=k[x_0,\cdots,x_n]$. So, $\phi$ is induced by the natural map $S\twoheadrightarrow S/I$, and thus, $\LL=\OO_X(1)$ with $s_i=x_i+I\in(S/I)_1$.
    \begin{itemize}
        \item[(1)] Given distinct closed points $P,Q\in X$, corresponding to the homogeneous maximal ideals $\mm_P\neq\mm_Q$ of $S/I$. Note that $S/I=k[s_0,\cdots,s_n]$ is finitely generated by $(S/I)_1$ over $(S/I)_0=k$, so $\exists\ \bar{f}\in(S/I)_1$ for which $\bar{f}\in\mm_P$ but $\bar{f}\notin\mm_Q$, i.e. $P\in V_+(\bar{f})$ but $Q\notin V_+(\bar{f})$. Take $s\coloneqq\bar{f}\in k\spn\{s_i\}_{i=0}^n=V$, then we have the desired result.
        \item[(2)] For any closed point $P\in X\subset\PP_k^n$. As $k$ is algebraically closed, we may assume $P=[a_0:\cdots:a_n]$ with $a_i\in k$. Further, we can always apply a suitable automorphism of $\PP_k^n$, see \textit{(7.1.1)}, and assume $P=[1:0:\cdots:0]$. It follows that  $P=(0,\cdots,0)$ on the open affine $U_0=\Spec k[y_1,\cdots,y_n]$ ($y_i\coloneqq x_i/x_0$), and that $\LL|_{U_0}\simeq\OO_X|_{U_0}$. Now, $P\in U_0$ corresponds to the maximal ideal 
        \begin{align*}
            \mm_{U_0,P}=\<y_1,\cdots,y_n\>\triangleleft k[y_1,\cdots,y_n]=S_{(x_0)}=\OO_{U_0,P}.
        \end{align*}
        Finally, since $\phi$ is a closed immersion, the map $\phi_P^\#:\OO_{U_0,P}=\OO_{\PP_k^n,P}\to\OO_{X,P}\simeq\LL_P$ is surjective. Therefore, the image of $\{\bar{y}_i\}_{i=1}^n$ under $\phi_P^\#$ generates $\mm_P/\mm_P^2$ over $k$.
    \end{itemize}
    
    ($\Leftarrow$) Conversely, suppose $\phi$ satisfies (1) and (2).
    \begin{itemize}
        \item Note that $X$ is projective over $k$, so it is proper over $k$ by \textit{(4.9)}, i.e. the composition $X\xrightarrow{\phi}\PP_k^n\to\Spec k$ is proper. Further, $\PP_k^n\to\Spec k$ is separated, it follows from \textit{(4.8e)} that $\phi$ is also proper. Thus, by \textit{[4.4]}, $\phi$ is a closed map (morphism) and its image $\phi(X)\subset\PP_k^n$ is closed.
        \item If $P\in X$ is a closed point, then $\LL\simeq\phi^\ast(\OO(1))$ induces the isomorphism
        \begin{align*}
            \OO(1)_{\phi(P)}/\mm_{\phi(P)}\OO(1)_{\phi(P)}\isolongto\LL_P/\mm_P\LL_P,\quad(x_i)_{\phi(P)}+\mm_{\phi(P)}\OO(1)_{\phi(P)}\longmapsto(s_i)_P+\mm_P\LL_P
        \end{align*}
        of $k$-vector spaces. Hence, if $f\in S_1$, we have
        \begin{align*}
            [\phi^\ast(f)]_P\in\mm_P\LL_P\ \Longleftrightarrow\ [f]_{\phi(P)}\in\mm_{\phi(P)}\OO(1)_{\phi(P)}\ \Longleftrightarrow\ [f]_{\phi(P)}=0\ \text{in}\ \kappa(\phi(P))=k.
        \end{align*}
        \item Now, if $P,Q\in X$ are distinct closed points, by (1), $\exists\ s\in V$ for which $s_P\in\mm_P\LL_P$ but $s_Q\notin\mm_Q\LL_Q$. This means, $\exists\ f\in S_1$ for which $\phi^\ast(f)=s$ and that $[f]_{\phi(P)}=0$ but $[f]_{\phi(Q)}\neq0$ in $k$. Thus, $\phi(P)=\phi(Q)$, i.e. $\phi$ is injective on closed points. By \textit{(7.A)}, $\phi$ is injective (as a map of sets), and thus, $\phi$ is a topological closed embedding.
        \item Write $Y=\PP_k^n$, it remains to show that the sheaf map $\phi^\#:\OO_Y\to\phi_\ast\OO_X$ is surjective. We check on the stalk maps. However, $\phi$ maps every closed point $X$ to a closed point of $Y$, so it reduces to show that the map $\phi_P^\#:\OO_{Y,\phi(P)}\to\OO_{X,P}$ is surjective for every closed point $P\in X$. We complete by the following lemma \textit{(7.4)}. 
        \item[(i)] Clearly, we have $\OO_{Y,\phi(P)}/\mm_{Y,\phi(P)}\simeq k\simeq\OO_{X,P}/\mm_{X,P}$.
        \item[(ii)] By (2) and the fact that $\LL$ is invertible, we see that
        \begin{multline*}
            \mm_{Y,\phi(P)}=k\spn\{[f]_{\phi(P)}:f\in S_1\}\\ \twoheadrightarrow
            k\spn\{[\phi^\ast(f)]_P+\mm_P:f\in S_1\}
            =k\spn\{\bar{s}_P:s\in V,\ s_P\in\mm_P\}=\mm_P/\mm_P^2.
        \end{multline*}
        \item[(iii)] Note from \textit{[4.9e]} that $\phi$ is also a projective morphism, and that $\OO_X\in\qcoh(X)$, it follows from \textit{(5.20)} that $\phi_\ast\OO_X\in\coh(Y)$. Therefore, $\OO_{X,P}\in\fg\MOD_{\OO_{Y,\phi(P)}}$.
    \end{itemize}
\end{proof}

\begin{lemma}{7.4}
    Let $f:A\xrightarrow[\Ring]{}B$ be a local homomorphism of local Noetherian rings such that
    \begin{itemize}
        \item[(i)] $A/\mm_A\isolongto B/\mm_B$ is a ring isomorphism.
        \item[(ii)] $\mm_A\twoheadrightarrow\mm_B/\mm_B^2$ is surjective.
        \item[(iii)] $B$ is a finitely generated $A$-module.
    \end{itemize}
    Then, $f$ is surjective.
\end{lemma}
\begin{proof}
    Consider the ideal $\bb\coloneqq\mm_AB\triangleleft B$, we have $\bb\subset\mm_B$. From (ii), we have $\mm_B=\bb+\mm_B^2$, it follows from \textit{Nakayama Lemma} for $(B,\mm_B)$ that $\bb=\mm_B$. Now, apply \textit{Nakayama Lemma} to $B\in\MOD_A$. Since $B\in\fg\MOD_A$ by (iii), $1\in B$ gives a generator for $B/\mm_AB=B/\mm_B\underset{\text{(i)}}{=}A/\mm_A$, i.e. $1\in B$ generates $B$ as an $A$-module. Thus, $f$ is surjective.
\end{proof}






\newpage
\section{Ample Invertible Sheaves}
\begin{definition}{(Ample Invertible Sheaves)}
    Fix a Noetherian scheme $X$.\\ An invertible sheaf $\LL\in\MOD_{\OO_X}$ is \textit{ample} if for every $\FF\in\coh(X)$, there is an integer $n_0>0$ (depending on $\FF$) such that the sheaf $\FF\otimes\LL^{\otimes n}$ is globally generated for every $n\geq n_0$.
\end{definition}

\begin{fact}{7.B}
    \begin{itemize}
        \item[(i)] Let $S$ be a graded ring, which is finitely generated by $S_1$ over a Noetherian ring $S_0=A$. Let $X=\Proj S$, then $\OO_X(1)$ is a very ample invertible sheaf relative to $\Spec A$.
        \item[(ii)] Let $X\in\Sch_A$ and $\LL\in\MOD_{\OO_X}$ be an invertible sheaf. Then, $\LL$ is very ample relative to $A\ \Longleftrightarrow$ There exists $n>0$ and global sections $\{s_i\}_{i=0}^n$, which generate $\LL$, such that the corresponding morphism $X\to\PP_A^n$ is an immersion.  
    \end{itemize}
\end{fact}
\begin{proof}
    (i) Note that $\Proj S\to\Spec A$ is projective. By \textit{(5.16.1)}, $\OO_X(1)$ is very ample over $A$. (ii) follows from \textit{(7.1)} and the \textit{definition of very ampleness}.
\end{proof}

\begin{proposition}{7.5}
    Let $\LL\in\MOD_{\OO_X}$ be an invertible sheaf on a Noetherian scheme $X$. T.F.A.E.
    \begin{itemize}
        \item[(i)] $\LL$ is ample. \qquad(ii) $\LL^{\otimes m}$ is ample for all $m>0$.\qquad(iii) $\LL^{\otimes m}$ is ample for some $m>0$.
    \end{itemize}
\end{proposition}
\begin{proof}
    (i) $\Rightarrow$ (ii). Given $\FF\in\coh(X)$, then $\FF\otimes\LL^{\otimes n}$ is globally generated for all large $n$, so is $\FF\otimes(\LL^{\otimes m})^{\otimes n}$. By definition, $\LL^m$ is ample.
    
    (ii) $\Rightarrow$ (iii). This is clear.
    
    (iii) $\Rightarrow$ (i). Suppose $\LL^{\otimes m}$ is ample. Given $\FF\in\coh(X)$, then $\exists\ n_0>0$ such that $\FF\otimes(\LL^{\otimes m})^{\otimes n}$ is globally generated for every $n\geq n_0$. Consider $\FF\otimes\LL\in\coh(X)$. As $\LL^{\otimes m}$ is ample, $\exists\ n_1>0$ such that $(\FF\otimes\LL)\otimes(\LL^{\otimes m})^{\otimes n}$ is globally generated for every $n\geq n_1$. Continuing this process, for each $1\leq k\leq m-1,\ \exists\ n_k>0$ such that $(\FF\otimes\LL^{\otimes k})\otimes(\LL^{\otimes m})^{\otimes n}\simeq\FF\otimes\LL^{\otimes(mn+k)}$ is globally generated for every $n\geq n_k$. Take $N\coloneqq m\cdot\max\{n_k:0\leq k\leq m-1\}>0$, then $\FF\otimes\LL^{\otimes n}$ is globally generated for every $n\geq N$. Therefore, $\LL$ is ample.
\end{proof}

\begin{theorem}{7.6}
    Let $X$ be a (Noetherian) scheme of finite type over a Noetherian ring $A$, and $\LL\in\MOD_{\OO_X}$ be an invertible sheaf. Then
    \begin{align*}
        \LL\ \text{is ample $\Longleftrightarrow\ \LL^{\otimes m}$ is very ample over $\Spec A$ for some $m>0$}.
    \end{align*}
\end{theorem}
\begin{proof}
    ($\Leftarrow$) Suppose $\LL^{\otimes m}$ is very ample over $A$. Let $X\xrightarrow[\Sch]{i}\PP_A^n$ be the immersion (with $n\geq1$) such that $\LL^{\otimes m}\simeq i^\ast(\OO(1))$. 
    \begin{itemize}
        \item Let $\overline{X}$ be the closure of $i(X)$ in $\PP_A^n$. Consider the closed immersion $\overline{X}\xrightarrow[\Sch]{j}\PP_A^n$ such that $i$ factors through $j$ via the open immersion $X\xrightarrow[\Sch]{k}\overline{X}$, i.e. $i=jk$. 
        \item Now, $\overline{X}\to\Spec A$ is projective. By \textit{(5.17)}, $\OO_{\overline{X}}(1)=j^\ast(\OO(1))$ is ample on $\overline{X}$.
        \item For any $\FF\in\coh(X)$, it follows from \textit{[]5.15]} that $\exists\ \overline{\FF}\in\coh(\overline{X})$ which extends the coherent sheaf $\FF$, i.e. $\overline{\FF}|_X\simeq\FF$.
        \item Now, if $\overline{\FF}\otimes\OO_{\overline{X}}(l)$ is globally generated, then so is
        \begin{align*}
            (\overline{\FF}\otimes\OO_{\overline{X}}(l))|_X &\simeq\FF\otimes(\OO_{\overline{X}}(1)|_X)^{\otimes l}\simeq\FF\otimes(k^\ast j^\ast\OO(1))^{\otimes l}\simeq\FF\otimes(i^\ast\OO(1))^{\otimes l}\\ &\simeq\FF\otimes(\LL^{\otimes m})^{\otimes l}
        \end{align*}
        for every large $l>0$. Thus, $\LL^{\otimes m}$ is ample, and hence, $\LL$ is ample by \textit{(7.5)}.
    \end{itemize}
    
    ($\Rightarrow$) Suppose $\LL$ is ample.\\
    
    (I) Given any $P\in X$. Choose an open affine $U\in\NN_P(X)$ on which $\LL|_U\simeq\OO_X|_U$.
    \begin{itemize}
            \item Let $Y\coloneqq X\setminus U$ be the closed subset in $X$, and let $\II_Y\triangleleft\OO_X$ be the sheaf of ideal with reduced induced structure. Since $X$ is Noetherian and of finite type over $A$, it follows from \textit{(5.9)} that $\II_Y\in\coh(X)$. So, $\II_Y\otimes\LL^{\otimes n}$ is globally generated by $\{s_i\}$ for some large $n>0$. Now, $(\II_Y\otimes\LL^{\otimes n})_P=\OO_{X,P}\spn\{(s_i)_P\}$, so we have $(s_i)_P\notin\mm_P(\II_Y\otimes\LL^{\otimes n})_P$ for some $i$. (Otherwise, we have $1=\sum_{\textit{finite}}a_i(s_i)_P\in\mm_P(\II_Y\otimes\LL^{\otimes n})_P$). Write $s\coloneqq s_i$.
            \item As $\LL$ is an invertible sheaf, $-\otimes\LL$ is an exact functor, whence $\II_Y\otimes\LL^{\otimes n}\xhookrightarrow{}\OO_X\otimes\LL^{\otimes n}\simeq\LL^{\otimes n}$, i.e. we may regard $s\in\Gamma(X,\II_Y\otimes\LL^{\otimes n})\subset\Gamma(X,\LL^{\otimes n})$. Consider the set $X_s\coloneqq\{Q\in X:s_Q\notin\mm_Q\LL_Q^{\otimes n}\}$, which is open in $X$, and that $P\in X_s$.
            \item Further, we have $X_s\subset U$. Indeed, if $Q\notin U$, then 
            \begin{align*}
                Q\in Y=\Supp(\OO_X/\II_Y)=\{x\in X:\II_{Y,x}\neq\OO_{X,x}\},
            \end{align*}
            i.e. $\II_{Y,Q}\triangleleftneq\OO_{X,Q}$, whence $\II_{Y,Q}\subset\mm_Q$. Thus, $s_Q\in\mm_Q\LL_Q^{\otimes n}$, i.e. $Q\notin X_s$.
            \item Moreover, as $U\simeq\Spec A$ is open affine and that $\LL|_U\simeq\OO_X|_U$, the section $s|_U\in\Gamma(U,\LL^{\otimes n})$ corresponds to a global section $f\in\Gamma(U,\OO_X|_U)\simeq\Gamma(U,\OO_X^{\otimes n})$. This implies that $X_s=U\cap X_s=U_f\simeq\Spec A_f$ is affine.
            \item \textit{Conclusion.} For any $P\in X,\ \exists\ n>0$ and a global section $s\in\Gamma(X,\LL^{\otimes n})$ such that $X_s\in\NN_P(X)$ is a open affine neighborhood of $P$.
        \end{itemize}
        
    (II) Reduction to $\LL$.
    \begin{itemize}
            \item Note that $X$ is of finite type over $A$, so $X$ is quasi-compact, i.e. we can cover $X$ by finitely many such open affine $\{X_{s_i}\}_{i=1}^r$ with each $s_i\in\Gamma(X,\LL^{\otimes n_i})$.
            \item For any $k>0$, note that 
            \begin{align*}
                Q\in X_{s_i}\ \Longleftrightarrow\ (s_i)_Q\notin\mm_Q\LL_Q^{\otimes n_i}\ \Longleftrightarrow\ (s_i^k)_Q\notin\mm_Q\LL_Q^{\otimes kn_i}\ \Longleftrightarrow\ Q\in X_{s_i^k},
            \end{align*}
            which shows that $X_{s_i}=X_{s_i^k}$ where $s_i^k\in\Gamma(X,\LL^{\otimes kn_i})$. Now, pick $n\coloneqq n_1\cdots n_r>0$ and $t_i\coloneqq s_i^{n/n_i}\in\Gamma(X,\LL^{\otimes n})$ for each $i$. Then $X_i\coloneqq X_{t_i}=X_{s_i}$ still gives a finite open affine cover of $X$.
            \item From \textit{(7.5)}, $\LL^{\otimes n}$ is ample. Moreover, note that we are only trying to show that $\LL^{\otimes m}$ is very ample over $A$ for some $m>0$. So, we may replace $\LL^{\otimes n}$ by $\LL$ and assume that $X$ is covered by open affines $X_i=X_{s_i}$ where $s_i\in\Gamma(X,\LL)$ ($1\leq i \leq r$).
        \end{itemize}
    
    (III) Define an $A$-morphism $X\to\PP_A^N$.
    \begin{itemize}
        \item For each $i$, we let $B_i\coloneqq\Gamma(X_i,\OO_X|_{X_i})$, i.e. $X_i=\Spec B_i$. Since $X$ is of finite type over $A$, we have $B_i\in\fg\ALG_A$ by \textit{[3.3]}. Let $\{b_{ij}\}_{j=1}^{k_i}$ denote the set of generators for $B_i$ over $A$. Apply \textit{(5.14b)} to each $b_{ij}\in\Gamma(X_i,\OO_X)$, then $\exists\ m_{ij}>0$ and a global section $c_{ij}\in\Gamma(X,\OO_X\otimes\LL^{\otimes m_{ij}})=\Gamma(X,\LL^{\otimes m_{ij}})$ such that $c_{ij}|_{X_i}=s_i^{m_{ij}}b_{ij}$. As usual, we take $m\geq\max\{m_{ij}\}$ such that $s_{ij}^mb_{ij}$ extends to $c_{ij}\in\Gamma(X,\LL^{\otimes m})$ for every $i,j$.
        \item Note that the global sections $\{s_i^{m}\}_{i=1}^r$ and $\{c_{ij}\}_{i,j}$ generated the invertible sheaf $\LL^{\otimes m}\in\MOD_{\OO_X}$. By \textit{(7.1)}, they define an $A$-morphism $\phi:X\to\PP_A^N$ such that $\LL^{\otimes m}\simeq\phi^\ast(\OO(1))$, (where $N=r+\sum_{i=1}^rk_i-1$).
    \end{itemize}
    
    (IV) Show that $\phi$ is an immersion.
    \begin{itemize}
        \item Let $\{x_i\}_{i=1}^r$ and $\{x_{ij}\}_{i,j}$ denote the set of homogeneous coordinates of $\PP_A^N$ via $\phi$, i.e. $s_i^m=\phi^\ast(x_i)$ and $c_{ij}=\phi^\ast(x_{ij})$ by \textit{(7.1)}. 
        \item For each $i$, let $U_i\coloneqq D_+(x_i)\subset\PP_A^N$ (which is open affine), then $\phi^{-1}(U_i)=X_{s_i^m}=X_i$. Further, the corresponding ring homomorphism
        \begin{align*}
            \phi^{\#}_{U_i}:A[\{y_k\};\{y_{kj}\}]\longrightarrow B_i,\quad y_k=\frac{x_k}{x_i}\longmapsto\frac{s_k^m}{s_i^m},\quad y_{kj}=\frac{x_{kj}}{x_i}\longmapsto\frac{c_{kj}}{s_i^m}=b_{kj}
        \end{align*}
        is surjective, since $y_{ij}\mapsto b_{ij}$ and $\{b_{ij}\}_{j=1}^{k_i}$ generates $B_i$ over $A$.
        \item By \textit{[2.18c]}, the morphism $X_i\to U_i$ between affine schemes via $\phi$ is a closed immersion. Thus, $X=\bigcup_{i=1}^rX_i\xrightarrow{\phi}\bigcup_{i=1}^rU_i$ is a closed immersion followed by an open immersion $\bigcup_{i=1}^rU_i\xhookrightarrow{}\PP_A^N$, i.e. $X\xrightarrow[\Sch]{\phi}\PP_A^N$ is an immersion. Therefore, $\LL^{\otimes m}$ is very ample over $A$.
    \end{itemize}
\end{proof}

\begin{example}{7.6.1}
    Let $X=\PP_A^n$ for any Noetherian ring $A$. Then,
    \begin{align*}
        \OO(l)\ \text{is ample}\ \Longleftrightarrow\ \OO(l)\ \text{is very ample}\ \Longleftrightarrow\ l>0.
    \end{align*}
\end{example}
\begin{proof}
    For any ring $A$, set $S=A[x_0,\cdots,x_n]$. Given $d>0$, let $\{M_i\}_{i=0}^N$ denote the set of monomials in $x_0,\cdots,x_n$ of degree $d$ (in a given order where $N=\binom{n+d}{n}-1$). Then the ring homomorphism 
    \begin{align*}
        \theta:A[y_0,\cdots,y_N]\longrightarrow S^{(d)},\quad y_i\longmapsto M_i
    \end{align*}
    gives rise to the closed immersion $\Proj S^{(d)}\xhookrightarrow{}\PP_A^N$ by \textit{[3.12d]}. This corresponds to the $d$-uple embedding
    \begin{align*}
        \rho_d:X=\Proj S\isolongto Y=\Proj S^{(d)}\xhookrightarrow{}\PP_A^N
    \end{align*}
    by \textit{[5.13]}. It follows that $\OO_X(d)\simeq\OO_Y(1)$ is very ample over $A$. In particular, if $A$ is Noetherian, then $\OO_X(n)$ is ample for every $n>0$ by \textit{(7.5)} and \textit{(7.6)}. On the other hand, the sheaf $\OO_X(l)$ has no global sections for $l<0$, one sees easily that $\OO_X(l)$ cannot be ample for $l\leq0$. Thus, we have the desired result. 
\end{proof}

\newpage
\section{Linear Systems}
\begin{nt}{}
    Throughout this subsection, we always assume $X$ is a \textit{nonsingular} (all local rings are regular) projective variety over an algebraically closed field $k$. Then,
    \begin{itemize}
        \item Recall that \textit{"regular $\Rightarrow$ UFD"}, so $X$ is nonsingular implies that $X$ is locally factorial. By \textit{(6.11)}, we have $\Div(X)\simeq\CDiv(X)$, and $\Cl(X)\simeq\CaCl(X)$.
        \item As $X$ is integral, $\Cl(X)\simeq\Pic(X)$ by \textit{(6.15)}.
        \item For any invertible $\OO_X$-module $\LL$, note that $\LL\in\coh(X)$. It follows from \textit{(5.19)} that $\Gamma(X,\LL)$ is a finite-dimensional $k$-vector space.
    \end{itemize}
\end{nt}

\begin{definition}{(Divisor of Zeros)}
    Fix an invertible sheaf $\LL\in\MOD_{\OO_X}$. Let $0\neq s\in\Gamma(X,\LL)$ be a nonzero global section. Choose an open cover $\{U_i\}$ of $X$ on which $\phi_i:\LL|_{U_i}\isoto\OO_X|_{U_i}$. Then, $\phi_i(s|_i)\in\Gamma(U_i,\OO_X)\subset\KK_X(U_i)$ is nonzero (since $\LL(U)\to\LL(V)$ is injective whenever $V\subset U$ as $X$ is integral). Thus, $\{(U_i,\phi_i(s|_i))\}$ defines an effective Cartier divisor (depending only on $\LL$ and $s$), called the \textit{divisor of zeros} of $s$, and is denoted by $(s)_0$.
\end{definition}

\begin{proposition}{7.7}
    Let $X$ be a nonsingular projective variety over the algebraically closed field $k$.\\ Fix $D_0\in\CDiv(X)$ and let $\LL$ denote the corresponding invertible sheaf $\LL(D_0)\in\Pic(X)$. Then,
    \begin{itemize}
        \item[(a)] For each $0\neq s\in\Gamma(X,\LL)$, the divisor of zeros $(s_0)\in\CDiv(X)$ is effective, and $(s)_0\sim D_0$.
        \item[(b)] Every effective divisor $D\in\CDiv(X)$ linearly equivalent to $D_0$ is $(s)_0$ for some $s\in\Gamma(X,\LL)$.
        \item[(c)] Given two nonzero sections $s,s'\in\Gamma(X,\LL)$. Then, $(s)_0=(s')_0\ \Longleftrightarrow\ s'=\lambda s$ for some $\lambda\in k^\times$.
    \end{itemize}
\end{proposition}
\begin{proof}
    (a) Given a nonzero global section $s\in\Gamma(X,\LL)$.
    \begin{itemize}
        \item Recall that $\LL=\LL(D_0)$ is an $\OO_X$-submodule of $\KK_X$, so $s\in\Gamma(X,\KK)$ is nonzero, i.e. we have  $s\in\KK^\ast(X)=K^\times$.
        \item Suppose $D_0$ is represented by $\{(U_i,f_i)\}$ where $f_i\in\KK^\ast(U_i)=K^\times$. We may choose $U_i\neq\emptyset$ such that $\LL|_{U_i}\simeq\OO_X|_{U_i}$, and that $\LL(U_i)=\LL(D_0)(U_i)$ is generated by $f_i^{-1}$.
        \item We have a local isomorphism $\phi:\LL(D_0)\to\OO_X$ for which $\phi_{U_i}:\LL(U_i)\isoto\OO_X(U_i)$ sends $s|_i\mapsto f_i\cdot s|_i$. Therefore, $(s)_0$ is determined by $\{(U_i,f_i\cdot s|_i)\}=D_0+\dv(s)$, whence $(s)_0\sim D_0$.
    \end{itemize}
    
    (b) Since $D\sim D_0$, we choose an $s\in\KK^\ast(X)$ such that $D-D_0=\dv(s)$. Then, $\dv(s)+D_0=D\geq0$ (i.e. effective), and $\dv(s)\geq-D_0$. Note that $D_1\geq D_2\ \Longleftrightarrow\ \LL(D_1)\supset\LL(D_2)$. Now, $\LL(-D_0)\subset\LL(\dv(s))=\OO_X\spn\{s^{-1}\}$, whence $s\in\Gamma(X,\LL(D_0))$ is nonzero. Therefore, from the proof of \textit{(a)}, we have $D=D_0+\dv(s)=(s)_0$.\\
    
    (c) Note that $X\in\Var_k$ is a projective variety, so $\OO_X(X)=k$ by \textit{(I, 3.4)}. Thus, 
    \begin{align*}
        (s)_0=(s')_0\ \Longleftrightarrow\ \dv(s)=\dv(s')\ \Longleftrightarrow\ s/s'\in\OO_X^\ast(X)=\OO_X(X)^\times=k^\times.
    \end{align*}
\end{proof}

\begin{definition}{(Linear Systems)}
    Fix a divisor $D_0\in\CDiv(X)$, and let $\LL=\LL(D_0)$.
    \begin{itemize}
        \item A \textit{complete linear system} of $D_0$ on $X$ is the set
        \begin{align*}
            |D_0|\coloneqq\{D\in\CDiv(X):\text{effective such that}\ D\sim D_0\},
        \end{align*}
        which is from \textit{(7.7)} in 1-1 correspondence with the set
        \begin{align*}
            (\Gamma(X,\LL)\setminus\{0\})/k^\times\overset{\text{homeom.}}{\simeq}\PP(\Gamma(X,\LL)).
        \end{align*}
        \item A \textit{linear system} $\dd$ of $D_0$ on $X$ is a subset of $|D_0|$, which is a linear subspace for $\PP(\Gamma(X,\LL))$.
        Note that the map $(\cdot)_0:\Gamma(X,\LL)\to\CDiv(X)$ gives a 1-1 correspondence:
        \begin{align*}
            \{V\subset\Gamma(X,\LL):k\text{-vector subspace}\} &\longleftrightarrow\ \{\dd\subset|D_0|:\text{linear system}\},\\ \{s\in\Gamma(X,\LL):(s)_0\in\dd\}\cup\{0\} &\longmapsfrom\qquad \dd,\\ V\qquad &\longmapsto\ \{(s)_0:0\neq s\in V\}.
        \end{align*}
        \item The \textit{dimension} of the linear system $\dd$ is defined to be $\dim\dd\coloneqq\dim\PP(V)=\dim V-1$.
    \end{itemize}
\end{definition}

\begin{definition}{(Base Point)}
    \newline A point $P\in X$ is a \textit{base point} of a linear system $\dd$ if $P\in\Supp D$ for every $D\in\dd$. Here, $\Supp D$ denotes the set $\Supp D=\bigcup\{Y\subset X:\text{prime divisor of}\ D\}$.
\end{definition}

\begin{fact}{7.C}
    For any nonzero $s\in\Gamma(X,\LL)$, the support $\Supp(s)_0=X\setminus X_s$, where $X_s$ denotes the open subset $X_s=\{P\in X:s_P\notin\mm_P\LL_P\}$.
\end{fact}
\begin{proof}
    This question is local, so we may assume $X=\Spec A$ so that $\LL=\OO_X$. Then, we have $0\neq s\in A,\ X_s=D(s)$, so $X\setminus X_s=V(s)$. Note that
    \begin{align*}
        s_P\notin\mm_P\LL_P=\mm_P\ \Longleftrightarrow\ s(P)=s_P\mod\mm_P\neq0\quad\text{in}\ k(P)=\OO_{X,P}/\mm_P=k.
    \end{align*}
    Thus, there are no prime divisors on $X_s$, which implies that the prime divisors of $(s)_0$ are the irreducible components of $V(s)$, whence $\Supp(s)_0=V(s)$.
\end{proof}

\begin{lemma}{7.8}
    Let $\dd$ be a linear system on $X$ corresponding to the $k$-vector subspace $V\subset\Gamma(X,\LL)$. Then,
    \begin{itemize}
        \item[(i)] A point $P\in X$ is a base point of $\dd\ \Longleftrightarrow\ s_P\in\mm_P\LL_P$ for every $s\in V$.
        \item[(ii)] $\dd$ is base-point-free  $\Longleftrightarrow\ \LL$ is generated by the global sections in $V$.
    \end{itemize}
\end{lemma}
\begin{proof}
    (i) This is clear from \textit{(7.C)}. For (ii), we see that:
    \begin{align*}
        \dd\subset|D_0|\ \text{is base-point-free} &\Longleftrightarrow\ \text{For every}\ P\in X,\ P\ \text{is not a base point}.\\  &\Longleftrightarrow\ \text{For every}\ P\in X,\ \exists\ s\in V\ \text{such that}\ s_P\notin\mm_P\LL_P.\\  &\Longleftrightarrow\ \text{For every}\ P\in X,\ \LL_P=\OO_{X,P}\spn\{s:s\in V\}.\\  &\Longleftrightarrow\ \LL\ \text{is globally generated by}\ \{s:s\in V\}.
    \end{align*}
\end{proof}
\newpage

\section{The Relative Proj \& Projective Space Bundles}
\begin{axiom}{($\dagger$)}
    \begin{itemize}
        \item $X$ is a Noetherian scheme, and $\Ss\in\qcoh(X)$ having a structure of graded $\OO_X$-algebras, i.e. we have $\Ss\simeq\bigoplus_{d\geq0}\Ss_d$, where $\Ss_d$ is the \textit{homogeneous part} of \textit{degree} $d$.
        \item Assume further that $\Ss$ is locally generated by $\Ss_1\in\coh(X)$ as an $\Ss_0=\OO_X$-algebra. Thus, we have $\Ss_d\in\coh(X)$ for every $d\geq0$.
    \end{itemize}
\end{axiom}

\begin{summary}{The Relative $\PROJ$ Construction}
    \newline Fix a scheme $X$, and a quasi-coherent algebra $\Ss\in\GR\ALG_{\OO_X}$.
    \begin{itemize}
        \item For any open affine $U=\Spec A$ of $X$, we have $\Ss(U)\in\GR\ALG_{\OO_X(U)}$. Consider $\Proj\Ss(U)$ with the natural morphism
        \begin{align*}
            \Proj\Ss(U)\xrightarrow[\Sch]{\pi_U}\Spec\OO_X(U)\simeq U.
        \end{align*}
        Checking through the \textit{glueing data}, we may glue these $\{\pi_U\}$ along such open affine cover $\{U\}$ to obtain a scheme, denoted by $\PROJ\Ss$, together with a morphism $\PROJ\Ss\xrightarrow[\Sch]{\pi}X$ and open immersions $\Proj\Ss(U)\xrightarrow[\Sch]{\psi_U}\PROJ\Ss$, such that $\{\Im\psi_U\}$ covers $\PROJ\Ss$, and is compatible on the intersections.
        \item Such $\pi$ is unique with the property that for any open affine $U=\Spec A$ of $X$, the square
        \begin{center}
            \begin{tikzcd}
                \Proj\Ss(U)\FibSq{dr}\rar["\psi_U",hook]\dar["\pi_U"] &\PROJ\Ss\dar["\pi"]\\ U\rar["\iota_U",hook] &X
            \end{tikzcd}
        \end{center}
        is a fibre square, which is compatible with all open affine subsets $V\subset U$. Further, we have
        \begin{align*}
            \pi^{-1}(U)=\Im\psi_U\simeq\Proj\Ss(U).
        \end{align*}
        \item From \textit{(5.12c)}, the invertible sheaves $\OO(1)$ on each $\Proj\Ss(U)$ are compatible under this construction, so they glue together to give an invertible sheaf $\OO(1)$ on $\PROJ\Ss$ canonically.
    \end{itemize}
\end{summary}

\begin{fact}{7.D}
    Consider the relative construction $\PROJ\Ss\xrightarrow[\Sch]{\pi}X$ as above. Then, we have the following results.
    \begin{itemize}
        \item[(i)] $\pi$ is separated.
        \item[(ii)] If $\Ss\in\ALG_{\OO_X}$ is locally finitely generated, then $\pi$ is of finite type.
        \item[(iii)] If, in addition to (ii), that $X\Sch$ is Noetherian, then so is $\PROJ\Ss$.
        \item[(iv)] If, in addition to (iii), that $\Ss$ satisfies ($\dagger$), then $\pi$ is proper.
    \end{itemize}
\end{fact}
\begin{proof}
    These are in fact the more detailed results of \textit{(7.10a)}. To see (i), we use the fact of \textit{(4.1)} and the fact that "separated is local on the base \textit{(4.6f)}", whence each $\pi_U$ is separated, and so is $\pi$.
\end{proof}

\begin{example}{7.8.7}
    If $\Ss=\OO_X[T_0,\cdots,T_n]$ is the polynomial algebra, then $\PROJ\Ss=\PP_X^n$ with its twisting sheaf $\OO(1)$ defined earlier.
\end{example}

\begin{lemma}{7.9}
    Let $X$ and $\Ss$ satisfy ($\dagger$), and let $\LL\in\MOD_{\OO_X}$ be an invertible sheaf. We define a new sheaf
    \begin{align*}
        \Ss\ast\LL\coloneqq\bigoplus_{d\geq0}(\Ss_d\otimes_{\OO_X}\LL^{\otimes d})\in\GR\ALG_{\OO_X}.
    \end{align*}
    Then $\Ss\ast\LL$ also satisfies ($\dagger$), and there is a canonical isomorphism
    \begin{align*}
        \phi:Y'=\PROJ(\Ss\ast\LL)\xrightarrow[\Sch]{~~\sim~~}Y=\PROJ\Ss,
    \end{align*}
    commuting with the projections $Y\xrightarrow[\Sch]{\pi}X$ and $Y'\xrightarrow[\Sch]{\pi'}X$, and having the property that 
    \begin{align*}
        \OO_{Y'}(1)\simeq\phi^\ast\OO_Y(1)\otimes_{\OO_{Y'}}(\pi')^\ast\LL.
    \end{align*}
\end{lemma}
\begin{proof}
    We divide into the following steps.\\
    
    (I) For any open affine $U\subset X$ satisfying $\LL|_U\simeq\OO_X|_U,\ \exists$ canonical $U$-isomorphism
    \begin{align*}
        \phi_U:\Proj(\Ss\ast\LL)(U)\xrightarrow[\Sch_U]{\sim}\Proj\Ss(U),
    \end{align*}
    compatible with all open affine subsets $V\subset U$.
    
    \textsc{proof.} Clearly, we have $\Ss\ast\LL\in\qcoh(X)$. Given an $\OO_X(U)$-basis $\{\eta\}$ of $\LL(U)$. Define a map 
    \begin{align*}
        \alpha_U:\Ss(U)\longrightarrow(\Ss\ast\LL)(U),\quad s=\sum_{d\geq0}s_d\longmapsto\sum_{d\geq0}s_d\otimes\eta^{\otimes d}.
    \end{align*}
    This is a graded $\OO_X(U)$-algebra isomorphism since for each $d\geq0$, the map $\Ss_d(U)\to(\Ss\ast\LL)_d,\ s\mapsto s\otimes\eta^{\otimes d}$ comes from the isomorphism $(\Ss_d\otimes\LL^{\otimes d})|_U\simeq\Ss_d|_U\otimes\LL|_U^{\otimes d}\simeq\Ss_d|_U\otimes\OO_X|_U^{\otimes d}\simeq\Ss_d|_U$. Thus, $\alpha_U$ induces an isomorphism $\phi_U:\Proj(\Ss\ast\LL)(U)\xrightarrow[\Sch_U]{\sim}\Proj\Ss(U)$.
    
    \textit{Check.} The isomorphism $\phi_U$ is independent of the choice of the basis $\eta$.
    \begin{itemize}
        \item Suppose $\{\eta'\}$ is another $\OO_X(U)$-basis of $\LL(U)$, then $\exists\ f\in\OO_X(U)^\times$ for which $\eta'=f\cdot\eta$. We also have a graded $\OO_X(U)$-algebra isomorphism 
        \begin{align*}
            \alpha'_U:\Ss(U)\xrightarrow{~~\sim~~}(\Ss\ast\LL)(U),\quad s=\sum_{d\geq0}s_d\longmapsto\sum_{d\geq0}s_d\otimes\eta'^{\otimes d}.
        \end{align*}
        It follows that $\alpha'_U(s)=\sum_{d\geq0}f^d(s_d\otimes\eta^{\times d})=\sum_{d\geq0}\alpha_U(s_d)$.
        \item \textit{Fact.} For any homomorphism $S\xrightarrow[\GR\Ring]{\alpha}T$, set $G(\psi)\coloneqq\{\pp\in\Proj T:\pp\not\supset\alpha(S_+)\}$ (as in \textit{[2.14]}). Suppose $f\in T_0$ is a unit, we consider the graded homomorphism
        \begin{align*}
            \alpha':S\longrightarrow T,\quad s=\sum_{d\geq0}s_d\longmapsto\sum_{d\geq0}f^d\alpha(s_d).
        \end{align*}
        Then, $G(\alpha)=G(\alpha')$, and the induced morphisms $G(\alpha)\xrightarrow[\Sch]{}\Proj S$ are the same.
        \item Now, the question reduces to prove the fact above. Consider the graded homomorphism
        \begin{align*}
            \theta:T\longrightarrow T,\quad t=\sum_{d\geq0}t_d\longmapsto\sum_{d\geq0}f^dt_d.
        \end{align*}
        As $f\in T_0$ is a unit, $\theta$ has an inverse homomorphism given by $\sum_{d\geq0}t_d\mapsto\sum_{d\geq0}f^{-d}t_d$. Thus, $\theta$ is an automorphism. Now, $\alpha'=\theta\alpha$, so the result clearly follows.
    \end{itemize}
    Finally, for the \textit{compatibility}, note that if $V\subset U$ is affine, $\{\eta|_V\}$ is an $\OO_X(V)$-basis for $\LL(V)$.\\
    
    (II) The isomorphisms $\phi_U$ glue to give an isomorphism $\phi:\PROJ(\Ss\ast\LL)\xrightarrow[\Sch]{\sim}\PROJ\Ss$ commuting with the projections, i.e. the following diagram commutes
    \begin{center}
        \begin{tikzcd}
            \PROJ(\Ss\ast\LL)\arrow[rr,"\phi"]\arrow[rd,"\pi'"'] &\Comm{d} &\PROJ\Ss\arrow[ld,"\pi"]\\ &X &
        \end{tikzcd}
    \end{center}
    \textsc{proof.} Cover $X$ with open affines $U$ satisfying $\LL|_U\simeq\OO_X|_U$. For each such $U$, we have the composition
    \begin{align*}
        \Proj(\Ss\ast\LL)(U)\xrightarrow[\sim]{~~\phi_U~~}\Proj\Ss(U)\xrightarrow[]{~~\pi_U~~}U.
    \end{align*}
    Since $\phi_U$ is compatible with every open affine subset $V\subset U$ by (I), they agree on the intersections, which implies that these $\phi_U$ can be glued. Thus, $\exists$ isomorphism $\PROJ(\Ss\ast\LL)\xrightarrow[\Sch]{\phi}\PROJ\Ss$ such that
    \begin{center}
        \begin{tikzcd}
            \PROJ(\Ss\ast\LL)\rar["\phi", "\sim"']\Comm{dr} &\PROJ\Ss\\ \Proj(\Ss\ast\LL)(U)\rar["\phi_U"',"\sim"]\uar["\psi_U'"] &\Proj\Ss(U)\uar["\psi_U"']
        \end{tikzcd}
    \end{center}
    which is in fact a fibre square, where $\psi_U,\psi_U'$ are defined in the construction. Now, $\pi\phi$ is another morphism $\Proj(\Ss\ast\LL)\xrightarrow[\Sch]{}X$ satisfying $(\pi\phi)\psi'_U=\pi\psi_U\phi_U=\iota_U(\pi_U\phi_U)$. From the \textit{uniqueness} of the relative construction $\PROJ$, we have $\pi'=\pi\phi$.\\
    
    (III) \textit{Check.} $\Ss\ast\LL\in\GR\ALG_{\OO_X}$ is locally generated by $(\Ss\ast\LL)_1=\Ss_1\otimes\LL$ over $\OO_X$, and thus, $\Ss\ast\LL$ satisfies ($\dagger$).\\
    
    (IV) Show the isomorphism of sheaves.
    
    \textsc{proof.} Cover $X$ with open affines $U$ satisfying $\LL|_U\simeq\OO_X|_U$. For each such $U$, consider the following commutative diagram
    \begin{center}
        \begin{tikzcd}[column sep=small]
            \Proj(\Ss\ast\LL)(U)\arrow[rr,"\phi_U","\sim"']\arrow[dr,"\pi'_U"']\arrow[dd,hook,"\psi'_U"] & &
            \Proj\Ss(U)\arrow[dd,hook,"\psi_U"]\arrow[ld,"\pi_U"]\\
            &U\\
            \PROJ(\Ss\ast\LL)\arrow[rr,near start,"\phi","\sim"']\arrow[dr,"\pi'"'] & &\PROJ\Ss\arrow[ld,"\pi"]\\
            &X\arrow[leftarrow,crossing over]{uu}[hook,near start,description]{\iota_U}
        \end{tikzcd}
    \end{center}
    Write $Y=\PROJ\Ss,\ Y'=\PROJ(\Ss\ast\LL)$ and $Y_U=\Proj\Ss(U),\ Y'_U=\Proj(\Ss\ast\LL)(U)$. Then we have the following two isomorphisms:
    \begin{align*}
        \Big(\phi^\ast\OO_Y(1)\Big)\Big|_{Y'_U} &\simeq\phi^\ast\Big(\OO_Y(1)|_{Y_U}\Big)\simeq\phi^\ast(\phi_U)_\ast\OO_{Y_U}(1)\simeq(\phi_U')_\ast\phi_U^\ast\OO_{Y_U}(1)\\ &\simeq(\phi'_U)_\ast\OO_{Y_U'}(1)\simeq\OO_{Y'}(1)|_{Y_U'}\tag{i}\\
        \Big(\pi'\,^\ast\LL\Big)\Big|_{Y'_U} &\simeq(\pi_U')^\ast(\LL|_U)\isolongto(\pi'_U)^\ast(\OO_X|_U)\simeq\OO_{Y'}|_{Y_U'}\tag{ii}
    \end{align*}
    Thus, the following isomorphism:
    \begin{multline*}
        \alpha_U:\Big(\phi^\ast\OO_Y(1)\otimes\pi'\,^\ast\LL\Big)\Big|_{Y'_U}\simeq\Big(\phi^\ast\OO_Y(1)\Big)\Big|_{Y'_U}\otimes\Big(\pi'\,^\ast\LL\Big)\Big|_{Y'_U}\isolongto\OO_{Y'}(1)|_{Y_U'}\otimes\OO_{Y'}|_{Y_U'}\simeq\OO_{Y'}(1)
    \end{multline*}
    for each such $U$. Further, this is independent of the chosen $\OO_X(U)$-basis for $\LL(U)$ (similar to the proof of (I)), and is compatible with all open affine subsets $V\subset U$. Thus, these $\{\alpha_U\}$ can be glued along such cover $\{U\}$ of $X$ and obtain an isomorphism $\alpha:\phi^\ast\OO_Y(1)\otimes\pi'\,^\ast\LL\xrightarrow[\GR\ALG_{\OO_{Y'}}]{\sim}\OO_{Y'}(1)$.
\end{proof}

\begin{fact}{7.E}
    Let $X$ be a Noetherian scheme and suppose $\FF\in\coh(X)$ is globally generated by the sections $\{s_i\}_{i\in I}$. Then $\FF$ can be generated by a finite subset of $\{s_i\}_{i\in I}$.
\end{fact}
\begin{proof}
    From assumption, $\exists$ surjective morphism $\OO_X^{\oplus I}\xrightarrow[\MOD_{\OO_X}]{\phi}\FF$. As $X$ is Noetherian, it can be covered by finitely many open affines $U_j=\Spec A_j$ ($1\leq j\leq n$) where each $A_j$ is Noetherian. Now, for each $1\leq j\leq n$, $\phi$ restricts to the surjective morphism
    \begin{align*}
        \widetilde{A}_j^{\oplus I}=\OO_X^{\oplus I}|_{U_j}\twoheadrightarrow\FF|_{U_j}\underset{\text{(5.4)}}{=}\widetilde{M}_j
    \end{align*}
    where $M_j\in\fg\MOD_{A_j}$. From \textit{(5.2)}, the homomorphism $A_j^{\oplus I}\xrightarrow[\MOD_{A_j}]{}M_j$ is also surjective. Since $M_j$ is finitely generated over the Noetherian ring $A_j$, it follows that $M_j$ is a Noetherain $A_j$-module, and thus, any $A_j$-submodule of $M_j$ is finitely generated. Thus, $\exists$ finite subset $J_j\subset I$ such that $\{s_i|_j\}_{i\in J_j}$ generates $M_j$, so $\FF_P=\OO_{X,P}\spn\{[s_i|_j]_P:i\in J_j\}$ for all $P\in U_j$. Finally, take $J\coloneqq\bigcup_{j=1}^nJ_j$, then $\FF$ is globally generated by finitely many sections $\{s_i\}_{i\in J}$.
\end{proof}

\begin{proposition}{7.10}
    Let $X,\ \Ss$ satisfy ($\dagger$), with $Y=\PROJ\Ss\xrightarrow[\Sch]{\pi}X$ and $\OO_Y(1)$ constructed as above. Then:
    \begin{itemize}
        \item[(a)] $\pi$ is proper, and thus, separated of finite type.
        \item[(b)] If $X$ admits an ample invertible sheaf $\LL$, then $\pi$ is projective, and $\OO_Y(1)\otimes\pi^\ast(\LL^{\otimes n})\in\MOD_{\OO_Y}$ is very ample over $X$ for suitable $n>0$.
    \end{itemize}
\end{proposition}
\begin{proof}
    (a) For any open affine subset $U\subset X$, note that $\Proj\Ss(U)\xrightarrow[\Sch]{\pi_U}U$ is projective \textit{(4.8.1)}, and thus, proper \textit{(4.9)}. But, \textit{proper is local on the base (4.8f)}, so $\pi$ is proper.\\
    
    (b) Suppose $\LL\in\MOD_{\OO_X}$ is an ample invertible sheaf. Since $\Ss_1\in\coh(X),\ \exists n>0$ such that $\Ss_1\otimes\LL^{\otimes n}$ is globally generated. Further, as $X$ is Noetherian, $\Ss_1\otimes\LL^{\otimes n}$ can be generated by finitely many global sections, i.e. $\exists$ surjective morphism $\OO_X^{\oplus(N+1)}\xrightarrow[\MOD_{\OO_X}]{}\Ss_1\otimes\LL^{\otimes n}$ for some $N>0$. This allows us to define a surjective morphism
    \begin{align*}
        \OO_X[T_0,\cdots,T_N]\xrightarrow[\GR\ALG_{\OO_X}]{}\bigoplus_{d\geq0}\Ss_d\otimes\LL^{\otimes nd}=\Ss\ast\LL^{\otimes n},
    \end{align*}
    which gives rise to a closed immersion (as closed immersion is local on the base) defined by
    \begin{align*}
        i:Y=\PROJ\Ss\xrightarrow[\text{(7.9)}]{\phi^{-1}\ \sim}\PROJ(\Ss\ast\LL^{\otimes n})\xhookrightarrow[\text{[3.12]}]{\text{closed}}\PROJ\OO_X[T_0,\cdots,T_N]\underset{\text{(7.8.7)}}{=}\PP_X^N.
    \end{align*}
    Therefore, $\pi:Y\xhookrightarrow{i}\PP_X^N\to X$ is projective.
    
    Now, using the notations in \textit{(7.9)} with $Y'=\PROJ(\Ss\ast\LL^{\otimes n})$, then $i$ is an immersion satisfying
    \begin{align*}
        i^\ast\OO_{\PP_X^N}(1) &\simeq(\phi^{-1})^\ast\OO_{Y'}(1)\underset{\text{(7.9)}}{\simeq}(\phi^{-1})^\ast\Big(\phi^\ast\OO_Y(1)\otimes\pi'\,^\ast\LL^{\otimes n}\Big)\\
        &\underset{\text{[5.1f*]}}{\simeq}(\phi^{-1})^\ast\phi^\ast\Big(\OO_Y(1)\otimes\pi^\ast\LL^{\otimes n}\Big)\simeq\OO_Y(1)\otimes\pi^\ast\LL^{\otimes n}.
    \end{align*}
    Thus, $\OO_Y(1)\otimes\pi^\ast\LL^{\otimes n}$ is very ample over $X$.
\end{proof}

\begin{definition}{(Projective Space Bundles)}
    Given a Noetherian scheme $X$ and  a locally free sheaf $\EE\in\coh(X)$. Let $\Ss\coloneqq S(\EE)=\bigoplus_{d\geq0}S^d(\EE)$, then $\Ss\in\GR\ALG_{\OO_X}$ satisfies ($\dagger$). We define the \textit{projective space bundle} associated to $\EE$ to be the space $\PP(\EE)=\PROJ\Ss$, which comes with a projection $\PP(\EE)\xrightarrow[\Sch]{\pi}X$ and an invertible sheaf $\OO(1)$.
\end{definition}

\begin{remark*}{}
    If $\EE\in\MOD_{\OO_X}$ is a free sheaf of rank $n+1$ over some open set $U\subset X$, then we have
    \begin{align*}
        \pi^{-1}(U)\simeq\Proj S(\EE)(U)\simeq\Proj\OO_X(U)[x_0,\cdots,x_n]\simeq\PP_U^n.
    \end{align*}
\end{remark*}

\begin{proposition}{7.11}
    Let $\widetilde{X}=\PP(\EE)\xrightarrow[\Sch]{\pi}X$ be a projective space bundle. Then
    \begin{itemize}
        \item[(a)] If $\rank_{\OO_X}\EE\geq2$, there is a canonical isomorphism $\Ss\xrightarrow[\GR\ALG_{\OO_X}]{\sim}\bigoplus_{l\in\mathbb{Z}}\pi_\ast(\OO_{\widetilde{X}}(l))$.
        \item[(b)] There is a natural surjective morphism $\pi^\ast\EE\xrightarrow[\MOD_{\OO_Y}]{\phi}\OO_Y(1)$.
    \end{itemize}
\end{proposition}
\begin{proof}
    (a) Say $\rank_{\OO_X}\EE=r+1$. Cover $X$ with open affines $U=\Spec A$ on which $\EE|_U\simeq\OO_X|_U^{\oplus r}$. For each such $U$, write $S=A[x_0,\cdots,x_r]$, then $\pi^{-1}(U)\simeq\Proj S$, so that the isomorphism
    \begin{align*}
        S\xrightarrow[\text{(5.13)}]{\sim}\Gamma_\ast(\OO_{\pi^{-1}(U)})=\bigoplus_{l\in\mathbb{Z}}\Gamma(\pi^{-1}(U),\OO_{\widetilde{X}}(l))=\bigoplus_{l\in\mathbb{Z}}\Gamma(U,\pi_\ast\OO_{\widetilde{X}}(l))
    \end{align*}
    of graded $S$-algebras induces the isomorphism
    \begin{align*}
        \Ss|_U\simeq S(\EE|_U)\simeq A[x_0,\cdots,x_r]\widetilde{}\xrightarrow[\text{(5.2)}]{\sim}\Big(\bigoplus_{l\in\mathbb{Z}}\Gamma(U,\pi_\ast\OO_{\widetilde{X}}(l))\Big)\widetilde{}\underset{\text{(5.15)}}{\simeq}\Big(\bigoplus_{l\in\mathbb{Z}}\pi_\ast\OO_{\widetilde{X}}(l))\Big)\Big|_U.
    \end{align*}
    Thus, these glue to give an isomorphism $\Ss\xrightarrow[\GR\ALG_{\OO_X}]{\sim}\bigoplus_{l\in\mathbb{Z}}\pi_\ast(\OO_{\widetilde{X}}(l))$.\\
    
    (b) This is just the relative version of \textit{(5.16.2)}, i.e. the twisting sheaf $\OO(1)$ on $\PP^n$ is globally generated by $\{x_i\}_{i=0}^n$.
\end{proof}

\begin{proposition}{7.12}
    Let $\widetilde{X}=\PP(\EE)\xrightarrow[\Sch]{\pi}X$ be a projective space bundle.\\ Then to give a morphism $Y\xrightarrow[\Sch_X]{f}\PP(\EE)$, it is equivalent to give an invertible sheaf $\LL\in\MOD_{\OO_Y}$ and a surjective morphism $g^\ast\EE\xrightarrow[\MOD_{\OO_Y}]{\psi}\LL$. \textit{Note.} This is a local version of (7.1).
\end{proposition}
\begin{proof}
    ($\Rightarrow$) Given a morphism $Y\xrightarrow[\Sch_X]{f}\PP(\EE)$. By \textit{(7.11b)}, $\exists$ surjective morphism $\pi^\ast\EE\xrightarrow[\MOD_{\OO_{\widetilde{X}}}]{\phi}\OO_{\widetilde{X}}(1)$. Applying the exact functor $f^\ast(-)$ to the morphism gives a surjective map 
    \begin{align*}
        g^\ast(\EE)=f^\ast\pi^\ast\EE\xrightarrow[\MOD_{\OO_Y}]{f^\ast\phi}f^\ast\OO_{\widetilde{X}}(1).
    \end{align*}
    Thus, we take $\psi=f^\ast\phi$ and $\LL=f^\ast\OO_{\widetilde{X}}(1)\in\MOD_{\OO_Y}$.\\
    
    ($\Leftarrow$) Given an invertible sheaf $\LL\in\MOD_{\OO_Y}$ with a surjective morphism $g^\ast\EE\xrightarrow[\MOD_{\OO_Y}]{\psi}\LL$. I claim that there is a unique morphism $Y\xrightarrow[\Sch_X]{f}\PP(\EE)=\widetilde{X}$ such that $\LL\simeq f^\ast\OO_{\widetilde{X}}(1)$, and that $\psi$ is obtained by applying $f^\ast$ to $\psi$.
    \begin{itemize}
        \item \textit{Note.} This is in fact a local question on $X$, so taking an open affine cover $\{U:\EE|_U\ \text{is free}\}$ of $X$ reduces the statement to \textit{(7.1)}.
        \item If $X=\Spec A$ and $\EE\simeq\OO_X^{\oplus(n+1)}$, then $\widetilde{X}=\PP_A^n$. Now that $g^\ast\EE\in\MOD_{\OO_Y}$ is free and $\psi$ is surjective, so $\psi$ is uniquely determined by $n+1$ global sections generating $\LL$, say $s_i\in\Gamma(Y,\LL)$ ($0\leq i\leq n$). From \textit{(7.1)}, these sections $\{s_i\}_{i=0}^n$ uniquely determine a morphism $Y\xrightarrow[\Sch_X]{f}\PP_A^n=\widetilde{X}$ satisfying $\LL\simeq f^\ast\OO_{\widetilde{X}}(1)$ and $s_i=f^\ast(x_i)$. 
        \item Further, there is a natural surjective morphism $\pi^\ast\EE\xrightarrow[\MOD_{\OO_Y}]{\phi}\OO_Y(1)$ by \textit{(7.11b)}. It follows that $f^\ast\phi$ is another surjective morphism $g^\ast\EE\xrightarrow[\MOD_{\OO_Y}]{}f^\ast\OO_{\widetilde{X}}(1)\simeq\LL$ determined by the global sections $\{s_i\}_{i=0}^n$ generating $\LL$. Thus, we have $\psi=f^\ast\phi$ by \textit{uniqueness}.
        \item Finally, the uniqueness of such morphism $f$ follows from \textit{(7.1)}.
    \end{itemize}
\end{proof}

\chapter{Differentials}
\section{K\"{a}hler Differential}
\begin{definition}{(Deriviations)}
    Given $A\in\Ring,\ B\in\ALG_A$ and $M\in\MOD_B$. An $A$-\textit{deriviation} of $B$ into $M$ is a map $d:B\to M$ satisfying:
    \begin{itemize}
        \item[(D1)] $d$ is additive.
        \item[(D2)] $d(bb')=bdb'+b'db$ for every $b,b'\in B$.
        \item[(D3)] $da=0$ for every $a\in A$.
    \end{itemize}
    The set of all $A$-deriviations of $B$ into $M$ is denoted by $\Der_A(B,M)$.
\end{definition}

\begin{definition}{(Modules of Relative Differential Forms)}
    Given $A$ and $B$ as above. The \textit{module of relative differential forms} of $B$ over $A$ is a pair $\<\Omega_{B/A},d\>$, where $\Omega_{B/A}\in\MOD_B$ and $d\in\Der_A(B,\Omega_{B/A})$, satisfying the \textit{Universal Property}: for any $M\in\MOD_B$, and any $d'\in\Der_A(B,M)$, there exists a unique homomorphism $f:\Omega_{B/A}\xrightarrow[\MOD_B]{}M$ such that $d'=fd$.
\end{definition}

\begin{remark}{(Existence of $\Omega_{B/A}$)}
    Let $F\in\MOD_B$ be a free module with a $B$-basis of symbols $\{db:b\in B\}$. Consider the set
    \begin{align*}
        F'\coloneqq B\spn\{d(b+b')-db-db',\ d(bb')-bdb'-b'db,\ da\ |\ b,b'\in B,\ a\in A\},
    \end{align*}
    which is a $B$-submodule of $F$. We define $\Omega_{B/A}\coloneqq F/F'\in\MOD_B$ with $A$-deriviation given by
    \begin{align*}
        d:B\longrightarrow F/F',\quad b\longmapsto db+F'.
    \end{align*}
    Further, we have a natural isomorphism $\Der_A(B,M)\underset{\Ab}{\isolongto}\Hom_B(\Omega_{B/A},M)$ in $M\in\MOD_B$.
\end{remark}
\begin{proof}
    We show that such construction satisfies the \textit{Universal Property}. Given any pair $\<M,d'\>$, we define
    \begin{align*}
        f:F/F'\longrightarrow M,\quad db+F'\longmapsto d'b,
    \end{align*}
    which is clearly $B$-linear and we have $fd=d'$. Suppose there is another homomorphism $f':F/F'\xrightarrow[\MOD_B]{}M$ such that $f'd=d'$. Then we have $f'd=d'=fd$. But $d$ is surjective, which implies that the homomorphism $f'=f$ is unique.
\end{proof}


\begin{proposition}{8.1}
    Let $B\in\ALG_A$, and let $f:B\otimes_AB\xrightarrow[\ALG_B]{} B$ be the natural "diagonal homomorphism" defined by $b\otimes b\mapsto b$, and set $I\coloneqq\ker f$. Regarding $B\otimes_AB\in\MOD_A$, then $I/I^2$ has a $B$-module structure given by
    \begin{align*}
        B\times I/I^2\longrightarrow I/I^2,\quad b\cdot(b'\otimes b''+I^2)\longmapsto bb'\otimes b''+I^2.
    \end{align*}
    Now, define a map 
    \begin{align*}
        d:B\longrightarrow I/I^2,\quad b\longmapsto(1\otimes b-b\otimes1)+I^2.
    \end{align*}
    Then, the pair $\<I/I^2,d\>$ gives a relative differential form $\Omega_{B/A}$.
\end{proposition}
\begin{proof}
    (I) \textit{Check.} $d\in\Der_A(B,I/I^2)$ is an $A$-deriviation.
    \begin{itemize}
        \item[(D1)] This is easy.
        \item[(D2)] Note that $(1\otimes b-b\otimes1)(1\otimes b'-b'\otimes1)\in I^2$, which implies $1\otimes bb'+bb'\otimes1=b\otimes b'+b'\otimes b$ in $I/I^2$. Thus, for any $b,b'\in B$, we have
        \begin{align*}
            bdb'+b'db &=b(1\otimes b'-b'\otimes1)+b'(1\otimes b-b\otimes1)+I^2\\ &=b\otimes b'-bb'\otimes1+b'\otimes b-bb'\otimes1+I^2\\ &=1\otimes bb'-bb'\otimes1+I^2=d(bb').
        \end{align*}
        \item[(D3)] Clearly, $da=1\otimes a-a\otimes1+I^2=0+I^2$ for all $a\in A$.
    \end{itemize}
    
    (II) \textit{Claim.} The pair $\<I/I^2,d\>$ satisfies the \textit{Universal Property} of $\Omega_{B/A}$.
    \begin{itemize}
        \item \textit{Claim}. The kernel $I$ is generated by $\{1\otimes b-b\otimes1:\ b\in B\}$.
        
        Clearly, $1\otimes b-b\otimes1\in I$ for every $b\in B$. On the other hand, we note that for any $b\otimes b'\in B\otimes_AB$, we have
        \begin{align*}
            b\otimes b'=bb'\otimes1+b(1\otimes b'-b'\otimes1).
        \end{align*}
        Now, suppose $\sum_ia_i(b_i\otimes b_i')\in I$, then $\sum_ia_ib_ib_i'=f(\sum_ia_i(b_i\otimes b_i'))=0$, whence
        \begin{align*}
            \sum a_i(b_i\otimes b_i')=\underbrace{\sum a_i(b_ib_i'\otimes1)}_{=0}+\sum a_ib_i(1\otimes b_i'-b_i'\otimes1)\in B\spn\{1\otimes b-b\otimes1:\ b\in B\}.
        \end{align*}
        \item For any pair $\<M,d'\>$ with $M\in\MOD_B$ and $d'\in\Der_A(B,M)$, we define
        \begin{align*}
            h:I/I^2\longrightarrow M,\quad1\otimes b-b\otimes1+I^2\longmapsto d'b,
        \end{align*}
        which is clearly $B$-linear and satisfies $d'=hd$. Further, $d$ is surjective, which shows that such $h$ must be unique.
    \end{itemize}
\end{proof}

\begin{proposition}{8.2}
    If $A',B\in\ALG_A$, and $B'\coloneqq B\otimes_AA'$. Then, $\Omega_{B'/A'}\simeq\Omega_{B/A}\otimes_BB'$ (as $B'$-modules). Further, if $S\subset B$ is multiplicatively closed, then $\Omega_{S^{-1}B/A}\simeq S^{-1}\Omega_{B/A}$.
\end{proposition}
\begin{proof}
    Let $\<\Omega_{B/A},d\>$ be the differential form of $B/A$. First, note that
    \begin{align*}
        \Omega_{B/A}\otimes_BB'=\Omega_{B/A}\otimes_B(B\otimes_AA')\simeq\Omega_{B/A}\otimes_AA'.
    \end{align*}
    \begin{itemize}
        \item We consider the map $d'\coloneqq d\otimes\mathds{1}:B'=B\otimes_AA'\to\Omega_{B/A}\otimes_AA'$, which clearly satisfies (D1) and (D3). For (D2), we see that
        \begin{align*}
            (d\otimes\mathds{1})((b_1\otimes a_1')(b_1\otimes a_2')) &=(d\otimes\mathds{1})(b_1b_2\otimes a_1'a_2')=d(b_1b_2)\otimes a_1'a_2'\\ &=(b_1db_2+b_2db_1)\otimes a_1'a_2'\\ &=b_1db_2\otimes a_1'a_2'+b_2db_1\otimes a_1'a_2'\\ &=(b_1\otimes a_1')(db_2\otimes a_2')+(b_2\otimes a_2')(db_1\otimes a_1'),
        \end{align*}
        for any $b_1\otimes a_1',b_2\otimes a_2'\in B\otimes_AA'=B$. Therefore, $d'\in\Der_{A'}(B',\Omega_{B/A}\otimes_AA')$.
        \item \textit{Claim.} The pair $\<\Omega_{B/A}\otimes_AA',d'\>$ satisfies the \textit{Universal Property} of $\Omega_{B'/A'}$.
        
        Given any $M\in\MOD_{B'}$ and any $D'\in\Der_{A'}(B',M)$. Clearly, the composition $D'i$ is an $A$-deriviation of $B$ into $M$, i.e. $D'i\in\Der_A(B,M)$, where $i:B\xhookrightarrow{}B',\ b\mapsto b\otimes1$ is the canonical injection of $B$-algebras. By the \textit{Universal Property} of $\Omega_{B/A}$, $\exists\ D\in\Der_A(\Omega_{B/A},M)$ for which $D'i=Dd$. Now, consider the $B'$-module homomorphism
        \begin{align*}
            E:\Omega_{B/A}\otimes_AA'\longrightarrow M,\quad c\otimes a\longmapsto a'Dc.
        \end{align*}
        Then, for any $b\otimes a'\in B'$, we have
        \begin{align*}
            Ed'(b\otimes a')=E(db\otimes a')=a'D(db)=a'D'i(b)=a'D'(b\otimes1)=D'(b\otimes a').
        \end{align*}
        Therefore, $Ed'=D'$ as required. The uniqueness is clear, so we get $\<\Omega_{B/A}\otimes_BB',d'\>\simeq\Omega_{B'/A'}$.
    \end{itemize}
\end{proof}

\begin{example}{8.2.1}
    Consider $B=A[x_1,\cdots,x_n]$. Then $\Omega_{B/A}\in\MOD_B$ is a free module of rank $n$ generated by $dx_1,\cdots,dx_n$. In general, for $B=A[\{x_i\}_{i\in I}]$, the $B$-module $\Omega_{B/A}$ is free with a $B$-basis $\{dx_i\}_{i\in I}$.
\end{example}
\begin{proof}
    Clearly, $\{dx_i\}_{i\in I}$ generates $\Omega_{B/A}$, so it suffices to show that $\{dx_i\}_{i\in I}$ is $B$-linearly independent. Suppose $\sum f_idx_i=0$ ($f_i\in B$). For any $j\in I$, consider the $j$th partial derivative $\partial/\partial x_j\in\Der_A(B,B)$. Then $\<B,\partial/\partial x_j\>$ gives a pair. By \textit{Universal Property} of $\Omega_{B/A},\ \exists$ unique $\Omega_{B/A}\xrightarrow[\MOD_B]{\theta_j}B$ such that $\theta_jd=\partial/\partial x_j$. It follows that $\theta_j(dx_i)=\partial x_i/\partial x_j=\delta_{ij}$ for every $i\in I$. Thus, we have
    \begin{align*}
        0=\theta_j(\sum f_idx_i)=\sum f_i\theta_j(dx_i)=\sum f_i\delta_{ij}=f_j
    \end{align*}
    for every $j\in I$. Therefore, $\{dx_i\}_{i\in I}$ is $B$-linearly independent.
\end{proof}

\begin{lemma}{8.A (Categorical Fact)}\newline
    Suppose we have morphisms $A\xrightarrow{v}B\xrightarrow{u}C$ in an Abelian category $\Af$ such that $uv=0$ and $u$ is epi. If the sequence
    \begin{align*}
        \Hom_\Af(A,T)\xleftarrow{-\circ v}\Hom_\Af(B,T)\xleftarrow{-\circ u}\Hom_\Af(C,T)
    \end{align*}
    is exact (in $\Ab$) for every $T\in\Ob\Af$. Then, $(B\xrightarrow{u}C)=\Coker(A\xrightarrow{v}B)$, and thus, (this is equivalent to that) the sequence $A\to B\to C\to0$ is exact in $\Af$.
\end{lemma}
\begin{proof}
    Consider $T\coloneqq\Coker(A\xrightarrow{v}B)\in\Af$ with $B\xrightarrow[\Af]{p}T$. Then, we have $pv=0$. By \textit{exactness}, $\exists\ \alpha\in\Hom_\Af(C,T)$ such that $\alpha u=p$. Note that the morphism $B\xrightarrow[\Af]{u}C$ satisfies $uv=0$, it follows from the \textit{Universal Property of Cokernel (or Coequalizer)}, $\exists$ unique $T\xrightarrow[\Af]{\theta}C$' with $\theta p=u$. Now, $u=\theta p=\theta\alpha u$, but $u$ is epi, so $\theta\alpha=\mathds{1}_C$; on the other hand, $\alpha\theta p=\alpha u=p$, but $\<T,p\>=\Coker v$, i.e. $p$ is epi, so $\alpha\theta=\mathds{1}_T$. Therefore, $C=T$, and thus, $(B\xrightarrow[\Af]{u}C)=\Coker(A\xrightarrow[\Af]{v}B)$. 
\end{proof}

\begin{remark*}{(8.A*)}
    As a special case of \textit{(8.A)}. Suppose the morphism $A\xrightarrow[\Af]{v}B$ is mono [resp. epi]. If one can show that the map $\Hom_\Af(A,\Coker v)\xleftarrow[\Ab]{-\circ v}\Hom_\Af(B,\Coker v)$ is epi [resp. $\Hom_\Af(A,A)\xleftarrow[\Ab]{-\circ v}\Hom_\Af(B,A)$ is mono], then $A\xrightarrow[\Af]{v}B$ is an isomorphism.
\end{remark*}

\begin{proposition}{8.3. (1st Fundamental Exact Sequence)}\newline
    Consider the ring homomorphisms $A\xrightarrow[]{\phi}\to B\xrightarrow[]{\psi}C$. Then there is a natural exact sequence of $C$-modules:
    \begin{align*}
        \Omega_{B/A}\otimes_BC\xrightarrow{~~v~~}\Omega_{C/A}\xrightarrow{~~u~~}\Omega_{C/B}\longrightarrow0.
    \end{align*}
\end{proposition}
\begin{proof}
    We apply the categorical fact \textit{(8.A)}. First, we define the maps
    \begin{align*}
        v:\Omega_{B/A}\otimes_BC &\longrightarrow\ \Omega_{C/A}, &u:\Omega_{C/A} &\longrightarrow\ \Omega_{C/B},\\
        d_{B/A}b\otimes c &\longmapsto c\cdot d_{C/A}\psi(b). & d_{C/A}c &\longmapsto\ d_{C/B}c,
    \end{align*}
    which are clearly $C$-linear.
    \begin{itemize}
        \item Note that $u$ is clearly surjective. Further, for any $b\in B$ and $c\in C$, we have
        \begin{align*}
            uv(d_{B/A}b\otimes c)=u(c\cdot d_{C/A}\psi(b))=c\cdot u(d_{C/A}\psi(b))=c\cdot\underbrace{d_{C/B}\psi(b)}_{=0}=0.
        \end{align*}
        \item Now, it suffices to show that the exactness of the sequence $\Hom_C(-,T)$ for any $T\in\MOD_C$, which can be rewrite as
        \begin{center}
            \begin{tikzcd}
                \Hom_C(\Omega_{B/A}\otimes_BC,T)\dar["\simeq"'] &\Hom_C(\Omega_{C/A},T)\lar["-\circ v"']\arrow[dd, equal] &\Hom_C(\Omega_{C/B},T)\lar["-\circ u"']\arrow[dd, equal]\\
                \Hom_B(\Omega_{B/A},T)\dar[equal] & &\\
                \Der_A(B,T) &\Der_A(C,T)\lar["-\circ\psi"'] &\Der_B(C,T)\lar["\alpha"'],
            \end{tikzcd}
        \end{center}
        where $\alpha$ is just the map $D\mapsto D$. Since we already have $uv=0$, it only remains to show that $\ker(-\circ\psi)\subset\Im\alpha$. Suppose $D\in\Der_A(C,T)$ with $D\psi=0$, i.e. the differential is $0$ on $B$, so we may extend $D$ to a $B$-deriviation of $C$, i.e. $D\in\Der_B(C,T)$.
    \end{itemize}
\end{proof}

\begin{proposition}{8.4. (2nd Fundamental Exact Sequence)}\newline
    Let $B\in\ALG_A,\ I\triangleleft B$, and $C=B/I$. Then there is a natural exact sequence of $C$-modules:
    \begin{align*}
        I/I^2\xrightarrow{~~\delta~~}\Omega_{B/A}\otimes_BC\xrightarrow{~~v~~}\Omega_{C/A}\longrightarrow0,
    \end{align*}
    where $\delta$ is a $C$-module homomorphism defined by
    \begin{align*}
        \delta:I/I^2\longrightarrow\Omega_{B/A}\otimes_BC,\quad b+I^2\longmapsto d_{B/A}b\otimes1.
    \end{align*}
\end{proposition}
\begin{proof}
    We again apply the fact \textit{(8.A)} and use the notations in \textit{(8.3)}.
    \begin{itemize}
        \item Note that $\Omega_{C/A}$ is generated by $d_{C/A}\bar{b}$ (where $\bar{b}=b+I\in C,\ b\in B$). Further, we have $v(d_{B/A}b\otimes1)=d_{C/A}\psi(b)=d_{C/A}\bar{b}$. Thus, $v$ is surjective.
        \item For any $b\in I$, we have $\bar{b}=0$ in $C$, so $v\delta(b+I^2)=v(d_{B/A}b\otimes1)=d_{C/A}\bar{b}=0.$
        \item For any $T\in\MOD_C$, we again rewrite the sequence $\Hom_C(-,T)$:
        \begin{center}
            \begin{tikzcd}
                \Hom_C(I/I^2,T)\dar["\simeq"] &\Hom_C(\Omega_{B/A}\otimes_BC,T)\lar["\delta"']\dar["\simeq"] &\Hom_C(\Omega_{C/A},T)\lar["-\circ v"']\dar[equal]\\ \Hom_B(I,T) &\Der_A(B,T)\lar["(\cdot)|_I"] &\Der_A(B/I,T)\lar["-\circ\psi"'],\\
                D|_I &D\lar[mapsto].
            \end{tikzcd}
        \end{center}
        Again, we have only to show that  $\ker(\cdot)|_I\subset\Im(-\circ\psi)$. Indeed, if $D\in\Der_A(B,T)$ with $D|_I=0$, this clearly lifts to an $A$-deriviation $B/I\to T$.
    \end{itemize}
\end{proof}

\begin{corollary}{8.5}
    Let $B\in\fg\ALG_A$ or $B=S^{-1}R$ for some $R\in\fg\ALG_A$, then $\Omega_{B/A}\in\fg\MOD_B$.
\end{corollary}
\begin{proof}
    In the first case, write $B=A[x_1,\cdots,x_n]/I$. From \textit{(8.4)}, the map $\Omega_{A[x_1,\cdots,x_n]/A}\otimes_{A[x_1,\cdots,x_n]}B\xrightarrow{v}\Omega_{B/A}$ is surjective. Further, $\{dx_i\otimes1\}_{i=1}^n$ is a $B$-basis of $\Omega_{A[x_1,\cdots,x_n]/A}\otimes_{A[x_1,\cdots,x_n]}B$, which implies that $\Omega_{B/A}$ is finitely generated by $\{d_{B/A}\bar{x}_i\}_{i=1}^n$ over $B$. Now, for the second case, we see that
    \begin{align*}
        \Omega_{B/A}=\Omega_{S^{-1}R/A}\underset{\textit{(8.2)}}{\simeq}S^{-1}\Omega_{R/A}=S^{-1}R\spn\{d_{R/A}\bar{x}_i/1\}_{i=1}^n\in\fg\MOD_B.
    \end{align*}
\end{proof}

\begin{lemma}{(8.B)}
    Let $k\subset K\subset L$ be finitely generated field extensions. We have the following four cases:
    \begin{itemize}
        \item[(i)] If $L/K$ is purely transcendental, then $\rank_L\Omega_{L/k}=\rank_K\Omega_{K/k}+\trdeg_KL$.
        \item[(ii)] If $L/K$ is algebraic and separable, then $\rank_L\Omega_{L/k}=\rank_K\Omega_{K/k}$.
        \item[(iii)] If $L=K(t)$ is algebraic but purely inseparable over $K$, then $\rank_L\Omega_{L/k}\geq\rank_K\Omega_{K/k}$.
        \item[(iv)] If, in general, $L/K$ is algebraic and purely inseparable, then $\rank_L\Omega_{L/k}\geq\rank_K\Omega_{K/k}$.
    \end{itemize}
\end{lemma}
\begin{proof}
    Note that $L/k$ is finitely generated, so $\rank_L\Omega_{L/k}<\infty$ by \textit{(8.5)}. Put $r\coloneqq\trdeg_KL$, which is finite since $L/K$ is finitely generated. Further, we may assume $K\subsetneq L$, otherwise the result is clear.\\
    
    (i) $L/K$ is purely transcendental, say $L=K(t_1,\cdots,t_r)$. Put $B\coloneqq K[x_1,\cdots,x_r]$, then $L\simeq\Frac B$ (over $K$). In this case, the exact sequence in \textit{(8.3)} becomes
    \begin{align*}
        0\longrightarrow\Omega_{K/k}\otimes_KB\xrightarrow{~~v~~}\Omega_{B/k}\xrightarrow{~~u~~}\Omega_{B/K}\longrightarrow0,
    \end{align*}
    and it splits since $v$ has a left inverse given by:
    \begin{align*}
        w:\Omega_{B/k}\longrightarrow\Omega_{K/k}\otimes_KB,\quad d_{B/k}f\longmapsto\sum_{\alpha=(\alpha_1,\cdots,\alpha_r)}d_{K/k}f_\alpha\otimes x_1^{\alpha_1}\cdots x_r^{\alpha_r}.
    \end{align*}
    Now, we get
    \begin{align*}
        \Omega_{B/k}\simeq(\Omega_{K/k}\otimes_KB)\oplus\Omega_{B/K}\underset{\textit{(8.2.1)}}{=}(\Omega_{K/k}\otimes_KB)\oplus Bdx_1\oplus\cdots\oplus Bdx_r,\\
        \Omega_{L/k}\underset{\textit{(8.2)}}{\simeq}\Frac\Omega_{B/k}\simeq(\Omega_{K/k}\otimes_KL)\oplus Ldt_1\oplus\cdots\oplus Ldt_r,
    \end{align*}
    whence $\rank_L\Omega_{L/k}=\rank_K\Omega_{K/k}+r$.\\
    
    (ii) If $L/K$ is algebraic and separable, since $L/K$ is finitely generated, it is also finite separable. By \textit{Primitive Element Theorem}, $\exists\ t\in L$ such that $L=K(t)$. Let $f\in K[x]$ be the minimal polynomial of $t$ of degree $n\geq2$, then $\{t^k\}_{k=0}^{n-1}$ is a $K$-basis for $L$. Now, to see the ranks are equal, it suffices to show that $\Omega_{K/k}\otimes_KL\xrightarrow[\MOD_L]{v}\Omega_{L/K}$ defined in \textit{(8.3)} is an isomorphism.
    \begin{itemize}
        \item \textit{Claim.} $\Omega_{L/K}=0$. Given any $d\in\Der_K(L,\Omega_{L/K})$. For any $\alpha\in L$, since $L/K$ is finite separable, we have $g(\alpha)=0$ but $g'(\alpha)\neq0$ for some polynomial $g\in K[x]$. It follows that $0=df(\alpha)=f'(\alpha)d\alpha$, whence $d\alpha=0$. But $\Omega_{L/K}$ is generated by the set $\{d\alpha:\ \alpha\in L\}$, thus, $\Omega_{L/K}=0$.
        \item Now, apply $\Omega_{L/K}=0$ to the exact sequence in \textit{(8.3)}, then $v$ is surjective. Finally, to show that $v$ is injective, it is enough to show that the map
        \begin{center}
            \begin{tikzcd}
                \Hom_L(\Omega_{K/k}\otimes_KL,T)\dar["\simeq"] &\Hom_L(\Omega_{L/k},T)\lar["-\circ v"']\dar[equal]\\ \Der_k(K,T) &\Der_k(L,T)\lar[]
            \end{tikzcd}
        \end{center}
        is surjective for any $T\in\MOD_L$. Now, for any $D\in\Der_k(L,T)$, we define
        \begin{align*}
            E:L\longrightarrow T,\quad\ \sum_{k=0}^{n-1}a_kt^k\longmapsto\ \sum_{k=0}^{n-1}t^kDa_k.
        \end{align*}
        It is easy to show that $E\in\Der_k(L,T)$ and that $E\mapsto D$.
    \end{itemize}
    
    (iii) Suppose $L=K(t)$ is algebraic but purely inseparable over $K$. In this case, we may assume that $t\notin K,\ \Char k=p>0$, and $e\geq1$ is the minimal for which $a=t^{p^e}\in K$. I claim that
    \begin{align*}
        \rank_L\Omega_{L/k}=\left\{\begin{array}{lc}
        \rank_K\Omega_{K/k}+1, &\text{if}\ d_{K/k}a=0; \\ 
        \rank_K\Omega_{K/k}, &\text{if}\ d_{K/k}a\neq0.\\ 
        \end{array}\right.
    \end{align*}
    Note that $f\coloneqq x^{p^e}-a\in K[x]$ is the minimal polynomial of $t$, and that $L\simeq K[x]/\<f\>$. Now, 
    \begin{align*}
        \Omega_{L/K} &\underset{\textit{(8.4)}}{\simeq}(\Omega_{K[x]/k}\otimes_{K[x]}L)/\Im\delta \underset{\text{(i)}}{\simeq}[(\Omega_{K/k}\otimes_KK[x]\oplus K[x]dx)\otimes_{K[x]}L]/\Im\delta\\  &\simeq[(\Omega_{K/k}\otimes_KL)\oplus Ldx]/\Im\delta
    \end{align*}
    as $L$-modules, where $\delta\bar{f}=d_{K/k}f\otimes1=-d_{K/k}a\otimes1$. If $d_{K/k}a=0$, then clearly, $\rank_L\Omega_{L/k}=\rank_K\Omega_{K/k}+1$; otherwise, $\delta\bar{f}$ generates an $L$-subspace of $\Omega_{L/k}$ of rank $1$, so $\rank_L\Omega_{L/k}=\rank_K\Omega_{K/k}+1-1=\rank_K\Omega_{K/k}$.\\
    
    (iv) In general, if $L/K$ is algebraic and purely inseparable, say $L=K(t_1,\cdots,t_m)$ with $t_i\notin K$. We have a chain of purely inseparable field extensions:
    \begin{align*}
        K\subset K(t_1)\subset K(t_1,t_2)\subset\cdots\subset K(t_1,\cdots,t_m)=L.
    \end{align*}
    Therefore, we have the inequality $\rank_L\Omega_{L/k}\geq\rank_K\Omega_{K/k}$ by (iii).
\end{proof}

\begin{proposition}{8.C}
    Let $k\subset K\subset L$ be finitely generated field extensions. Then
    \begin{align*}
        \rank_L\Omega_{L/k}\geq\rank_K\Omega_{K/k}+\trdeg_KL,
    \end{align*}
    and the equality holds if $L$ is separably generated over $K$.
\end{proposition}
\begin{proof}
    First, we write $L=K(t_1,\cdots,t_n)$ such that $\{t_i\}_{i=1}^r$ is a transcendence basis for $L/K$. Write $E=K(t_1,\cdots,t_r)$ so that $E/K$ is purely transcendental and $L/E$ is algebraic. From \textit{(6.Bi)}, we have $\rank_E\Omega_{E/k}=\rank_K\Omega_{K/k}+r$, so we may now reduce to the case where $L/K$ is algebraic.
    
    Let $K_s$ denote the set of all separable elements of $L$ over $K$. Since $L/K$ is algebraic, now $K_s/K$ is finite separable, and $L/K_s$ is finitely generated and purely inseparable. It follows from \textit{(8.Bii)} that $\rank_{K_s}\Omega_{K_s/k}=\rank_K\Omega_{K/k}$, so we reduce to the case (iv) of \textit{(6.B)}. Therefore, we obtain the inequality.
    
    Now, suppose $L/K$ is separably generated, there exists a transcendence basis $\{t_i\}_{i=1}^r$ of $K/k$ for which $E=K(t_1,\cdots,t_r)/K$ is purely transcendental and $L/E$ is separable algebraic. Therefore, 
    \begin{align*}
        \rank_L\Omega_{L/k}\underset{\textit{(6.Bii)}}{=}\rank_E\Omega_{E/k}\underset{\textit{(6.Bi)}}{=}\rank_K\Omega_{K/k}+r
    \end{align*}
    as required.
\end{proof}

\begin{theorem}{8.6}
    Let $K/k$ be a finitely generated field extension. Then,
    \begin{itemize}
        \item[(i)] $\rank_K\Omega_{K/k}\geq\trdeg_kK$.
        \item[(ii)] The equality holds $\Longleftrightarrow\ K/k$ is separably generated.
        \item[(iii)] Moreover, if $K/k$ is algebraic, then $\Omega_{K/k}=0\ \Longleftrightarrow\ K/k$ is separable.
    \end{itemize}
\end{theorem}
\begin{proof}
    Now, (i) and the direction ($\Leftarrow$) of (ii) are just the special case of \textit{(8.C)}. Further, we have shown the direction ($\Leftarrow$) of (iii) in \textit{(8.Bii)}. It remains to show the direction ($\Rightarrow$) of both (ii) and (iii).
    
    (iii) Suppose $\Omega_{K/k}=0$, i.e. $\rank_k\Omega_{K/k}=0$, so $\trdeg_kK=0$. Now, for any intermediate field $k\subset E\subset K$, we have $\rank_k\Omega_{E/k}=0$. One can show that the cases (i), (iii) and (iv) in \textit{(6.B)} cannot happen, i.e. $K/k$ is separable and algebraic. (For example, the case (iii), since $\Omega_{E/k}=0$, we can only have $d_{E/k}a=0$, whence $\rank_K\Omega_{K/k}=1>0$, a contradiction).
    
    (ii) Now, assume that $\rank_K\Omega_{K/k}=\trdeg_kK=r\geq1$. Let $\{dx_i\}_{i=1}^r$ be a $K$-basis for $\Omega_{K/k}$, then $\{dx_i\otimes1\}_{i=1}^r$ forms a $K$-basis for $\Omega_{k(x_1,\cdots,x_r)/k}\otimes_{k(x_1,\cdots,x_r)}K$. By \textit{(8.3)}, this gives $\Omega_{K/k(x_1,\cdots,x_r)}=0$, whence $K/k(x_1,\cdots,x_r)$ is separable and algebraic by (iii). Further, as $\trdeg_kK=r$, the set $\{x_i\}_{i=1}^r$ must form a transcendence basis for $K/k$, and hence, a separating transcendence basis. Therefore, $K/k$ is separably generated.
\end{proof}

\begin{proposition}{8.7}
    Let $B$ be a local ring containing its residue field $k=B/\mm$. Then, the map\\ $\delta:\mm/\mm^2\xrightarrow[\MOD_k]{}\Omega_{B/k}\otimes_Bk$ defined in \textit{(8.4)} is an isomorphism.
\end{proposition}
\begin{proof}
    Note that $\Omega_{k/k}=0$, so \textit{(8.4)} implies that $\delta$ is surjective. We have only to show that $\delta$ is injective. 
    \begin{itemize}
        \item Recall that the dual functor $(\cdot)^\ast=\Hom_D(-,D)$ over a division ring $D$ is a contravariant exact functor, and thus, it suffices to show that the map
        \begin{align*}
            \delta':\Hom_k(\Omega_{B/k}\otimes_Bk,k)\xrightarrow[\Ab]{-\circ\delta}\Hom_k(\mm/\mm^2,k)
        \end{align*}
        is surjective. 
        \item We note that
        \begin{align*}
            \Hom_k(\Omega_{B/k}\otimes_Bk,k)\simeq\Hom_B(\Omega_{B/k},k)=\Der_k(B,k).
        \end{align*}
        It follows that if $d\in\Der_k(B,k)$, then $\delta'(d)$ is obtained by restricting to $\mm$, and that $d(\mm^2)=0$.
        \item We have a short exact sequence $0\to\mm\to B\to k\to0$ of $B$-modules, but the fact that our $k\subset B$ implies that such exact sequence splits, and hence, $B\simeq\mm\oplus k$.
        \item Now, given any $h\in\Hom_k(\mm/\mm^2,k)$, we define the $k$-deriviation
        \begin{align*}
            d:B\longrightarrow k,\quad b\coloneqq c+\lambda\longmapsto h(c+\mm^2)=h(\bar{c}).
        \end{align*}
        It is clear that $\delta'(d)=h$. Further, we see that
        \begin{align*}
            d(bb') &=d((c+\lambda)(c'+\lambda'))=d(cc'+c\lambda'+c'\lambda+\underbrace{\lambda\lambda'}_{\in k})=h(\underbrace{\overline{cc}'}_{\in\mm^2}+\lambda'\bar{c}+\lambda\bar{c}')\\ &=h(\bar{c}\bar{c}'+\bar{c}'\bar{c}+\lambda'\bar{c}+\lambda\bar{c}') \underset{k\text{-linear}}{=}(c'+\lambda')h(\bar{c})+(c+\lambda)h(\bar{c}')=b'db+bdb',
        \end{align*}
        for any $b=c+\lambda,b'=c'+\lambda'\in B$. Thus, $d\in\Der_k(B,k)$, and therefore, $\delta'$ is surjective.
    \end{itemize}
\end{proof}

\begin{theorem}{8.8}
    Let $B$ be a local ring containing its residue field $k=B/\mm$. Assume further that $k$ is perfect, and that $B=S^{-1}R$ for some $R\in\fg\ALG_k$. Then,
    \begin{align*}
        \Omega_{B/k}\in\MOD_B\ \text{is a free module with}\ \rank_B\Omega_{B/k}=\dim B\ \Longleftrightarrow\ B\ \text{is a regular local ring}.
    \end{align*}
\end{theorem}
\begin{proof}
    ($\Rightarrow$) Note that such $B$ is a Noetherian ring from our assumption. Suppose that $\Omega_{B/k}$ is a free $B$-module of rank $\dim B$. Then
    \begin{align*}
        \dim_k\mm/\mm^2\underset{\textit{(8.7)}}{=}\dim_k(\Omega_{B/k}\otimes_Bk)=\rank_k\Omega_{B/k}=\dim B.
    \end{align*}
    Therefore, $B$ is a regular local ring. In particular, this implies that $B$ is an integral domain.
    
    ($\Leftarrow$) Suppose that $B$ is a regular local ring with $\dim B=r$. Then, 
    \begin{align*}
        \dim_k(\Omega_{B/k}\otimes_Bk)\underset{\textit{(8.7)}}{=}\dim_k\mm/\mm^2\underset{B:\ \text{regular}}{=}\dim B=r.
    \end{align*}
    On the other hand, let $K=\Frac B$, then $\Omega_{B/k}\otimes_BK=\Frac\Omega_{B/k}\underset{\textit{(8.2)}}{=}\Omega_{K/k}$. Now that $k$ is perfect, it follows from \textit{(I, 4.8)} that $K/k$ is separably generated, and thus,
    \begin{align*}
        r=\dim B\underset{\textit{(I, 1.8)}}{=}\dim_K\Omega_{K/k}\underset{\textit{(8.6)}}{=}\trdeg_kK.
    \end{align*}
    Finally, our assumption of $B$ gives that $\Omega_{B/k}\in\fg\MOD_B$. We conclude that $\Omega_{B/k}$ is a free module of rank $r$ by using the following lemma \textit{(8.9)}.
\end{proof}

\begin{lemma}{8.9}
    Let $A$ be a Noetherian local domain, with residue field $k$ and $K=\Frac A$. Suppose that $M\in\fg\MOD_A$ satisfies $\dim_k(M\otimes_Ak)=\dim_K(M\otimes_AK)=r$, then $M$ is free of rank $r$.
\end{lemma}
\begin{proof}
    Since $\dim_k(M\otimes_Ak)=r$, the \textit{Nakayama Lemma} states that $M$ can be generated by $r$ elements, i.e. $\exists$ surjective map $A^{\oplus r}\xrightarrow[\MOD_A]{\phi}M$. Put $I\coloneqq\ker\phi$. Note that $K$ is a flat $A$-module (as it is a field), i.e. $-\otimes_AK$ is an exact functor, we have an exact sequence
    \begin{align*}
        0\longrightarrow I\otimes_AK\longrightarrow K^{\oplus r}\xrightarrow{\phi\otimes\mathds{1}}M\otimes_AK\longrightarrow0
    \end{align*}
    of $A$-modules. Further, we have $\dim_K(M\otimes_AK)=r$, i.e. $M\otimes_AK\simeq K^{\oplus r}$, so $S^{-1}I=I\otimes_AK=0$ where $S=A\setminus\{0\}$. Thus, $\exists\ s\in S$ for which $sI=0$. But $I$ is torsion-free (as it is a submodule of a free module), so $I=0$, and thus, $M\simeq A^{\oplus r}$ is free of rank $r$.
\end{proof}
\newpage
\section{Sheaves of Differentials}

Let $f:X\xrightarrow[\Sch]{}Y$ be a morphism. Note that the diagonal morphism $\Delta:X\xrightarrow[\Sch]{}X\times_YX$ is an immersion, i.e. $\Delta$ gives an isomorphism of $X$ onto its image $\Delta(X)$, which is a \textit{locally closed subscheme} of $X\times_YX$, i.e. a closed subscheme of an open subset $W\subset X\times_YX$.

\begin{definition}{(Sheaf of Relative Differentials)}
    \newline Let $\II$ be the ideal sheaf of $\Delta(X)$ corresponding to the closed immersion $\Delta(X)\xrightarrow[\Sch]{i}W$. We define the \textit{sheaf of relative differentials} of $X$ over $Y$ (via $f$) to be the sheaf $\Omega_{X/Y}\coloneqq\Delta^\ast(\II/\II^2)\in\MOD_{\OO_X}$.
\end{definition}

\begin{remark}{8.9.1}
    Recall that $\II=\ker(\OO_W\overset{i^\#}{\twoheadrightarrow}i_\ast\OO_{\Delta(X)})\in\qcoh(W)$ by \textit{(5.9)}, so $\II/\II^2\in\qcoh(W)$. Now, we factorize $\Delta$ into
    \begin{center}
        \begin{tikzcd}
         &W\arrow[dr,hook,"\text{open}"]&\\ X\arrow[ur,"\widetilde{\Delta}"]\arrow[rr,"\Delta"'] &\Comm{u} &X\times_YX
        \end{tikzcd}
    \end{center}
    Now, we have
    \begin{align*}
        \Omega_{X/Y}=\Delta^\ast(\II/\II^2)=\widetilde{\Delta}^\ast(\II/\II^2)\underset{\textit{(5.8)}}{\in}\qcoh(X).
    \end{align*}
    Further, if $Y$ is a Noetherian scheme and $f$ is of finite type, then $X\times_YX$ is also Noetherian, whence $\Omega_{X/Y}\in\coh(X)$.
\end{remark}

\begin{remark}{8.9.2}
    If $U=\Spec A$ is an open affine of $Y$, we choose an open affine $V=\Spec B$ of $X$ for which $f(V)\subset U$, then we have $\Omega_{V/U}\simeq(\Omega_{B/A})\widetilde{}$. This shows that we may define $\Omega_{X/Y}$ by covering $X$ and $Y$ with open affines $V$ and $U$, respectively, and gluing the corresponding sheaves $(\Omega_{B/A})\widetilde{}$. Now, the deriviation $d\in\Der_A(B,\Omega_{B/A})$ patch together to give a map $d:\OO_X\xrightarrow[\Ab_X]{}\Omega_{X/Y}$, which is a deriviation of the local rings at each point.
\end{remark}
\begin{proof}
    Note that $V\times_UV\simeq\Spec(B\otimes_AB)$ is an open affine subset of $X\times_YX$, and that $\Delta:X\xrightarrow[\Top]{\sim}\Delta(X)$ is a homoemorphism onto to its image. Further, $f|_V$ is separated by \textit{(4.1)}, so we have the diagram:
    \begin{center}
        \begin{tikzcd}
            B &V\Comm{dr}\rar["\Delta"]\dar["\Delta_{V/U}","\text{closed}"'] &\Delta(V)\dar["\simeq"]\\ B\otimes_AB\uar[two heads,"\phi"] &V\times_UV &\Delta(X)\cap(V\times_UV)\lar[hook',"j"',"\text{closed}"]
        \end{tikzcd}
    \end{center}
    From \textit{(5.10)}, the closed subscheme $\Delta(V)\subset V\times_UV$ is uniquely determined by an ideal $I\triangleleft B\otimes_AB$. Moreover, $\II_{\Delta(V)}=\ker j^\#=(\ker\phi)\widetilde{}$ is the ideal sheaf of $\Delta(V)$ in $V\times_UV$, it follows from \textit{(5.9)} that its defining ideal must be $I=\ker\phi$, i.e. $\II_{\Delta(V)}=\widetilde{I}$. Therefore, 
    \begin{align*}
        \Omega_{V/U}=\Delta^\ast_{V/U}(\II_{\Delta(V)}/\II_{\Delta(V)}^2)=\Delta_{V/U}^\ast(I/I^2)\widetilde{}\underset{\textit{(5.2)}}{\simeq}(I/I^2\otimes_{B\otimes_AB}B)\widetilde{}=\,_B(I/I^2)\widetilde{}\underset{\textit{(8.1)}}{=}(\Omega_{B/A})\widetilde{}.
    \end{align*}
\end{proof}

\begin{proposition}{8.10}
    Given a morphism $X\xrightarrow[\Sch]{f}Y$, and let $X'=X\times_YY'\xrightarrow[\Sch]{f'}Y'$ be the base extension of $f$ through $Y'\xrightarrow[\Sch]{g}Y$, i.e. the fibre square
    \begin{center}
        \begin{tikzcd}
            X\times_YY'\FibSq{dr}\rar["f'"]\dar["g'"'] & Y'\dar["g"]\\ X\rar["f"'] &Y
        \end{tikzcd}
    \end{center}
    Then, we have $\Omega_{X'/Y'}\simeq(g')^\ast\ \Omega_{X/Y}$.
\end{proposition}
\begin{proof}
    From \textit{(8.9.2)}, the question is local, so we may reduce to the affine case. Suppose $Y,\ Y',\ X,\ X'$ are the spectra of the rings $A,\ A',\ B,\ B'$, respectively. Then $B'\simeq B\otimes_AA'$. Thus, we have
    \begin{align*}
        \Omega_{X'/Y'}\underset{\textit{(8.9.2)}}{\simeq}(\Omega_{B'/A'})\widetilde{}\underset{\textit{(8.2)}}{\simeq}(\Omega_{B/A}\otimes_BB')\widetilde{}\underset{\textit{(5.2e)}}{\simeq}(g')^\ast(\Omega_{B/A})\widetilde{}\underset{\textit{(8.9.2)}}{\simeq}(g')^\ast\ \Omega_{X/Y}.
    \end{align*}
\end{proof}

\begin{proposition}{8.11}
    Given morphisms $X\xrightarrow[\Sch]{f}Y\xrightarrow[\Sch]{g}Z$. Then there is an exact sequence of $\OO_X$-modules:
    \begin{align*}
        f^\ast\Omega_{Y/Z}\longrightarrow\Omega_{X/Z}\longrightarrow\Omega_{X/Y}\longrightarrow0.
    \end{align*}
\end{proposition}
\begin{proof}
     From \textit{(8.9.2)}, the question is local, so we may reduce to the affine case. Suppose $Z,\ Y,\ X$ are the spectra of the rings $A,\ B,\ C$, respectively. By \textit{(8.3)}, we have the exact sequence $\Omega_{B/A}\otimes_BC\to\Omega_{C/A}\to\Omega_{C/B}\to0$ of $C$-modules. Taking the exact functor $M\mapsto\widetilde{M}$ gives
     \begin{center}
         \begin{tikzcd}
            f^\ast(\Omega_{B/A})\widetilde{}\rar[equal,"\textit{(5.2)}"']\dar[equal] &(\Omega_{B/A}\otimes_BC)\widetilde{}\rar[] &(\Omega_{C/A})\widetilde{}\rar[]\dar[equal] &(\Omega_{C/B})\widetilde{}\rar[]\dar[equal] &0\\
            f^\ast\Omega_{Y/Z}\arrow[rr] & &\Omega_{X/Z}\rar[] &\Omega_{X/Y}\rar[] &0
         \end{tikzcd}
     \end{center}
     an exact sequence of $\OO_X$-modules.
\end{proof}

\begin{proposition}{8.12}
    Given a morphism $X\xrightarrow[\Sch]{f}Y$, and let $Z\subset X$ be a closed subscheme with its ideal sheaf $\II$. Then there is an exact sequence of $\OO_Z$-modules:
    \begin{align*}
        \II/\II^2\xrightarrow{~~\delta~~}\Omega_{X/Y}\otimes\OO_Z\longrightarrow\Omega_{Z/Y}\longrightarrow0.
    \end{align*}
\end{proposition}
\begin{proof}
    From \textit{(8.9.2)}, the question is local, so we may reduce to the affine case. Suppose $Y,\ X$ are the spectra of the rings $A,\ B$, respectively. Now, the closed subscheme $Z\subset X$ is determined by an ideal $I\triangleleft B$, for which $Z\simeq\Spec B/I$ and $\II\simeq\widetilde{I}$ by \textit{(5.10)}. Now, \textit{(8.4)} gives the exact sequence $I/I^2\xrightarrow{\delta}\Omega_{B/A}\otimes_BC\to\Omega_{C/A}\to0$ of $C=B/I$-modules. Taking the exact functor $M\mapsto\widetilde{M}$ gives
    \begin{center}
        \begin{tikzcd}
            (I/I^2)\widetilde{}\rar[]\dar[equal]  &(\Omega_{B/A}\otimes_BC)\widetilde{}\rar[]\dar[equal] &(\Omega_{C/A})\widetilde{}\rar[]\dar[equal] &0\\
            \II/\II^2\rar[] &\Omega_{X/Z}\otimes_{\OO_X}\OO_Z\rar[] &\Omega_{Z/Y}\rar[] &0
        \end{tikzcd}
    \end{center}
    an exact sequence of $\OO_Z$-modules.
\end{proof}

\begin{example}{8.12.1}
    If $X=\AA_Y^n$ for any scheme $Y$, then $\Omega_{X/Y}\in\MOD_{\OO_X}$ is a free module of rank $n$, generated by global sections $\{dx_i\}_{i=1}^n$. 
\end{example}
\begin{proof}
    Consider the fibre square
    \begin{center}
        \begin{tikzcd}
            X=\AA_\mathbb{Z}^n\times_\mathbb{Z}Y\FibSq{dr}\rar["f"]\dar["g"'] &Y\dar[]\\ Z=\AA_\mathbb{Z}^n\rar[] &\Spec\mathbb{Z}=S
        \end{tikzcd}
    \end{center}
    We have 
    \begin{align*}
        \Omega_{X/Y}\underset{\textit{(8.10)}}{\simeq}g^\ast\Omega_{Z/S}\underset{\textit{(8.9.2)}}{\simeq}g^\ast(\Omega_{\mathbb{Z}[x_1,\cdots,x_n]/\mathbb{Z}})\widetilde{}.
    \end{align*}
    Further, $(\Omega_{\mathbb{Z}[x_1,\cdots,x_n]/\mathbb{Z}})\widetilde{}$ is a free $\OO_Z$-module generated by the global sections $\{dx_i\}_{i=1}^n$, whence $\Omega_{X/Y}$ is a free $\OO_X$-module generated by the global sections $\{g^\ast dx_i\}_{i=1}^n$.
\end{proof}

\begin{theorem}{8.13}
    Let $A\in\Ring,\ Y=\Spec A$, and $X=\PP_A^n$. Then there is an exact sequence of $\OO_X$-modules:
    \begin{align*}
        0\longrightarrow\Omega_{X/Y}\longrightarrow\OO_X(-1)^{\oplus(n+1)}\longrightarrow\OO_X\longrightarrow0.
    \end{align*}
\end{theorem}
\begin{proof}
    (I) Construct the exact sequence.
    \begin{itemize}
        \item Write $S=A[x_0,\cdots,x_n]$, then $X=\Proj S$. Let $E\coloneqq S(-1)^{\oplus(n+1)}\in \GR\MOD_S$ be the free graded module with an $S$-basis $\{e_i\}_{i=0}^n$ where $e_i=(0,\cdots,\underset{i\text{th}}{1},\cdots,0)\in S_0^{\oplus(n+1)}=E_1$. 
        \item For each $x_i\in S_1$, we have a natural graded homomorphism 
        \begin{align*}
            \alpha_i:S(-1)\longrightarrow S,\quad1\in S(-1)_1\longmapsto x_i\in S_1.
        \end{align*}
        Then $\alpha\coloneqq\alpha_0\oplus\cdots\oplus\alpha_n$ gives a graded homomorphism $E\to S$. Let $M\coloneqq\ker\alpha$, then we have the exact sequence $0\to M\to E\to S$ of graded $S$-modules. By taking the functor $N\mapsto\widetilde{N}$, it induces an exact sequence of $\OO_X$-modules:
        \begin{align*}
            0\longrightarrow\widetilde{M}\longrightarrow\widetilde{E}=\OO_X(-1)^{\oplus(n+1)}\longrightarrow\OO_X\longrightarrow0.
        \end{align*}
        \item Note that the map $\alpha$ is NOT surjective (as $E_0=S_{-1}^{\oplus(n+1)}=\emptyset$), but it is surjective in $E_{\geq1}\to S_{\geq1}$, so the corresponding sheaf map $\widetilde{E}\to\widetilde{S}$ is still surjective!
    \end{itemize}
    
    (II) \textit{Claim.} $\widetilde{M}\simeq\Omega_{X/Y}$ as $\OO_X$-modules.
    \begin{itemize}
        \item First, we localize at each $x_i$. Then, $\alpha_{x_i}:E_{x_i}\to S_{x_i}$ is clearly a surjective graded  homomorphism of free $S_{x_i}$-modules. We see in the following that $M_{x_i}$ is also a free $S_{x_i}$-modules of rank $n$ with an $S_{x_i}$-basis $\{\frac{e_jx_i-e_ix_j}{x_i}:j\neq i\}$. Obviously, we have $\alpha_{x_i}(\frac{e_jx_i-e_ix_j}{x_i})=\frac{x_jx_i-x_ix_j}{x_i}=0$ for every $j\neq i$. On the other hand, suppose $s\in M_{x_i}=\ker\alpha_{x_i}$, we write $s\coloneqq\sum_{j\neq i}\frac{a_je_j}{x_i^k}$ ($a_j\in S$), which sends via $\alpha_{x_i}$ to $\sum_{j\neq i}\frac{a_je_j}{x_i^k}=0$. So, $\exists\ l>0$ for which $x_i^l\sum_{j\neq i}a_jx_j=0$. Now, %since
        %\begin{align*}
            %x_i^l\cdot\bigg(x_i^{k}\sum_{j\neq i}a_je_j-x_i^k\sum_{j\neq i}a_j(e_j-e_ix_j)\bigg)=x_i^l\cdot\bigg(e_ix_i^{k}\sum_{j\neq i}a_jx_j\bigg)=0,
        %\end{align*}
        we obtain
        \begin{align*}
            s=\frac{\sum_{j\neq i}a_je_jx_i}{x_i^{k+1}}-\underbrace{\frac{e_i\sum_{j\neq i}a_jx_j}{x_i^{k+1}}}_{=0}=\sum_{j\neq i}\frac{a_j}{x_i^k}\bigg(\frac{e_jx_i-e_ix_j}{x_i}\bigg)\in S_{x_i}\spn\bigg\{\frac{e_jx_i-e_ix_j}{x_i}:j\neq i\bigg\}.
        \end{align*}
        Moreover, it is easy to check that they are $S_{x_i}$-linearly independent, whence they form a basis over $S_{x_i}$.
        \item Let $U_i\coloneqq D_+(x_i)$ for each $i$. Now, it follows from the previous discussion that $\widetilde{M}|_{U_i}\simeq(M_{(x_i)})\widetilde{}$ is a free $\OO_X|_{U_i}$-module generated by the sections $\{\frac{e_jx_i-e_ix_j}{x_i^2}:j\neq i\}$ on $U_i$. (Here we need the additional factor $1/x_i$ since $x_i\in S_1\subset E_2$, and $x_i^2\in S_1\cdot E_2\subset E_3$). 
        \item For each $i$, we define a map $\phi_i:\Omega_{X/Y}|_{U_i}\to\widetilde{M}|_{U_i}$ as follows. Note that $U_i\simeq\Spec S_{(x_i)}=\Spec A[\frac{x_0}{x_i},\cdots,\frac{x_n}{x_i}]$, so $\Omega_{X/Y}|_{U_i}\simeq(\Omega_{S_{(x_i)/A}})\widetilde{}$ is a free $\OO_X|_{U_i}$-module generated by the global sections $\{d(x_j/x_i):j\neq i\}$ on $U_i$. Now, we define
        \begin{align*}
            \beta_i:\Omega_{S_{(x_i)}/A}\longrightarrow M_{(x_i)},\quad d\Big(\frac{x_j}{x_i}\Big)\longmapsto\frac{e_jx_i-e_ix_j}{x_i^2},
        \end{align*}
        which is thus an $S_{(x_i)}$-module isomorphism. From \textit{[5.3]}, $\beta_i$ corresponds to an isomorphism $\phi_i:\Omega_{X/Y}|_{U_i}\to\widetilde{M}|_{U_i}$ of free $\OO_X|_{U_i}$-modules.
        \item \textit{Claim.} The isomorphisms $\phi_i$ glue to give an isomorphism $\Omega_{X/Y}\xrightarrow[\MOD_{\OO_X}]{\phi}\widetilde{M}$.
        
        We look at the intersection $U_{ij}=U_i\cap U_j\simeq\Spec S_{(x_ix_j)}$, we have $\frac{x_k}{x_i}=\frac{x_k}{x_j}\frac{x_j}{x_i}$ for any $k$, which gives 
        \begin{align*}
            d\Big(\frac{x_k}{x_i}\Big)=\frac{x_k}{x_j}d\Big(\frac{x_j}{x_i}\Big)+\frac{x_j}{x_i}d\Big(\frac{x_k}{x_j}\Big)\qquad\text{in}\ \Omega_{S_{(x_ix_j)}/A}.
        \end{align*}
        (Here, $d=d_{S_{x_ix_j}/A}$). Now, we have
        \begin{align*}
            \phi_i|_{ij}\bigg(d\Big(\frac{x_k}{x_i}\Big)-\frac{x_k}{x_j}d\Big(\frac{x_j}{x_i}\Big)\bigg) &=\frac{e_kx_i-e_ix_k}{x_i^2}-\frac{x_k}{x_j}\frac{e_jx_i-e_ix_j}{x_i^2}=\frac{x_j}{x_i}\frac{e_kx_j-e_jx_k}{x_j^2}=\phi_j|_{ij}\bigg(\frac{x_j}{x_i}d\Big(\frac{x_k}{x_j}\Big)\bigg)
        \end{align*}
        for every $k$. Therefore, the ismorphisms $\{\phi_i\}$ glue on $\{U_i\}$ and obtain such isomorphism $\phi$.
    \end{itemize}
\end{proof}

\newpage
\section{Nonsingular Varieties}
\begin{definition}{(Nonsingular Varieties)}
    \begin{itemize}
        \item \textbf{(Abstract Varieties).} From now on, our \textit{(abstract) variety} means an integral separated scheme of finite type over an algebraically closed field $k$.
        \item \textbf{(Nonsingularity).} An (abstract) variety $X$ over $k$ is \textit{nonsigular} if all its local rings $\OO_{X,P}$ are regular. 
    \end{itemize}
\end{definition}

\begin{theorem}{8.14}
    Any localization of a regular local ring at a prime ideal is again a regular local ring.
\end{theorem}

\begin{theorem}{8.15}
    Let $X$ be an irreducible separated scheme of finite type over $k$. Then,
    \begin{align*}
        \Omega_{X/k}\in\MOD_{\OO_X}\ \text{is locally free of rank}\ n=\dim X\ \Longleftrightarrow\ X\ \text{is a nonsingular variety over}\ k.         
    \end{align*}
\end{theorem}
\begin{proof}
    Note that the question is local by covering $X$ with open affine subsets, so we may assume $X$ is affine, say $X=\Spec A$. 
    \begin{itemize}
        \item If $P\in X$ is a closed point, we have $\dim\OO_{X,P}\underset{\textit{[3.20a]}}{=}\dim X=n,\ \OO_{X,P}/\mm_P\simeq k$. Write $B=\OO_{X,P}=A_\mm$, where $\mm=\mm_P$ corresponds to the point $P$. Thus,
        \begin{align*}
            (\Omega_{X/k})_P\underset{\textit{(8.9.2)}}{\simeq}(\Omega_{A/k})\widetilde{}_P\underset{\textit{(5.1)}}{\simeq}(\Omega_{A/k})_\mm\underset{\textit{(8.2)}}{\simeq}\Omega_{B/k}.
        \end{align*}
        \item Since $X$ is of finite type over $k$, we have $A\in\fg\ALG_k$. Further, $k$ is perfect as it is algebraically closed, so we see from \textit{(8.8)} that
        \begin{align*}
            \Omega_{B/k}\in\MOD_B\ \text{is free of rank}\ n\ \Longleftrightarrow\ B\ \text{is regular}.
        \end{align*}
        \item Thus, we have
        \begin{align*}
            \Omega_{X/k}\ \text{is locally free of rank}\ n &\underset{\textit{[5.7b]}}{\Longleftrightarrow}\ (\Omega_{X/k})_P\in\MOD_{\OO_{X,P}}\ \text{is free of rank $n$ for all}\ P\in X\\ &\Longleftrightarrow\ (\Omega_{A_\mm/k})\in\MOD_{A_\mm}\ \text{is free of rank $n$ for all closed point}\ \mm\triangleleft X\\ &\underset{\textit{(8.8)}}{\Longleftrightarrow}\ A_\mm\ \text{is regular for all closed point}\ \mm\triangleleft A\\
            &\underset{\textit{(8.14)}}{\Longleftrightarrow}\ A_\pp\ \text{is regular for all point}\ \pp\triangleleft A\\
            &\Longleftrightarrow\ X=\Spec A\ \text{is (integral and) nonsingular}.
        \end{align*}
        Here, for the second arrow, we note that is $M_\mm$ is a free $A_\mm$-module (where $\mm\triangleleft A$ maximal), then $M_\pp=(M_\mm)_\pp=M_\mm\otimes_AA_\pp$ is a free $A_\pp$-module for any prime $\pp\triangleleft A$ with $\pp\subset\mm$.
    \end{itemize}
\end{proof}

\begin{corollary}{8.16}
    If $X\in\Var_k$, then there is an open dense subset $U\subset X$ which is nonsingular. Note, this gives a new proof of \textit{(I, 5.3)}. 
\end{corollary}
\begin{proof}
    Suppose $\dim X=n$ and write $K=K(X)$. Then $\trdeg_kK=n$ by \textit{[3.20b}, where $K/k$ is a finitely generated field extension, whence $K/k$ is separably generated by \textit{(I, 4.8)}. Now, we have $\rank_K\Omega_{K/k}\underset{\textit{(8.6)}}{=}\trdeg_kK=n$. On the other hand, we note that $(\Omega_{X/k})_\eta=\Omega_{K/k}$ where $\eta\in X$ is the generic point. Clearly, the field $K$ is a regular local ring, so \textit{(8.8)} shows that $(\Omega_{X/k})_\eta=\Omega_{K/k}$ is a free $\OO_{X,\eta}=K$-module of rank $n$. By \textit{[5.7a]}, $\exists\ U\in\NN_\eta(X)$ on which $\Omega_{X/k}|_U\in\MOD_{\OO_X|_U}$ is free of rank $n$. Thus, $U$ is a nonsingular variety over $k$ by \textit{(8.15)}. Clearly, such $U$ is a dense subset as it is a nonempty open set in the irreducible topological space $X$. Done!    
\end{proof}

\begin{theorem}{8.17}
    Let $X$ be a nonsingular variety over $k$, and let $Y\subset X$ be an irreducible closed subscheme with its ideal sheaf $\II$. Then, $Y$ is nonsingular if and only if 
    \begin{itemize}
        \item[(1)] $\Omega_{Y/k}\in\MOD_{\OO_Y}$ is locally free.
        \item[(2)] The sequence of \textit{(8.12)} is exact on the left, i.e.
        \begin{align*}
            0\longrightarrow\II/\II^2\xrightarrow{~~\delta~~}\Omega_{X/k}\otimes\OO_Y\longrightarrow\Omega_{Y/k}\longrightarrow0.
        \end{align*}
        Furthermore, in this case, $\II$ is locally generated by $r=\codim(Y,X)$ elements, and $\II/\II^2\in\MOD_{\OO_Y}$ is a locally free sheaf of rank $r$.
    \end{itemize}
\end{theorem}
\begin{proof}
    ($\Leftarrow$) Suppose both (1) and (2) are satisfied.
    \begin{itemize}
        \item From (1), $\Omega_{Y/k}$ is locally free. By \textit{(8.15)}, it suffices to show that $\rank_{\OO_Y}\Omega_{Y/k}=\dim Y$. Say $q=\rank_{\OO_Y}\Omega_{Y/k}$.
        \item As $X$ is nonsingular, we see also from \textit{(8.15)} that $\Omega_{X/k}\in\MOD_{\OO_X}$ is locally free of rank $n=\dim X$. By (2), it follows that $\rank_{\OO_Y}(\II/\II^2)=n-q$. By \textit{Nakayama Lemma}, $\II$ can be locally generated by $n-q$ elements. Thus, $\dim Y\geq n-(n-q)=q$ by \textit{[I, 1.9]}.
        \item For any closed point $P\in Y$, write $B\coloneqq\OO_{Y,P}$, then
        \begin{align*}
            \dim Y\underset{\text{(I,5.2)}}{\leq}\dim_k(\mm_P/\mm_P^2)\underset{\text{(8.7)}}{=}\rank_k(\Omega_{B/k}\otimes_Bk)=\rank_B\Omega_{Y/k}=\rank_{\OO_Y}\Omega_{Y/k}=q.
        \end{align*}
        Therefore, $Y$ is nonsingular, and that $\rank_{\OO_Y}(\II/\II^2)=n-q=\dim X-\dim Y=\codim(Y,X)$.
    \end{itemize}
    
    ($\Rightarrow$) Conversely, suppose $Y$ is a nonsingular variety.
    \begin{itemize}
        \item Again, (1) follows from \textit{(8.15)}, and that $\rank_{\OO_Y}\Omega_{Y/k}=\dim Y$, say $q$. Further, we have the exact sequence $\II/\II^2\xrightarrow{~~\delta~~}\Omega_{X/k}\otimes\OO_Y\xrightarrow{~~\phi~~}\Omega_{Y/k}\longrightarrow0$ of $\OO_Y$-modules from \textit{(8.12)}. So, it suffices to show that $\delta$ is injective.
        \item Let $y\in Y$ be any closed point. We have the short exact sequence
        \begin{align*}
            0\longrightarrow\Im\delta=\ker\phi\longrightarrow\Omega_{X/k}\otimes\OO_Y\xrightarrow{~~\delta~~}\Omega_{Y/k}\longrightarrow0
        \end{align*}
        of $\OO_Y$-modules (as $\phi$ is surjective). Note that $\Omega_{Y/k}\in\coh(Y)$ by \textit{(8.9.1)}, so $\ker\phi\in\coh(Y)$ by \textit{(5.7)} Thus, $(\Im\delta)_y=(\ker\phi)_y$ is a free $\OO_{Y,y}$-module of rank $r=n-q$. From \textit{[5.7a]}, $\exists\ U\in\NN_y(Y)$ on which $(\ker\phi)|_U\in\MOD_{\OO_Y|_U}$ is free of rank $r$. Now, we may choose some sections $x_i\in\Gamma(U,\II)$ ($1\leq i\leq r$) such that $\{dx_i\}_{i=1}^r$ forms an $\OO_Y|_U$-basis of $(\ker\phi)|_U$.
        \item Now, consider the ideal sheaf $\II'$ generated by the sections $\{x_i\}_{i=1}^r$ over $\OO_X$, and let $Y'\subset X$ be the corresponding closed subscheme. Then, we have $\II'\subset\II$, and thus,
        \begin{align*}
            Y'=\Supp(\OO_X/\II')=\{y\in X:\II_y'\subsetneq\OO_{X,y}\}\supset\{y\in X:\II_y\subsetneq\OO_{X,y}\}=\Supp(\OO_X/\II)=Y.
        \end{align*}
        \item From this construction, $\{dx_i\}_{i=1}^r$ generates a free $\OO_{Y'}$-submodule of $\Omega_{X/k}\otimes\OO_{Y'}$ of rank $r$ in the neighborhood $U\in\NN_y(Y')$. Apply the exact sequence of \textit{(8.12)} again to $Y'$, then the map $\delta:\II'/\II'^2\to\Omega_{X/k}\otimes\OO_{Y'}$ is injective since $(\Im\delta)_y=\Im(\delta_y)$ is free of rank $r$ on each closed point $y\in Y$. Further, this shows that $\Omega_{Y'/k}\in\MOD_{\OO_{Y'}}$ is locally free of rank $n-r=q$.
        \item Finally, applying these results to the proof of ($\Leftarrow$) yields that $Y'$ is nonsingular (thus, irreducible and integral) of dimension $\dim Y'=q$. So, $Y\subset Y'$ are integral schemes having the same dimension, so we must have $Y=Y'$ and $\II=\II'$. Therefore, $\delta$ is injective.
    \end{itemize}
\end{proof}

\begin{lemma}{8.D}
    Let $(A,\mm)$ be a regular local ring, and given a nonzero element $a\in\mm$. Then, $A/\<a\>$ is a regular local ring $\Longleftrightarrow\ a\notin\mm^2$.
\end{lemma}
\begin{proof}
    Suppose $A$ is regular of dimension $n\geq1$. Denote $V$ for the $k$-vector space
    \begin{align*}
        \frac{\mm/\<a\>}{(\mm/\<a\>)^2}\simeq\frac{A/\<a\>}{\mm/\<a\>}\otimes_{A/\<a\>}\mm/\<a\>\simeq k\otimes_{A/\<a\>}\mm/\<a\>.
    \end{align*}
    
    ($\Rightarrow$) Suppose $A/\<a\>$ is regular. Suppose on the contrary that $a\in\mm^2$, then we have $(\mm/\<a\>)^2=\mm^2/\<a\>$, whence $V\simeq\mm/\mm^2$. Thus, $\dim A/\<a\>=\dim_kV=\dim A=n$. However, if we choose a maximal chain
    \begin{align*}
        \pp_0/\<a\>\subsetneq\pp_1/\<a\>\subsetneq\cdots\subsetneq\pp_n/\<a\>=\mm/\<a\>
    \end{align*}
    of prime ideals of $A/\<a\>$, as $a\neq0$ and $A$ is integral (since it is regular), we obtain a chain $0\subsetneq\pp_0\subsetneq\pp_1\subsetneq\cdots\subsetneq\pp_n=\mm$ of prime ideals of $A$, whence $\dim A\geq n+1$, a contradiction.\\
    
    ($\Leftarrow$) \textit{Claim.} $A/\<a\>$ is regular of dimension $n-1$. Clearly, $A/\<a\>$ is also a local ring with maximal ideal $\mm/\<a\>$. If $a\notin\mm^2$, we may extend $\bar{a}\in\mm/\mm^2$ to a $k$-basis $\{\bar{a}_i\}_{i=1}^n$ of $\mm/\mm^2$ (where $a_i\in\mm$ and $\bar{a}=\bar{a}_1$). Then, $\{1\otimes(a_i+\<a\>)\}_{i=2}^n$ generates the $k$-vector space $V$, whence $\dim_kV\leq n-1$. On the other hand, consider a maximal chain $0=\pp_0\subsetneq\pp_1\subsetneq\cdots\subsetneq\pp_n=\mm$ of prime ideals of $A$. (Note. $A$ is regular, thus integral). Moreover, $A$ is Noetherian with $0\neq a\in\mm$, so we may modify the chain $\{\pp_i\}_{i=0}^n$ and assume that $a\in\pp_1$. Thus, 
    \begin{align*}
        \pp_1/\<a\>\subsetneq\pp_2/\<a\>\subsetneq\cdots\subsetneq\pp_n/\<a\>=\mm/\<a\>
    \end{align*}
    is a chain of prime ideal of $A/\<a\>$, whence $\dim A/\<a\>\geq n-1$. However, we already have $\dim A/\<a\>\leq\dim_kV$. Combining these all, we obtain $\dim A/\<a\>=\dim_kV=n-1$.
\end{proof}

\begin{theorem}{8.18. (Bertini Theorem)}
    \newline Let $X$ be a nonsingular closed subvariety of $\PP_k^n$. Then 
    \begin{itemize}
        \item[(i)] There exists a hyperplane $H\subset\PP_k^n$, not containing $X$, such that the scheme $H\cap X$ is regular at every point.
        \item[(ii)] The set of such hyperplanes forms an open dense subset of the complete linear system $|H|$, considered as a projective space.
    \end{itemize}
\end{theorem}
\begin{proof}
    (I) For each closed point $x\in X$, we consider the set
    \begin{align*}
        B_x\coloneqq\{\text{hyperplanes}\ H\ |\ H\supset X,\ \text{or}\ H\not\supset X\ \text{with}\ x\in H\cap X\ \text{non-regular}\}.
    \end{align*}
    Note that a hyperplane $H\subset\PP_k^n$ is determined by a nonzero global section $f\in S_1=\Gamma(\PP_k^n,\OO(1))=:V$. Fix an $f_0\in V$ for which $x_0\notin H_0\coloneqq V_+(f_0)$. Now, we define a map of $k$-linear spaces as follows:
    \begin{center}
        \begin{tikzcd}
            V\rar[] &S_{(f_0)}\rar[equal] &\OO_{\PP_k^n}(D_+(f_0))\Comm{dr}\rar["i^\#"]\dar[] &\OO_X(X\setminus H_0)\dar[] &\\
            & &\OO_{\PP_k^n,x}\rar["i_x^\#"] &\OO_{X,x}\rar[two heads] &\OO_{X,x}/\mm_x^2,\\
            f\rar[mapsto] &\frac{f}{f_0}\rar[mapsto] &\Big[\frac{f}{f_0}\Big]_x\rar[mapsto] &\Big[i^\#\Big(\frac{f}{f_0}\Big)\Big]_x=:\phi_x(f) &
        \end{tikzcd}
    \end{center}
    where $i$ is the closed immersion $X\xhookrightarrow{}\PP_k^n=Y$. Now, the scheme $H\cap X$ is given by $V_+(f)\cap X=V_+(i^\#(f/f_0))$. Thus, we see that
    \begin{itemize}
        \item First, note that 
        \begin{align*}
            x\in H\setminus H_0\ \Longleftrightarrow\ \Big[\frac{f}{f_0}\Big]_x\in\mm_{Y,x}\ \Longleftrightarrow\ \phi_x(f)\in\mm_x.
        \end{align*}
        \item Suppose $X$ is defined by the homogeneous prime ideal $I\triangleleft S=k[x_0,\cdots,x_n]$, i.e. $\OO_X\simeq(S/I)\widetilde{}$. Then 
        \begin{align*}
            H\supset X\ \text{(as a closed subscheme)}\  &\Longleftrightarrow\ f\in I\ \Longleftrightarrow\ \Big[i^\#\Big(\frac{f}{f_0}\Big)\Big]_x=\Big[\frac{f}{f_0}\Big]_x+\<\phi_x(f)\>\\
            &\Longleftrightarrow\ \phi_x(f)=0\ \text{in}\ \mm_x=\mm_{Y,x}/\<\phi_x(f)\>.
        \end{align*}
        \item If $H\not\supset X$, then $x\in H\cap X\ \Longleftrightarrow\ \phi_x(f)=[i^\#(f/f_0)]_x\in\mm_x$ is nonzero. In this case, 
        \begin{align*}
            x\in H\cap X\ \text{is nonregular}\ &\Longleftrightarrow\ \OO_{H\cap X,x}=\OO_{X,x}/\<\phi_x(f)\>\ \text{is not a regular ring}\\ &\underset{\text{(8.D)}}{\Longleftrightarrow}\ \phi_x(f)\notin\mm_x^2.
        \end{align*}
    \end{itemize}
     Therefore, we conclude that $H\in B_x\ \Longleftrightarrow\ f\in\ker(\phi_x)$. \\
     
     (II) Since $x\in X$ is a closed point and $k$ is algebraically closed, the maximal ideal $\mm_x$ is generated by $\{a_jx_i-a_jx_i:1\leq i,j\leq n\}$ where $a_i\in k$ not all zero. Now, it is an algebraic problem: writing $A=\OO_{X,x},\ \mm=\mm_x,\ M=A/\mm^2$, then
     $\{(a_jx_i-a_ix_j)/f_0+\mm^2\}_{i,j}$ generates $\mm/\mm^2\subset M$ over $k$. In particular, $A^\times=A\setminus\mm$ implies that the image of $\phi_x(f_0)=f_0/f_0+\mm^2$ generates the $k$-vector space $M/\mm M\simeq k$. By \textit{Nakayama (AM, 2.8)}, $M$ is finitely generated by $\{\phi_x(a_jx_i-a_jx_i)\}_{i,j}\cup\{\phi_x(f_0)\}$ over $k$, and therefore, $\phi_x$ is surjective. Now, the \textit{Isomorphism Theorem} gives $V/\ker\phi_x\simeq A/\mm^2$. If $\dim X=r$, recall that $V=k\spn\{x_i\}_{i=0}^n$, we obtain
     \begin{align*}
         \dim_k(\ker\phi_x) &=\dim_kV-\dim_k(V/\ker\phi_x)=(n+1)-\dim_k(A/\mm^2)\\ &=(n+1)-\Big[\dim_k(\mm/\mm^2)+\dim_k\frac{A/\mm^2}{\mm/\mm^2}\ \Big]=(n+1)-\Big[\underset{\text{nonsingular}}{\dim X}+\dim_kk\ \Big]\\ &=(n+1)-(r+1)=n-r.
     \end{align*}
     Now, we may identify $B_x=\PP(\ker\phi_x)$, i.e. a linear system of hyperplanes in $|H_0|$, which has dimension $\dim B_x=n-r-1$. (cf. \S II.7).
\end{proof}

\newpage
\section{Applications}

We define some invariants of nonsingular varieties.

\begin{definition*}{}
    Let $X$ be a nonsingular variety over $k$ of dimension $n$.
    \begin{itemize}
        \item \textbf{(Tangent Sheaf).} The \textit{tangent sheaf} of $X$ is % defined to be 
        $\TT_X\coloneqq\HHom_{\OO_X}(\Omega_{X/k},\OO_X)\in\MOD_{\OO_X}$, which is locally free of rank $n$.
        \item \textbf{(Canonical Sheaf).} The \textit{canonical sheaf} of $X$ is defined to be $\omega_X\coloneqq\bigwedge^n\Omega_{X/k}\in\MOD_{\OO_X}$, which is an invertible sheaf.
        \item \textbf{(Geometric Genus).} If, in addition, $X$ is projective, we define the \textit{geometric genus} of $X$ to be $p_g\coloneqq\dim_k\Gamma(X,\omega_X)\geq0$.
    \end{itemize}
\end{definition*}

\begin{theorem}{8.19}
    Let $X$ and $X'$ be two birationally equivalent nonsingular projective varieties over $k$. Then, $p_g(X)=p_g(X')$.
\end{theorem}
\begin{proof}
    Recall that for $X$ and $X'$ to be birationally equivalent, there are rational maps $X\to X'$ and $X'\to X$ which are inverses to each other. We consider the rational map $X\to X'$ represented by the morphism $V\xrightarrow[\Sch]{f}X'$, where $V\subset X$ is the largest defining open set. By \textit{(8.11)}, we have a map $f^\ast\Omega_{X'/k}\xrightarrow[\MOD_{\OO_V}]{}\Omega_{V/k}$. Note that
    \begin{itemize}
        \item $\dim X=\dim X'=n$.
        \item $X$ and $X'$ being nonsingular yields that $\Omega_{V/k}=(\Omega_{X/k})|_V,\ \Omega_{X'/k}$, and thus $f^\ast\Omega_{X'/k}$ are locally free sheaves of rank $n$.
        \item \textit{Check.} $\omega_X|_V=(\bigwedge^n\Omega_{X/k})|_V\underset{\text{(?)}}{\simeq}\bigwedge^n(\Omega_{X/k}|_V)\underset{\text{(8.10)}}{\simeq}\bigwedge^n\Omega_{V/k}=\omega_V$
        \item \textit{Check.} $f^\ast\omega_{X'}=f^\ast(\bigwedge^n\Omega_{X'/k})\underset{\text{(?)}}{\simeq}\bigwedge^n(f^\ast\Omega_{X'/k})$.
    \end{itemize}
    Thus, this induces a natural map $f^\ast\omega_{X'}\simeq\bigwedge^n(f^\ast\Omega_{X'/k})\xrightarrow[\MOD_{\OO_V}]{}\bigwedge^n\Omega_{V/k}=\omega_V$. By \textit{Adjoint Property}, we have a morphism $\omega_{X'}\xrightarrow[\MOD_{\OO_{X'}}]{\phi}f_\ast\omega_V$, which gives a map $\Gamma(X',\omega_{X'})\xrightarrow[\MOD_k]{\phi^\#}\Gamma(V,\omega_V)$ of $k$-vector spaces.\\
    
    (I) \textit{Claim.} $\phi^\#$ is injective.
    \begin{itemize}
        \item As $f$ is birational, it follows from \textit{(I, 4.5)} that $\exists\ U\subset V$ open such that $f(U)\subset X'$ is open and that $f:U\xrightarrow[\Sch]{\sim}f(U)$ is an isomorphism. Thus, an isomorphism $\omega_V|_U\xleftarrow[]{\sim}\omega_{X'}|_{f(U)}$ of sheaves.
        \item Now, suppose $s\in\ker\phi^\#$, we look at the following diagram:
        \begin{center}
            \begin{tikzcd}
                \Gamma(X',\omega_{X'})\Comm{dr}\rar[hook,"(\cdot)|_{f(U)}"]\dar["\phi^\#"] &\Gamma(f(U),\omega_{X'})\Comm{dr}\rar["\sim"]\dar["\phi_{f(U)}^\#"] &\Gamma(f(U),\omega_{X'}|_{f(U)})\dar["\sim"]\\
                \Gamma(V,\omega_V)\rar[hook,"(\cdot)|_U"] &\Gamma(U,\omega_V)\rar["\sim"] &\Gamma(U,\omega_V|_U)
            \end{tikzcd}
        \end{center}
        Note that $\omega_{X'}$ and $\omega_V$ are both invertible sheaves on an integral scheme, so the restriction maps $(\cdot)|_{f(U)}$ and $(\cdot)|_U$ are both injective. Further, $\phi_{f(U)}^\#$ is an isomorphism, so $s|_{f(U)}=0$. But $f(U)\subset X$ is an open dense subset (as $X$ is irreducible), whence $s=0$. Therefore, $\phi^\#$ is injective.
    \end{itemize}
    
    (II) \textit{Claim.} Let $X$ be nonsingular and $X'$ be projective, with a rational morphism $f:X\to X'$ defined on the open set $V\subset X$. Set $Z\coloneqq X\setminus V$, then $\codim(Z,X)\geq2$.
    \begin{itemize}
        \item Note that $V\neq\emptyset$, and $X$ is irreducible, so $\codim(Z,X)\geq1$.
        \item Given any point $P\in X$ with $\dim\OO_{X,P}=1$. As $X$ is nonsingular, $\OO_{X,P}$ is a regular (Noetherian local domain), whence a DVR by \textit{(I, 6.2)}.
        \item Note that $f|_U:U\xrightarrow[\Top]{\sim}f(U)\subset X'$ is a homeomorphism, which must map the generic point $\eta\in U\subset V$ to the generic point $\eta'\in X'$.
        \item Recall that $f(P)\in X'$ and an inclusion $k(f(P))\xhookrightarrow{}k(P)$ (induced by $\OO_{X',f(P)}\to\OO_{V,P}$) determine the morphism $\Spec k(P)\xrightarrow[\Sch]{\alpha}X'$ by \textit{[2.7]}. Now, $X'\xrightarrow[\Sch]{\pi'}\Spec k$ is projective, thus, proper by \textit{(4.9)}. It follows from the \textit{Valuative Criterion (4.7)} that $\exists$ unique morphism $\Spec\OO_{X,P}\xrightarrow[\Sch]{\theta}X'$ such that the whole diagram below commute:
        \begin{center}
            \begin{tikzcd}
                \Spec k(P)\rar["\alpha"]\dar[hook,"i"'] &X'\dar["\pi'"',"\text{proper}"]\\ \Spec\OO_{X,P}\rar["\beta"']\arrow[ur,dashrightarrow,"\theta"] &\Spec k
            \end{tikzcd}
        \end{center}
        \item \textbf{(?)} Further, note that $\Spec\OO_{X,P}$ is the intersection of all neighborhoods of $P$ in $X$, so we may choose a small $W\in\NN_P(X)$ and consider the morphism $\theta':\Spec\OO_{X,P}\xhookrightarrow[\Sch]{j}W\xrightarrow[\mathfrak{Rat}]{f}X'$. \textit{Check.} $\theta'$ is also compatible with the diagram, and hence, $\theta=\theta'$ by uniqueness. So, $f$ extends to a morphism in a neighborhood of $P$. But $V\subset X$ is the largest defining domain, whence $P\in V$.
        \item \textit{Conclusion.} We see that $Z$ contains no point with $\dim\OO_{X,P}=1$. Thus, by \textit{[3.20(c)]}, we have $\codim(Z,X)\geq2$.
    \end{itemize}
    
    (III) \textit{Claim.} The map $\Gamma(X,\omega_X)\xrightarrow[\MOD_k]{}\Gamma(V,\omega_V)$ is bijective, and thus, a $k$-linear isomorphism.
    \begin{itemize}
        \item The \textit{injectivity} part is easy since $\omega_V\simeq\omega_X|_V$ is an invertible $\OO_V$-module on the integral scheme $V\subset X$. It suffices to show that: for any $x\in X,\ \exists$ open affine $U\in\NN_x(X)$ such that $\Gamma(U,\omega_X)\xrightarrow[\MOD_k]{}\Gamma(U\cap V,\omega_V)$ is bijective.
        \item Indeed, as $\omega_X$ is invertible, we may choose a small neighborhood $U=\Spec A$, on which $\omega_X|_U\simeq\OO_X|_U$. Then, we must show that $\Gamma(U,\OO_X|_U)\xrightarrow[\MOD_k]{}\Gamma(U\cap V,\omega_X|_{U\cap V})$ is bijective. At this point, we may replace $(X,V)$ by $(U=\Spec A,U\cap V)$.
        \item Since $X=\Spec A$ is nonsingular, every local ring $A_\pp$ ($\pp\in\Spec A$) is a regular domain, thus, UFD and integrally closed. So, $A$ is also an integrally closed domain (AM, Ch3).
        \item Now, from (II), $\codim(Z,X)\geq2$, so every prime ideal $\pp\triangleleft A$ with $\Ht\pp=1$ must be contained in $V$ (which gives the following first map). It follows that
        \begin{align*}
            \Gamma(V,\OO_X)\xhookrightarrow[]{}\bigcap_{\Ht\pp=1}A_\pp\underset{\text{(6.3)}}{=}A\xhookrightarrow{}K(X)
        \end{align*}
        as subrings of $K(X)$. Therefore, $A=\Gamma(X,\OO_X)\hookrightarrow\Gamma(V,\OO_X)$ is bijective.\\
        
        (IV) Finally, combining the results above gives
        \begin{align*}
            p_g(X)=\dim_k\Gamma(X,\omega_X)\underset{\text{(III)}}{=}\dim_k\Gamma(V,\omega_V)\underset{\text{(I)}}{\geq}\dim_k\Gamma(X',\omega_{X'})=p_g(X').
        \end{align*}
        We obtain the converse ($\leq$) by symmetry, whence $p_g(X)=p_g(X')$.
    \end{itemize}
\end{proof}

\begin{definition*}{}
    Let $Y$ be a nonsingular subvariety of a nonsingular variety $X$ over $k$.
    \begin{itemize}
        \item \textbf{(Conormal Sheaf).} The \textit{conormal sheaf} of $Y$ in $X$ is $\II/\II^2\in\MOD_{\OO_Y}$ as in \textit{(8.17)}.
        \item \textbf{(Normal Sheaf).} The \textit{normal sheaf} of $Y$ in $X$ is the dual $\NN_{Y/X}\coloneqq(\II/\II^2)^\vee\in\MOD_{\OO_Y}$, which is locally free of rank $r=\codim(Y,X)$.
    \end{itemize}
\end{definition*}

\begin{remark*}{}
    Note that the dual functor $(\cdot)^\vee=\Hom_{\OO_Y}(-,\OO_Y)$ of locally free $\OO_Y$-modules is an exact (contravariant) functor. Applying this to \textit{(8.17)}, we have the following exact sequence:
    \begin{align*}
        0\longleftarrow\NN_{Y/X}\longleftarrow\TT_X\otimes_{\OO_X}\OO_Y\longleftarrow\TT_Y\longleftarrow0.
    \end{align*}
    Here, the middle term comes from the isomorphism
    \begin{align*}
        (\Omega_{X/Y}\otimes_{\OO_X}\OO_Y)^\vee=(i^\ast\Omega)^\vee\simeq i^\ast(\Omega_{X/Y}^\vee)=i^\ast\TT_X=\TT_X\otimes_{\OO_X}\OO_Y,
    \end{align*}
    where $i$ is the closed immersion $Y\xhookrightarrow[\Sch]{i}X$.
\end{remark*}

\begin{proposition}{8.20}
    Let $Y$ be a nonsingular subvariety of codimension $r$ in a nonsingular variety $X$ over $k$. Then, $\omega_Y\simeq\omega_X\otimes_{\OO_X}\bigwedge^r\NN_{Y/X}$ (as $\OO_Y$-modules). In case $r=1$, consider $Y$ as a divisor, and let $\LL\in\MOD_{\OO_X}$ be the associated invertible sheaf $\LL(Y)$. Then, $\omega_Y\simeq\omega_X\otimes\LL\otimes\OO_Y$.
\end{proposition}
\begin{proof}
    (i) We have the exact sequence $0\to\II/\II^2\to\Omega_{X/k}\otimes\OO_Y\to\Omega_{Y/k}\to0$ of $\OO_Y$-modules from \textit{(8.17)}. Then, we have
    \begin{align*}
        \bigwedge\,^r(\II/\II^2)\otimes_{\OO_Y}\omega_Y &=\bigwedge\,^r(\II/\II^2)\otimes_{\OO_Y}\bigwedge\,^{n-r}\Omega_{Y/k}\underset{\text{[5.16d]}}{\simeq}\bigwedge\,^n(\Omega_{X/k}\otimes_{\OO_X}\OO_Y)\\ &=\bigwedge\,^n(i^\ast\Omega_{X/k})\underset{\text{[5.16e]}}{=}i^\ast(\bigwedge\,^n\,\Omega_{X/k})=\omega_X\otimes_{\OO_X}\OO_Y. \tag{$\star$}
    \end{align*}
    Thus, we see that
    \begin{align*}
        \omega_X\otimes_{\OO_X}\bigwedge\,^r\NN_{Y/X} &=\omega_X\otimes_{\OO_X}\bigwedge\,^r(\II/\II^2)^\vee\underset{\text{(Check)}}{\simeq}\omega_X\otimes_{\OO_X}(\bigwedge\,^r\,\II/\II^2)^\vee\\ &=(\omega_X\otimes_{\OO_X}\OO_Y)\otimes_{\OO_Y}(\bigwedge\,^r\,\II/\II^2)^\vee\\ &\underset{(\star)}{\simeq}\underbrace{\bigwedge\,^r\,\II/\II^2}_{\text{invertible}}\,\otimes_{\OO_Y}\,(\bigwedge\,^r\,\II/\II^2)^\vee\otimes_{\OO_Y}\omega_Y\simeq\omega_Y.
    \end{align*}
    (ii) In the special case $r=1$, we have $\II_Y\simeq\LL^\vee$ by \textit{(6.18)}. It follows that $\II/\II^2\simeq\LL^\vee\otimes_{\OO_X}\OO_Y$ and $\NN_{Y/X}\simeq\LL\otimes_{\OO_X}\OO_Y$. Now, applying (i) gives the desired result.
\end{proof}

\begin{example}{8.20.1}
    Let $X=\PP_k^n$ ($n\geq0$). Then, $\omega_X\simeq\OO_X(-n-1)$ and $p_g(\PP_k^n)=0$. Thus, every nonsingular projective \textit{rational} (i.e. birational to some $\PP_k^n$) variety $X$ has genus $p_g(X)=0$. This allows us to demonstrate the existence of nonrational varieties in all dimensions.
\end{example}
\begin{proof}
    We have the exact sequence $0\to\Omega_{X/k}\to\OO_X(-1)^{\oplus(n+1)}\to\OO_X\to0$ of locally free $\OO_X$-modules by \textit{(8.13)}. Note that $\bigwedge^1\OO_X=\OO_X$, so we see that
    \begin{align*}
        \omega_X &\simeq\bigwedge\,^1\,\OO_X\otimes_{\OO_X}\bigwedge\,^n\,\Omega_{X/k}\underset{\text{[5.16d]}}{\simeq}\bigwedge\,^{n+1}\,\OO_X(-1)^{\oplus(n+1)}\\ &\underset{\text{[5.16d]}}{\simeq}\bigwedge\,^n\,\OO_X(-1)^{\oplus n}\otimes_{\OO_X}\bigwedge\,^1\,\OO_X(-1)=\cdots\\ &=\Big(\bigwedge\,^1\,\OO_X(-1)\Big)^{\otimes(n+1)}=\OO_X(-n-1).
    \end{align*}
    Moreover, since $\OO_X(l)$ has no global sections for $l<0$, it follows that $p_g(\PP_k^n)=0$ for every $n\geq0$.
\end{proof}

\chapter{Derived Functors}

\begin{definition}{(Abelian Category)}\newline
    An \textit{abelian category} is a category $\Af$ such that: for any objects $A,B\in\Ob\Af$,
    \begin{itemize}
        \item $\Hom_\Af(A,B)$ has a structure of ablian group.
        \item The composition law is linear.
        \item Finite direct sums exist.
        \item Every morphism has a kernel and a cokernel.
        \item Every monomorphism is the kernel of its cokernel; every epimorphism is the cokernel of its kernel.
        \item Every morphism can be factored into an epimorphism followed by a monomorphism.
    \end{itemize}
\end{definition}

\begin{remark*}{}\mbox{}
    \begin{itemize}
        \item In an abelian category, a morphism is mono [resp. epi, isomorphism] $\Longleftrightarrow$ it is injective [resp. surjective, bijective].
    \end{itemize}
\end{remark*}

\begin{definition}{(Cochain Complex \& Cohomology)}\mbox{}
    \begin{itemize}
        \item A \textit{complex} $A^\bullet$ in an abelian category $\Af$ is a collection of objects $\{A^i\}_{i\in\mathbb{Z}}$ together with morphisms $A^i\xrightarrow[\Af]{d^i}A^{i+1}$ such that $d^{i+1}\circ d^i=0$ for all $i$.
        \item A \textit{morphism of complexes} $(A^\bullet,d_A)$ and $(B^\bullet,d_B)$ in $\Af$, denoted by $f:A^\bullet\to B^\bullet$, is a set of morphisms $f^i:A^i\to B^i$ for each $i$, which commute with the coboundary maps, i.e. $f^{i+1}d_A^i=d_B^if^i$ for all $i$.
        \item The \textit{$i$th cohomology object} $h^i(A^\bullet)$ of the complex $(A^\bullet,d)$ is defined to be $\ker d^i/\Im d^{i-1}$. 
        \item If $f:A^\bullet\to B^\bullet$ is a morphism of complexes, then $f$ indeces a natural map $h^i(f):h^i(A^\bullet)\to h^i(B^\bullet)$ for every $i$.
        \item If $0\to A^\bullet\to B^\bullet\to C^\bullet\to0$ is a short exact sequence of complexes, then there are natural maps $\delta^i:h^i(C^\bullet)\to h^{i+1}(A^\bullet)$ giving rise to a long exact sequence
        \begin{align*}
            \cdots\to h^i(A^\bullet)\to h^i(B^\bullet)\to h^i(C^\bullet)\xrightarrow{\delta^i} h^{i+1}(A^\bullet)\to\cdots
        \end{align*}
    \end{itemize}
\end{definition}

\begin{definition}{(Homotopy)}
    Let $f,g:A^\bullet\to B^\bullet$ be two morphisms of complexes.
    \begin{itemize}
        \item We say $f$ and $g$ are \textit{homotopic}, written $f\sim g$, if there is a collection of morphisms $k^i:A^i\to B^{i-1}$ for each $i$ (which need not commute with the $d^i$) such that $f-g=dk-kd$, (precisely,  $f^i-g^i=d_B^{i-1}k^i+k^{i+1}d_A^i$ for every $i$).
        \item The collection of morphisms $k=\{k^i\}_{i\in \mathbb{Z}}$ is called a \textit{homotopy operator}.
        \item If $f\sim g$, then $f$ and $g$ induce the same morphism $h^i(f)=h^i(g)$ for every $i$.
    \end{itemize}
\end{definition}

\begin{remark}{(Covariant \& Contravariant Functors)}\mbox{}
    \begin{itemize}
        \item \textbf{(Additive Functor).} A covariant [resp. contravariant] functor $F:\Af\to\Bf$ between two abelian categories is \textit{additive} if for any objects $A,A'\in\Ob\Af$, the induced map $\Hom_\Af(A,A')\to\Hom_\Bf(FA,FA')$ [resp. $\Hom_\Af(A,A')\to\Hom_\Bf(FA',FA)$] is a homomorphism of abelian groups.
        \item From now on, we only consider abelian categories and the exact [resp. left exact, right exact] covariant [resp. contravariant] functors between abelian categories are assumed to be additive.
    \end{itemize}
\end{remark}

\begin{definition}{(Resolutions)}\newline
    A $(\ast)$-\textit{resolution} of an object $A\in\Ob\Af$ is a complex $C^\bullet$ together with a morphism $\epsilon:A\to C^0$ such that every $C^i\in\Ob\Af$ has property $(\ast)$, and that the sequence
    \begin{align*}
        0\to A\xrightarrow{\epsilon}C^0\to C^1\to C^2\to\cdots
    \end{align*}
    is exact.
\end{definition}

\begin{definition}{(Injective Objects)}
    \begin{itemize}
        \item An object $I\in\Ob\Af$ is \textit{injective} if the functor $\Hom_\Af(-,I)$ is exact.
        \item If every object of $\Af$ is isomorphic to a subobject of an injective object of $\Af$, we say that $\Af$ has \textit{enough injectives}.
    \end{itemize}
\end{definition}

\begin{lemma}{1.A}\mbox{}
    \begin{itemize}
        \item[(i)] If $\Af$ has enough injectives, then every object has an injective resolution.
        \item[(ii)] If $A\xrightarrow[\Af]{f}B$ is a morphism and both $A$ and $B$ admit injective resolutions $I^\bullet$ and $J^\bullet$, respectively. Then there exists a morphism $f^\bullet:I^\bullet\to J^\bullet$ of complexes such that the diagram 
        \begin{center}
            \begin{tikzcd}
                0\rar[] &A\rar["\epsilon_A"]\dar["f"] &I^0\rar["d_A^0"]\dar["f^0"] &I^1\rar["d_A^1"]\dar["f^1"]
                &I^2\rar["d_A^2"]\dar["f^2"]
                &\cdots\\
                0\rar[] &B\rar["\epsilon_B"] &J^0\rar["d_B^0"] &J^1\rar["d_B^1"]
                &J^2\rar["d_B^2"]
                &\cdots
            \end{tikzcd}
        \end{center}
        commutes.
        \item[(iii)] Any two injective resolutions of the same object are homotopy equivalent.
    \end{itemize}
\end{lemma}
\begin{proof}[Proof of (ii)]
    First, we show the existence of the cochain map $f^\bullet:I^\bullet\to J^\bullet$.
    \begin{itemize}
        \item Note that $J^0\in\Ob\Af$ is injective, the contravariant functor $\Hom_\Af(-,J^0)$ is exact. Apply to the exact sequence $0\to A\rightarrowtail I^0$, we have $\Hom_\Af(I^0,J^0)\overset{-\circ\epsilon_A}{\twoheadrightarrow}\Hom_\Af(A,J^0)\to0$, i.e. $-\circ\epsilon_B$ is epi. Thus, $\exists\ f^0\in\Hom_\Af(I^0,J^0)\mapsto\epsilon_Bf$, i.e. $f^0\epsilon_A=\epsilon_Bf$.
        \item Suppose the $(n-1)$th step holds, i.e. such $f^{n-1}$ exists such that $f^{i+1}d_I^i=d_J^if^i$ for every $0\leq i\leq n-2$. Note that the object $J^n\in\Ob\Af$ is injective, so $\Hom_\Af(-,J^n)$ is an exact functor, which gives the following exact sequence
        \begin{align*}
            \Hom_\Af(I^n,J^n)\xrightarrow{-\circ d_I^{n-1}} \Hom_\Af(I^{n-1},J^n) &\xrightarrow{-\circ d_I^{n-2}} \Hom_\Af(I^{n-2},J^n),\\
            d_J^{n-1}f^{n-1} &\longmapsto d_J^{n-1}f^{n-1}d_I^{n-2}=d_J^{n-1}d_J^{n-2}f^{n-2}=0.
        \end{align*}
        By exactness, $\exists\ f^n\in\Hom_\Af(I^n,J^n)\mapsto d_J^{n-1}f^{n-1}$, i.e. $f^nd_I^{n-1}= d_J^{n-1}f^{n-1}$.
    \end{itemize}
    Now, we show that such cochain map $f^\bullet$ exists up to homotopy. Suppose $g^\bullet:I^\bullet\to J^\bullet$ is another such cochain map. We construct the homotopy operater $h^\bullet:I^\bullet\to J^{\bullet-1}$.
    \begin{itemize}
        \item Define $h^i=0$ for every $i\leq0$.
        \item For $i=1$, we shall find an $h^1:I^1\to J^0$ satisfying $h^1d_I^0=f^0-g^0$. First, since $d_I^0\epsilon_A=0$, we may factor  $d_I^0:I^0\overset{p^0}{\twoheadrightarrow}\Coker\epsilon_A\overset{i^0}{\rightarrowtail}I^1$. Further, note that $(f^0-g^0)\epsilon_A=\epsilon_B(f-f)=0$, factor $f^0-g^0:I^0\overset{p^0}{\twoheadrightarrow}\Coker\epsilon_A\overset{j^0}{\rightarrowtail}J^0$. Apply the exact functor $\Hom_\Af(-,J^0)$ to $0\to\Coker\epsilon_A\overset{i^0}{\rightarrowtail}I^0$, we have $\Hom_\Af(I^1,J^0)\overset{-\circ i^0}{\twoheadrightarrow}\Hom_\Af(\Coker\epsilon_A,J^0)\to0$. So, $\exists\  h^1\in\Hom_\Af(I^1,J^0)\mapsto j^0$. Finally,
        \begin{align*}
            f^0-g^0=j^0p^0=h^1i^0p^0=h^1d_I^0.
        \end{align*}
        \item Suppose the $n$th step holds, i.e. we have such $h^i$ for all $i\leq n$ such that $f^i-g^i=h^{i+1}d_I^i+d_J^{i-1}h^i$ for all $i\leq n-1$. First, factor $d_I^n:I^n\overset{p^n}{\twoheadrightarrow}\Coker d_I^{n-1}\overset{i^n}{\rightarrowtail}I^{n+1}$ as $d_I^nd_I^{n-1}=0$. Note that
        \begin{align*}
            (f^n-g^n-d_J^{n-1}h^n)d_I^{n-1} &=d_J^{n-1}(f^{n-1}-g^{n-1})-d_J^{n-1}h^nd_I^{n-1}\\
            &=d_J^{n-1}(h^nd_I^{n-1}+d_J^{n-2}h^{n-1})-d_J^{n-1}h^nd_I^{n-1}=0.
        \end{align*}
        Factor $f^n-g^n-d_J^{n-1}h^n:I^n\overset{p^n}{\twoheadrightarrow}\Coker d_I^{n-1}\overset{j^n}{\rightarrowtail}J^n$. Apply the exact functor $\Hom_\Af(-,J^n)$ to $0\to\Coker d_I^{n-1}\overset{i^n}{\rightarrowtail}I^{n+1}$, we have $\Hom_\Af(I^{n+1},J^n)\overset{-\circ i^n}{\twoheadrightarrow}\Hom_\Af(\Coker d_I^{n-1},J^n)\to0$. So, $\exists\  h^{n+1}\in\Hom_\Af(I^{n+1},J^n)\mapsto j^n$. Finally,
        \begin{align*}
            f^n-g^n-d_J^{n-1}h^n=j^np^n=h^{n+1}i^np^n=h^{n+1}d_I^n.
        \end{align*}
    \end{itemize}
\end{proof}

\begin{definition}{(Derived Functor)}\newline
    Let $\Af$ be an abelian category with enough injectives, and let $F:\Af\to\Bf$ be a covariant left exact functor. We construct the \textit{right derived functors} $\{R^iF\}_{i\geq0}$ of $F$ as follows. For each $A\in\Ob\Af$, choose once and for all an injective resolution $I^\bullet$ of $A$. Then we define $R^iF(A)\coloneqq h^i(F(I^\bullet))$.
\end{definition}











\end{document}
